% Stéréochimie : isomérie et représentation des molécules — Chimie Tle C — Cours signé OnBuch+
% Programme officiel MINESEC, Terminale C, D, E. Compiler avec : tectonic L05-stereochimie.tex
\newcommand{\DOCMATIERE}{Chimie}
\newcommand{\DOCNIVEAU}{Tle C}
\newcommand{\DOCMODULE}{Module 1 · Chimie organique}
\newcommand{\DOCLECON}{05}
\newcommand{\DOCTITRE}{Stéréochimie : isomérie et représentation des molécules}
% OnBuch+ — préambule « cours signés » ENRICHI (compilé avec tectonic / XeLaTeX)
% Système visuel « Pop » d'OnBuch+ : fond crème (jamais blanc), bordures encre
% épaisses, ombres « dures » décalées, titres Archivo Black en capitales.
% Version enrichie pour les cours de chimie : chimie (mhchem, chemfig),
% graphes (pgfplots), boîtes pédagogiques supplémentaires (définition,
% propriété, méthode, attention, le savais-tu, exercice résolu, exercice,
% TP, bilan), page de garde refaite et signature OnBuch+ ancrée en bas.
%
% Macros attendues dans main.tex AVANT \input{preamble} :
%   \DOCMATIERE  \DOCNIVEAU  \DOCMODULE  \DOCLECON (numéro)  \DOCTITRE
\documentclass[11pt]{article}

\usepackage{fontspec}
\usepackage[french]{babel}
\usepackage[a4paper,margin=2cm,top=3.4cm,bottom=2.8cm,headheight=1.4cm,footskip=1.3cm]{geometry}
\usepackage{amsmath,amssymb,amsfonts}
\usepackage[table]{xcolor} % [table] : fournit aussi \rowcolors (alternance de lignes)
\usepackage{pagecolor}
\usepackage{tikz}
\usetikzlibrary{arrows.meta,positioning,calc,shapes.geometric,shapes.misc,shapes.arrows,decorations.pathmorphing,decorations.markings,patterns,shadows,mindmap,trees,fit,backgrounds,matrix,intersections}
\usepackage{pgfplots}
\pgfplotsset{compat=1.18}
\usepackage[version=4]{mhchem}
\usepackage{chemfig}
\usepackage{enumitem}
\usepackage{fancyhdr}
% [most] charge toutes les bibliothèques tcolorbox (listings, minted, raster,
% documentation...) et gonfle inutilement la mémoire TeX sur un document long
% (~300 boîtes) : on ne charge que ce qui est réellement utilisé.
\usepackage[skins,breakable]{tcolorbox}
\usepackage{graphicx}
\usepackage{colortbl}
\usepackage{booktabs}
\usepackage{tabularx}
\usepackage{array}
\usepackage{multicol}
\usepackage{siunitx}
% parse-numbers=false : accepte aussi \SI{2,5 \times 10^{-3}}{...}
\sisetup{locale=FR,output-decimal-marker={,},group-minimum-digits=4,parse-numbers=false}
\usepackage{qrcode}
% \degree : commande fréquemment utilisée par les modèles pour un angle ou une
% température en dehors de \SI{}{} (non fournie par les paquets chargés).
\providecommand{\degree}{^{\circ}}
\usepackage{lastpage}
\usepackage{hyperref}
\hypersetup{hidelinks}

% ============================================================
% Palette « Pop » OnBuch+ — mêmes hex que la classe P (plus_ui.dart)
% ============================================================
\definecolor{poporange}{HTML}{F59321}
\definecolor{popdark}{HTML}{E07A0C}
\definecolor{popcream}{HTML}{FFF6E8}
\definecolor{popink}{HTML}{1C1714}
\definecolor{poppurple}{HTML}{7A5AE0}
\definecolor{popgreen}{HTML}{2E9E6A}
\definecolor{poppink}{HTML}{E0457B}
\definecolor{popblue}{HTML}{2F7DE1}
\definecolor{popgold}{HTML}{C98A12}
\definecolor{popmuted}{HTML}{8A7F76}
\definecolor{popline}{HTML}{EADFD2}
\definecolor{popgray}{HTML}{8A7F76}
\colorlet{popgrey}{popgray}
% Alias tolérants (le modèle nomme parfois les couleurs différemment)
\colorlet{popwhite}{white}
\colorlet{popblack}{popink}
\colorlet{popred}{poppink}
\colorlet{popyellow}{popgold}
% Teintes claires pour fonds de tableaux / zones
\colorlet{popblueL}{popblue!10!white}
\colorlet{poporangeL}{poporange!14!white}
\colorlet{popgreenL}{popgreen!12!white}
\colorlet{poppurpleL}{poppurple!12!white}
\colorlet{poppinkL}{poppink!12!white}
\colorlet{popgoldL}{popgold!14!white}

\pagecolor{popcream}

% ============================================================
% Typographie
% ============================================================
\defaultfontfeatures{Ligatures=TeX}
\setmainfont{PlusJakartaSans-Medium.ttf}[Path=../../../fonts/,BoldFont=PlusJakartaSans-Bold.ttf]
% Maths en sans-serif (Fira Math) avec chiffres et lettres droites de Plus
% Jakarta, pour que les formules se fondent dans le texte.
\usepackage[math-style=ISO,bold-style=ISO]{unicode-math}
\setmathfont{FiraMath-Regular.otf}[Path=../../../fonts/]
\setmathfont{PlusJakartaSans-Medium.ttf}[Path=../../../fonts/,range={up/num,up/{Latin,latin},\mathup}]
\usepackage{icomma} % virgule décimale sans espace en mode math
% Caractères absents de la police de texte (sinon carré vide) : rendus par
% leur équivalent mathématique, en texte comme en formule.
\usepackage{newunicodechar}
\newunicodechar{α}{\ensuremath{\alpha}}
\newunicodechar{Α}{\ensuremath{\alpha}}
\newunicodechar{β}{\ensuremath{\beta}}
\newunicodechar{δ}{\ensuremath{\delta}}
\newunicodechar{✓}{\popcheck{popgreen}}
\newfontfamily\popdisplay{ArchivoBlack-Regular.ttf}[Path=../../../fonts/]
\newfontfamily\poplogo{SpaceGrotesk-Bold.ttf}[Path=../../../fonts/]
\newfontfamily\popsemi{PlusJakartaSans-SemiBold.ttf}[Path=../../../fonts/]
\newfontfamily\popxbold{PlusJakartaSans-ExtraBold.ttf}[Path=../../../fonts/]
\newfontfamily\popxboldls{PlusJakartaSans-ExtraBold.ttf}[Path=../../../fonts/,LetterSpace=3.0]

\setlist[enumerate]{leftmargin=1.7em,itemsep=5pt,topsep=4pt}
\setlist[itemize]{leftmargin=1.7em,itemsep=4pt,topsep=4pt,label={\color{poporange}\textbullet}}
\setlength{\parindent}{0pt}
\setlength{\parskip}{5pt}
\renewcommand{\arraystretch}{1.25}

% chemfig : liaisons un peu plus épaisses, encre Pop
\setchemfig{atom sep=2em,bond style={line width=0.9pt,popink},cram width=3pt}

% pgfplots : style Pop par défaut
\pgfplotsset{
  every axis/.append style={axis line style={popink,line width=1pt},tick style={popink},
    grid style={popline,line width=0.5pt},label style={font=\small},tick label style={font=\footnotesize}},
  popcurve/.style={poporange,line width=1.8pt,no markers,smooth},
}

% ============================================================
% Pill / squiggle
% ============================================================
\newcommand{\poppill}[3]{%
  \tikz[baseline=-0.55ex]\node[draw=popink,line width=1.2pt,rounded corners=2.6pt,%
    fill=#2,inner xsep=6.5pt,inner ysep=3pt]{%
    \popxboldls\fontsize{7}{7}\selectfont\color{#3}\MakeUppercase{#1}};%
}
\newcommand{\popsquiggle}[1]{%
  \tikz[baseline=-0.4ex]{
    \draw[#1,line width=2.6pt,line cap=round]
      plot [smooth] coordinates {(0,0.09) (0.34,0.01) (0.68,0.21) (1.02,0.01) (1.36,0.10)};
  }%
}
% Mot-clé surligné (marqueur orange sous le texte)
\newcommand{\cle}[1]{{\popxbold\color{popdark}#1}}

% ============================================================
% En-tête / pied
% ============================================================
\pagestyle{fancy}
\fancyhf{}
\renewcommand{\headrulewidth}{0pt}
\renewcommand{\footrulewidth}{0pt}
\lhead{\raisebox{-0.15ex}{\poplogo\fontsize{13}{13}\selectfont\color{popink}OnBuch\color{poporange}+}}
\rhead{\poppill{\DOCMATIERE\ \textbullet\ \DOCNIVEAU\ \textbullet\ Leçon \DOCLECON}{white}{popink}}
\lfoot{\popxbold\fontsize{7.6}{7.6}\selectfont\color{popmuted}\MakeUppercase{Cours signé OnBuch+}\ \textbullet\ {\popsemi\DOCTITRE}}
\rfoot{\tikz[baseline=-0.6ex]\node[rounded corners=8.5pt,fill=popink,minimum height=17pt,minimum width=17pt,inner xsep=5pt,inner sep=0]{\popxbold\fontsize{8}{8}\selectfont\color{popcream}\thepage\,/\,\pageref*{LastPage}};}

% ============================================================
% Titres
% ============================================================
\newcommand{\chaptitre}[2]{%
  \begin{tcolorbox}[enhanced,colback=poporange,colframe=popink,boxrule=2.6pt,arc=11pt,
      drop shadow=popink,left=15pt,right=15pt,top=12pt,bottom=11pt,before skip=3pt,after skip=18pt]
    \poppill{#2}{popink}{popcream}\par\vspace{9pt}
    {\raggedright\hyphenpenalty=10000\exhyphenpenalty=10000\popdisplay\fontsize{22}{24}\selectfont\color{popcream}\MakeUppercase{#1}\par}
  \end{tcolorbox}
}
% Section numérotée : pastille orange avec le numéro + titre Archivo
\newcounter{popsec}
\newcommand{\coursec}[1]{%
  \stepcounter{popsec}\setcounter{popsub}{0}%
  \par\vspace{16pt}\noindent
  \begin{minipage}{\linewidth}
  \tikz[baseline=-0.7ex]\node[draw=popink,line width=1.6pt,fill=poporange,rounded corners=4pt,minimum size=22pt,inner sep=0]{\popdisplay\fontsize{12}{12}\selectfont\color{popcream}\thepopsec};%
  \hspace{8pt}{\popdisplay\fontsize{15}{17}\selectfont\color{popink}\MakeUppercase{#1}}\par
  \vspace{4pt}\noindent{\color{poporange}\rule{36pt}{3.6pt}}
  \end{minipage}\par\nopagebreak\vspace{8pt}
}
\newcounter{popsub}
\newcommand{\courssub}[1]{%
  \stepcounter{popsub}%
  \par\vspace{9pt}\noindent{\popxbold\fontsize{11.5}{13}\selectfont\color{popink}{\color{poporange}\thepopsec.\thepopsub}\hspace{6pt}#1}\par\nopagebreak\vspace{4pt}
}

% ============================================================
% Boîtes pédagogiques — même recette : fond blanc, bordure épaisse colorée,
% ombre dure encre, badge plein. Titre optionnel [#1].
% ============================================================
\tcbset{popbase/.style={breakable,enhanced,colback=white,boxrule=2.3pt,arc=9pt,
  shadow={3pt}{-3pt}{0pt}{popink},left=14pt,right=14pt,top=11pt,bottom=11pt,before skip=11pt,after skip=15pt,
  fontupper=\popsemi\fontsize{10.6}{14.5}\selectfont\color{popink},
  fontlower=\popsemi\fontsize{10.6}{14.5}\selectfont\color{popink}}}
% Coche dessinée (le glyphe ✓ n'existe pas dans les polices du document)
\newcommand{\popcheck}[1]{\tikz[baseline=-0.5ex]\draw[#1,line width=1.6pt,line cap=round,line join=round](-0.13,0.0)--(-0.04,-0.09)--(0.13,0.1);}
\newcommand{\popboxhead}[3]{\poppill{#1}{#2}{white}\ifx\relax#3\relax\else\hspace{6pt}{\popxbold\fontsize{10.6}{12}\selectfont\color{popink}#3}\fi\par\vspace{7pt}}

\newtcolorbox{aretenir}[1][]{popbase,colframe=popgreen,before upper={\popboxhead{À retenir}{popgreen}{#1}}}
\newtcolorbox{exemplebox}[1][]{popbase,colframe=poporange,before upper={\popboxhead{Exemple}{poporange}{#1}}}
\newtcolorbox{definition}[1][]{popbase,colframe=popblue,before upper={\popboxhead{Définition}{popblue}{#1}}}
\newtcolorbox{propriete}[1][]{popbase,colframe=poppurple,before upper={\popboxhead{Propriété}{poppurple}{#1}}}
\newtcolorbox{methode}[1][]{popbase,colframe=popgold,before upper={\popboxhead{Méthode}{popgold}{#1}}}
\newtcolorbox{attention}[1][]{popbase,colframe=poppink,before upper={\popboxhead{Attention}{poppink}{#1}}}
\newtcolorbox{experience}[1][]{popbase,colframe=popblue,colback=popblueL,before upper={\popboxhead{Expérience}{popblue}{#1}}}
% « Le savais-tu ? » : carte violette pleine (contexte, culture, Cameroun)
\newtcolorbox{savaistu}[1][]{popbase,colback=poppurple,colframe=popink,
  fontupper=\popsemi\fontsize{10.4}{14.2}\selectfont\color{white},
  before upper={\poppill{Le savais-tu\,?}{white}{poppurple}\ifx\relax#1\relax\else\hspace{6pt}{\popxbold\color{white}#1}\fi\par\vspace{7pt}}}
% Exercice résolu : énoncé en haut, \tcblower puis la solution sur fond crème
\newcounter{popexo}\newcounter{popexr}
\newtcolorbox{exoresolu}[1][]{popbase,colframe=poporange,colbacklower=popcream,
  segmentation style={popink,line width=1.2pt,dashed},
  before upper={\stepcounter{popexr}\popboxhead{Exercice résolu \thepopexr}{poporange}{#1}},
  before lower={\poppill{Solution}{popink}{popcream}\par\vspace{6pt}}}
% Exercice (non résolu en place) — la correction est dans la partie Corrigés
\newtcolorbox{exercice}[1][]{popbase,colframe=popink,
  before upper={\stepcounter{popexo}\popboxhead{Exercice \thepopexo}{popink}{#1}}}
% Corrigé
\newtcolorbox{corrige}[1][]{popbase,colframe=popgreen,colback=popgreenL,
  before upper={\popboxhead{Corrigé}{popgreen}{#1}}}
% Objectifs / prérequis / bilan
\newtcolorbox{objectifs}{popbase,colframe=popink,colback=poporangeL,
  before upper={\popboxhead{Objectifs de la leçon}{popink}{}}}
\newtcolorbox{prerequis}{popbase,colframe=popink,colback=popblueL,
  before upper={\popboxhead{Prérequis}{popblue}{}}}
\newtcolorbox{bilan}{popbase,colframe=popink,colback=white,boxrule=2.8pt,
  before upper={\popboxhead{Fiche bilan}{poporange}{L'essentiel à connaître pour l'examen}}}
% Figure encadrée avec légende
\newtcolorbox{popfigure}[1][]{enhanced,breakable=false,colback=white,colframe=popline,boxrule=1.4pt,arc=8pt,
  left=10pt,right=10pt,top=10pt,bottom=8pt,before skip=10pt,after skip=12pt,halign=center,
  lower separated=false,halign lower=center,
  fontlower=\popsemi\fontsize{9}{11.5}\selectfont\color{popmuted},
  #1}
\newcommand{\legende}[1]{\par\vspace{6pt}{\popsemi\fontsize{9}{11.5}\selectfont\color{popmuted}#1}}

% ============================================================
% Signature OnBuch+
% ============================================================
\newcommand{\poplogoinline}[1]{{\poplogo\fontsize{#1}{#1}\selectfont\color{popink}OnBuch\color{poporange}+}}
\newcommand{\signatureonbuch}{%
  \begin{tcolorbox}[enhanced,colback=popink,colframe=popink,boxrule=2.2pt,arc=10pt,
      drop shadow=poporange,left=14pt,right=14pt,top=10pt,bottom=10pt,before skip=0pt,after skip=0pt]
    \tikz[baseline=-0.7ex]\node[circle,fill=popgreen,draw=popcream,line width=1.3pt,minimum size=22pt,inner sep=0]{\popcheck{white}};%
    \hspace{9pt}\begin{minipage}[c]{0.62\linewidth}
      {\popxbold\fontsize{12}{13}\selectfont\color{popcream}Signé\ \textbullet\ {\poplogo OnBuch\color{poporange}+}}\par\vspace{2pt}
      {\raggedright\popsemi\fontsize{8}{10.5}\selectfont\color{popline}\MakeUppercase{Équipe pédagogique OnBuch+}\\\MakeUppercase{Conforme au programme officiel MINESEC}\par}
    \end{minipage}\hfill
    {\poplogo\fontsize{22}{22}\selectfont\color{popcream}OnBuch\color{poporange}+}
  \end{tcolorbox}%
}

% ============================================================
% Page de garde
% #1 titre · #2 matière · #3 niveau · #4 module · #5 n° leçon · #6 durée ·
% #7 description / objectifs (texte ou itemize)
% ============================================================
\newcommand{\pagegarde}[7]{%
  \thispagestyle{empty}
  \newgeometry{a4paper,margin=1.7cm,top=1.6cm,bottom=1.4cm}%
  \begin{tikzpicture}[remember picture,overlay]
    \fill[poporange!18] ([xshift=-3.2cm,yshift=-3.6cm]current page.north east) circle (4.6cm);
    \fill[poppurple!14] ([xshift=2.4cm,yshift=7.6cm]current page.south west) circle (3.8cm);
  \end{tikzpicture}%
  {\poplogo\fontsize{30}{30}\selectfont\color{popink}OnBuch\color{poporange}+}\hfill
  \raisebox{0.6ex}{\poppill{Cours signé \textbullet\ Premium}{popink}{popcream}}\par
  \vspace{1pt}\popsquiggle{poppurple}\par
  \vspace{8pt}
  \begin{tcolorbox}[enhanced,colback=poporange,colframe=popink,boxrule=2.8pt,arc=13pt,
      drop shadow=popink,left=18pt,right=18pt,top=12pt,bottom=13pt,after skip=10pt]
    \poppill{#2 \textbullet\ #3}{popink}{popcream}\hspace{5pt}\poppill{#4}{white}{popink}\par\vspace{8pt}
    {\popdisplay\fontsize{14}{14}\selectfont\color{popink}LEÇON #5}\par\vspace{4pt}
    {\raggedright\hyphenpenalty=10000\exhyphenpenalty=10000\popdisplay\fontsize{25}{27}\selectfont\color{popcream}\MakeUppercase{#1}\par}
    \vspace{8pt}
    {\popxbold\fontsize{9.5}{9.5}\selectfont\color{popink}\MakeUppercase{Durée conseillée : #6}}
  \end{tcolorbox}
  \begin{tcolorbox}[enhanced,colback=white,colframe=popink,boxrule=2.2pt,arc=10pt,
      drop shadow=popink,left=16pt,right=16pt,top=10pt,bottom=10pt,after skip=8pt]
    \poppill{Dans cette leçon}{poppurple}{white}\par\vspace{5pt}
    \popsemi\fontsize{9.3}{12.6}\selectfont\color{popink}#7
  \end{tcolorbox}
  \noindent\begin{minipage}[t]{0.487\linewidth}
    \begin{tcolorbox}[enhanced,colback=popgreen,colframe=popink,boxrule=2.2pt,arc=10pt,
        drop shadow=popink,left=11pt,right=11pt,top=10pt,bottom=10pt,height=3.1cm,valign=center]
      \tcbox[colback=white,colframe=popink,boxrule=1.3pt,arc=5pt,boxsep=0pt,nobeforeafter,
        left=3pt,right=3pt,top=3pt,bottom=3pt,valign=center]{\qrcode[height=1.9cm]{https://play.google.com/store/apps/details?id=com.luvvix.onbuch&hl=fr}}%
      \hspace{8pt}\begin{minipage}[c]{0.5\linewidth}
        {\popdisplay\fontsize{11}{12.5}\selectfont\color{white}\MakeUppercase{Télécharge OnBuch}}\par\vspace{4pt}
        {\raggedright\popsemi\fontsize{8.4}{11}\selectfont\color{white}Gratuit \textbullet\ Google Play\\Tuteur IA, annales, concours, résultats\par}
      \end{minipage}
    \end{tcolorbox}
  \end{minipage}\hfill
  \begin{minipage}[t]{0.487\linewidth}
    \begin{tcolorbox}[enhanced,colback=poppurple,colframe=popink,boxrule=2.2pt,arc=10pt,
        drop shadow=popink,left=11pt,right=11pt,top=10pt,bottom=10pt,height=3.1cm,valign=center]
      \tikz[baseline=-0.7ex]\node[circle,fill=white,draw=popink,line width=1.3pt,minimum size=1.6cm,inner sep=0]{\popdisplay\fontsize{24}{24}\selectfont\color{poppurple}+};%
      \hspace{8pt}\begin{minipage}[c]{0.6\linewidth}
        {\popdisplay\fontsize{11}{12.5}\selectfont\color{white}\MakeUppercase{Passe à OnBuch+}}\par\vspace{4pt}
        {\raggedright\popsemi\fontsize{8.4}{11}\selectfont\color{white}Tous les cours signés, Tuteur IA illimité, fiches PDF — onglet OnBuch+ de l'app\par}
      \end{minipage}
    \end{tcolorbox}
  \end{minipage}
  \vspace{8pt}
  \popcontenu
  \vfill
  \signatureonbuch
  \restoregeometry
  \newpage
}

% Rangée « Ce cours contient » (page de garde)
\newcommand{\poptile}[3]{% #1 couleur #2 gros texte #3 légende
  \begin{tcolorbox}[enhanced,colback=#1,colframe=popink,boxrule=2pt,arc=9pt,shadow={2.5pt}{-2.5pt}{0pt}{popink},
      left=6pt,right=6pt,top=9pt,bottom=9pt,halign=center,valign=center,height=1.75cm,width=\linewidth,nobeforeafter]
    {\popdisplay\fontsize{12.5}{13}\selectfont\color{white}\MakeUppercase{#2}}\par\vspace{3pt}
    {\centering\popsemi\fontsize{7.8}{9.5}\selectfont\color{white}#3\par}
  \end{tcolorbox}}
\newcommand{\popcontenu}{%
  \par\noindent{\popxboldls\fontsize{7.5}{7.5}\selectfont\color{popmuted}CE COURS SIGNÉ CONTIENT}\par\vspace{7pt}
  \noindent\begin{minipage}[t]{0.235\linewidth}\poptile{popblue}{Cours}{complet, illustré, pas à pas}\end{minipage}\hfill
  \begin{minipage}[t]{0.235\linewidth}\poptile{poporange}{Méthodes}{et exercices résolus}\end{minipage}\hfill
  \begin{minipage}[t]{0.235\linewidth}\poptile{poppink}{Exercices}{type examen + corrigés}\end{minipage}\hfill
  \begin{minipage}[t]{0.235\linewidth}\poptile{popgreen}{Fiche bilan}{l'essentiel à retenir}\end{minipage}\par}

% Sommaire
\newtcolorbox{sommaire}{enhanced,colback=white,colframe=popink,boxrule=2.2pt,arc=10pt,
  drop shadow=popink,left=18pt,right=18pt,top=13pt,bottom=13pt,before skip=6pt,after skip=20pt,
  before upper={\poppill{Sommaire}{poppurple}{white}\par\vspace{9pt}\popsemi\fontsize{10.6}{14.5}\selectfont\color{popink}}}

% Signature de fin de cours — poussée tout en bas de la dernière page
\newcommand{\findecours}{%
  \par\vfill\nopagebreak
  \begin{center}\popsquiggle{poporange}\end{center}\vspace{4pt}
  \signatureonbuch
}

%% FIGFIX-BEGIN v3 — correctifs globaux des figures (docs/onbuch-generation/tools/figfix)
\usepackage{adjustbox}\usepackage{etoolbox}
% toute figure TikZ/pgfplots plus large que la ligne est réduite à la largeur disponible
\BeforeBeginEnvironment{tikzpicture}{\begin{adjustbox}{max width=\linewidth}}
\AfterEndEnvironment{tikzpicture}{\end{adjustbox}}
% légendes pgfplots lisibles par-dessus les courbes
\pgfplotsset{every axis/.append style={legend style={fill=white,fill opacity=0.92,text opacity=1,draw=black!15,inner sep=3pt}}}
% étiquettes d'arêtes (syntaxe edge["..."]) avec fond blanc
\tikzset{every edge quotes/.append style={fill=white,inner sep=1.2pt}}
% bulles de cartes mentales : texte à la ligne dans la bulle, bulle un peu plus grande
\tikzset{every concept/.append style={text width=2.0cm,align=center,minimum size=2.3cm,inner sep=2pt}}
%% FIGFIX-END
\begin{document}

\pagegarde{Stéréochimie : isomérie et représentation des molécules}{Chimie}{Tle C}{Module 1 · Chimie organique}{05}{4 h}{Cette leçon explore les différentes façons dont des molécules de même formule brute peuvent différer : par leur enchaînement d'atomes (isomérie de constitution) ou par leur arrangement dans l'espace (stéréo-isomérie). Elle fournit les outils de représentation (Cram, Fischer, Newman) indispensables pour décrire la chiralité, les énantiomères, les diastéréoisomères et les conformations, notions au cœur du programme de Terminale et de l'épreuve de chimie du Baccalauréat.
\vspace{4pt}
\begin{multicols}{2}\small
\begin{itemize}
\item À la fin de cette leçon, je sais distinguer les trois types d'isomérie de constitution (chaîne, position, fonction) à partir de formules semi-développées
\item À la fin de cette leçon, je sais représenter une molécule en perspective de Cram en respectant les conventions (traits pleins, hachurés, dans le plan)
\item À la fin de cette leçon, je sais identifier un carbone asymétrique dans une molécule et justifier sa chiralité
\item À la fin de cette leçon, je sais représenter deux énantiomères en perspective de Cram et en projection de Fischer
\item À la fin de cette leçon, je sais distinguer énantiomères et diastéréoisomères, y compris les isomères Z/E
\item À la fin de cette leçon, je sais décrire l'activité optique d'une molécule chirale (dextrogyre, lévogyre, racémique) et interpréter la loi de Biot
\end{itemize}\end{multicols}}

\begin{sommaire}
\begin{multicols}{2}
\begin{enumerate}
\item Isomérie de constitution
\item Représentation spatiale et chiralité
\item Activité optique et pouvoir rotatoire
\item Diastéréoisomères : Z/E et molécules à deux C*
\item Projection de Fischer
\item Conformations : Newman, éthane et cyclohexane
\item Importance de la stéréo-isomérie
\item Travaux pratiques
\item Méthodes et exercices résolus
\item Exercices
\item Corrigés des exercices
\item Fiche bilan
\end{enumerate}
\end{multicols}
\end{sommaire}

\begin{objectifs}
\begin{itemize}
\item À la fin de cette leçon, je sais distinguer les trois types d'isomérie de constitution (chaîne, position, fonction) à partir de formules semi-développées
\item À la fin de cette leçon, je sais représenter une molécule en perspective de Cram en respectant les conventions (traits pleins, hachurés, dans le plan)
\item À la fin de cette leçon, je sais identifier un carbone asymétrique dans une molécule et justifier sa chiralité
\item À la fin de cette leçon, je sais représenter deux énantiomères en perspective de Cram et en projection de Fischer
\item À la fin de cette leçon, je sais distinguer énantiomères et diastéréoisomères, y compris les isomères Z/E
\item À la fin de cette leçon, je sais décrire l'activité optique d'une molécule chirale (dextrogyre, lévogyre, racémique) et interpréter la loi de Biot
\item À la fin de cette leçon, je sais construire les représentations de Newman (éclipsée, décalée) de l'éthane et de ses dérivés
\item À la fin de cette leçon, je sais dessiner les conformations chaise et bateau du cyclohexane et comparer leur stabilité
\item À la fin de cette leçon, je sais citer des exemples concrets (thalidomide, ibuprofène, carvone) montrant l'importance de la stéréo-isomérie
\end{itemize}
\end{objectifs}

\begin{prerequis}
\begin{itemize}
\item Formules brutes, développées et semi-développées des molécules organiques (Première)
\item Géométrie tétraédrique du carbone à 4 liaisons simples et structure plane de la double liaison C=C
\item Nomenclature des alcanes, alcènes, alcools, aldéhydes, cétones et acides carboxyliques
\item Notion de molécule et de liaison covalente ; écriture des formules développées
\item Propagation rectiligne de la lumière et notion d'onde (physique, pour la lumière polarisée)
\end{itemize}
\end{prerequis}

\newpage

\coursec{Isomérie de constitution}

En 1960, un médicament contre les nausées des femmes enceintes, le \cle{thalidomide}, a provoqué la naissance de milliers d'enfants malformés : une seule des deux formes spatiales de la molécule était bénéfique, l'autre était toxique. Pourtant, ces deux molécules ont exactement la même formule brute et les mêmes atomes ! De même, l'orange et le citron vendus au marché de Mokolo doivent leurs odeurs différentes à deux molécules qui sont des images l'une de l'autre dans un miroir. Comment deux molécules « presque identiques » peuvent-elles avoir des effets si différents ? C'est toute la question de la \cle{stéréochimie}, c'est-à-dire de l'arrangement des atomes dans l'espace.

Mais avant d'explorer l'espace, il faut d'abord comprendre une idée plus simple : deux molécules peuvent avoir la même formule brute tout en ayant des \emph{formules développées différentes}. C'est le phénomène d'\cle{isomérie}, et son premier niveau, le plus fondamental, est l'\cle{isomérie de constitution}. C'est elle que nous étudions dans cette section.

\courssub{La notion d'isomérie : même recette, plats différents}

\textbf{Intuition d'abord.} Imagine deux cuisinières de Yaoundé qui reçoivent exactement les mêmes ingrédients : quatre tomates, deux oignons, un poisson. L'une prépare une sauce, l'autre un rôti. Mêmes ingrédients, plats différents, goûts différents. En chimie, les atomes sont les ingrédients : avec les mêmes atomes, on peut construire des molécules différentes, aux propriétés différentes.

\begin{definition}[Isomères]
Deux composés sont \cle{isomères} s'ils ont la \textbf{même formule brute} mais des \textbf{formules développées différentes} (c'est-à-dire un enchaînement différent des atomes). Des isomères ont des \textbf{propriétés physiques et/ou chimiques différentes} : températures de fusion et d'ébullition, densité, odeur, réactivité, activité biologique.
\end{definition}

\begin{attention}[Ce qui n'est PAS un isomère]
Deux molécules de formules brutes différentes ne sont \emph{jamais} isomères, même si elles se ressemblent. Par exemple \ce{CH4} (méthane) et \ce{C2H6} (éthane) ne sont pas isomères : ce sont deux alcanes différents. L'isomérie exige l'\textbf{égalité stricte des formules brutes}.
\end{attention}

On distingue deux grandes familles d'isomérie :
\begin{itemize}
\item l'\cle{isomérie de constitution} : les atomes ne sont pas reliés dans le même ordre (objet de cette section) ;
\item la \cle{stéréo-isomérie} : les atomes sont reliés dans le même ordre, mais arrangés différemment \emph{dans l'espace} (objet des sections suivantes).
\end{itemize}

L'isomérie de constitution se décline elle-même en trois types : l'\cle{isomérie de chaîne}, l'\cle{isomérie de position} et l'\cle{isomérie de fonction}. Étudions-les une par une.

\courssub{Isomérie de chaîne}

\begin{definition}[Isomérie de chaîne]
Deux isomères de \cle{chaîne} ont la même formule brute et la même fonction chimique, mais des \textbf{chaînes carbonées différentes} : linéaire ou ramifiée, ramifications placées différemment.
\end{definition}

\begin{exemplebox}[Le cas de \ce{C4H10}]
La formule brute \ce{C4H10} correspond à \textbf{deux} alcanes :
\begin{itemize}
\item le \textbf{butane} : \chemfig{CH_3-CH_2-CH_2-CH_3} (chaîne linéaire de 4 carbones) ;
\item le \textbf{2-méthylpropane} (isobutane) : \chemfig{CH_3-CH(-[2]CH_3)-CH_3} (chaîne de 3 carbones avec une ramification).
\end{itemize}
Ces deux gaz ont des propriétés différentes : le butane bout à \SI{-0,5}{\celsius}, le 2-méthylpropane à \SI{-11,7}{\celsius}. La ramification « compacte » la molécule, diminue les contacts entre molécules, donc abaisse la température d'ébullition.
\end{exemplebox}

\begin{exemplebox}[Les trois isomères du pentane \ce{C5H12}]
La formule \ce{C5H12} possède \textbf{exactement 3 isomères de chaîne} :
\begin{enumerate}
\item le \textbf{pentane} : \chemfig{CH_3-CH_2-CH_2-CH_2-CH_3} ;
\item le \textbf{2-méthylbutane} : \chemfig{CH_3-CH(-[2]CH_3)-CH_2-CH_3} ;
\item le \textbf{2,2-diméthylpropane} (néopentane) : \chemfig{C(-[2]CH_3)(-[4]CH_3)(-[6]CH_3)-CH_3}.
\end{enumerate}
\end{exemplebox}

\begin{center}
\begin{tabularx}{\linewidth}{|l|X|c|c|}
\hline
\rowcolor{popblueL} \textbf{Isomère} & \textbf{Formule semi-développée} & $\theta_{\text{éb}}$ (\si{\celsius}) & $\theta_{\text{fus}}$ (\si{\celsius}) \\
\hline
Pentane & \ce{CH3-CH2-CH2-CH2-CH3} & $36,1$ & $-129,7$ \\
\hline
2-méthylbutane & \ce{CH3-CH(CH3)-CH2-CH3} & $27,8$ & $-159,9$ \\
\hline
2,2-diméthylpropane & \ce{C(CH3)4} & $9,5$ & $-16,6$ \\
\hline
\end{tabularx}
\end{center}

\begin{propriete}[Effet de la ramification sur l'ébullition]
Pour des isomères de chaîne, \textbf{plus la molécule est ramifiée, plus sa température d'ébullition est basse}. La ramification rend la molécule plus compacte (plus « sphérique »), ce qui diminue la surface de contact entre molécules et donc l'intensité des interactions de Van der Waals.
\end{propriete}

\begin{popfigure}
\begin{tikzpicture}
\begin{axis}[
width=0.85\linewidth, height=6cm,
ybar, bar width=18pt,
symbolic x coords={Pentane, 2-méthylbutane, {2,2-diméthylpropane}},
xtick=data, x tick label style={font=\small, align=center},
ylabel={$\theta_{\text{éb}}$ (\si{\celsius})},
ymin=0, ymax=45,
nodes near coords, nodes near coords style={font=\small},
ymajorgrids, grid style={popline!40},
every axis/.append style={tick style={color=popdark}},
]
\addplot[fill=popblue!70, draw=popblue] coordinates {(Pentane,36.1) (2-méthylbutane,27.8) ({2,2-diméthylpropane},9.5)};
\end{axis}
\end{tikzpicture}
\legende{Températures d'ébullition des trois isomères de \ce{C5H12} : la ramification croissante abaisse régulièrement $\theta_{\text{éb}}$.}
\end{popfigure}

\courssub{Isomérie de position}

\begin{definition}[Isomérie de position]
Deux isomères de \cle{position} ont la même formule brute, la même chaîne carbonée et la même fonction chimique, mais le \textbf{groupe caractéristique} (ou la liaison multiple) n'occupe \textbf{pas la même position} sur la chaîne.
\end{definition}

\begin{exemplebox}[Butan-1-ol et butan-2-ol, \ce{C4H10O}]
Les deux alcools de formule \ce{C4H9OH} à chaîne linéaire :
\begin{itemize}
\item \textbf{butan-1-ol} : \chemfig{CH_3-CH_2-CH_2-CH_2-OH} (alcool \emph{primaire}) ;
\item \textbf{butan-2-ol} : \chemfig{CH_3-CH(-[2]OH)-CH_2-CH_3} (alcool \emph{secondaire}).
\end{itemize}
Même fonction alcool, même chaîne de 4 carbones, mais le groupe \ce{-OH} est en position 1 ou en position 2. Conséquence chimique majeure : l'oxydation ménagée du butan-1-ol conduit à un \textbf{aldéhyde} puis à un acide carboxylique, tandis que celle du butan-2-ol conduit à une \textbf{cétone}.
\end{exemplebox}

\begin{exemplebox}[Pent-1-ène et pent-2-ène, \ce{C5H10}]
\begin{itemize}
\item \textbf{pent-1-ène} : \chemfig{CH_2=CH-CH_2-CH_2-CH_3} ;
\item \textbf{pent-2-ène} : \chemfig{CH_3-CH=CH-CH_2-CH_3}.
\end{itemize}
La double liaison \ce{C=C} n'est pas à la même place : ce sont des isomères de position.
\end{exemplebox}

\begin{attention}[Piège classique de numérotation]
Le pent-1-ène et le « pent-4-ène » sont \textbf{la même molécule} ! On numérote toujours la chaîne de façon à donner à la liaison multiple (ou au groupe caractéristique) l'indice \emph{le plus petit}. Avant de conclure que deux écritures sont des isomères, vérifie la nomenclature : si les deux noms corrects sont identiques, c'est le même composé.
\end{attention}

\courssub{Isomérie de fonction}

\begin{definition}[Isomérie de fonction]
Deux isomères de \cle{fonction} ont la même formule brute mais des \textbf{groupes caractéristiques différents} : ils appartiennent à des \textbf{familles chimiques différentes}. Leurs propriétés chimiques sont alors profondément différentes.
\end{definition}

Voici les trois exemples canoniques du programme.

\begin{exemplebox}[Éthanol et méthoxyméthane, \ce{C2H6O}]
\begin{itemize}
\item \textbf{éthanol} (alcool) : \chemfig{CH_3-CH_2-OH}, liquide bouillant à \SI{78}{\celsius}, soluble dans l'eau, réagit avec le sodium ;
\item \textbf{méthoxyméthane} (éther) : \chemfig{CH_3-O-CH_3}, gaz à température ordinaire ($\theta_{\text{éb}} = \SI{-24}{\celsius}$), très peu réactif.
\end{itemize}
La différence d'ébullition s'explique par les \textbf{liaisons hydrogène} entre molécules d'alcool, impossibles entre molécules d'éther (pas de H porté par O).
\end{exemplebox}

\begin{exemplebox}[Propanal et propanone, \ce{C3H6O}]
\begin{itemize}
\item \textbf{propanal} (aldéhyde) : \chemfig{CH_3-CH_2-C(=[1]O)-[7]H} ;
\item \textbf{propanone} ou acétone (cétone) : \chemfig{CH_3-C(=[1]O)-[7]CH_3}.
\end{itemize}
Le test à la \textbf{liqueur de Fehling} ou au \textbf{réactif de Tollens} les distingue : le propanal, aldéhyde, donne un précipité rouge brique de \ce{Cu2O} avec la liqueur de Fehling chauffée ; la propanone, cétone, ne réagit pas. Tous deux donnent en revanche un précipité jaune avec la DNPH (test commun aux composés carbonylés).
\end{exemplebox}

\begin{exemplebox}[Acide éthanoïque et méthanoate de méthyle, \ce{C2H4O2}]
\begin{itemize}
\item \textbf{acide éthanoïque} : \chemfig{CH_3-C(=[1]O)-[7]OH} (acide carboxylique, le vinaigre en contient) ;
\item \textbf{méthanoate de méthyle} : \chemfig{H-C(=[1]O)-[7]O-CH_3} (ester).
\end{itemize}
L'acide rougit le papier pH et réagit avec les carbonates en dégageant \ce{CO2} ; l'ester, neutre, a une odeur fruitée.
\end{exemplebox}

\begin{popfigure}
\begin{center}
\begin{tikzpicture}[font=\small]
% Chaîne
\node[draw=popblue, fill=popblueL, rounded corners, thick, text width=3.4cm, align=center] (chaine) at (0,0) {\textbf{Chaîne}\\[2pt] \chemfig{CH_3-CH_2-CH_2-CH_3}\\ \chemfig{CH_3-CH(-[2]CH_3)-CH_3}\\ \ce{C4H10}};
% Position
\node[draw=popgreen, fill=popgreenL, rounded corners, thick, text width=3.4cm, align=center] (pos) at (5.6,0) {\textbf{Position}\\[2pt] \chemfig{CH_3-CH_2-CH_2-CH_2-OH}\\ \chemfig{CH_3-CH(-[2]OH)-CH_2-CH_3}\\ \ce{C4H10O}};
% Fonction
\node[draw=poporange, fill=poporangeL, rounded corners, thick, text width=3.4cm, align=center] (fonc) at (11.2,0) {\textbf{Fonction}\\[2pt] \chemfig{CH_3-CH_2-C(=[1]O)-[7]H}\\ \chemfig{CH_3-C(=[1]O)-[7]CH_3}\\ \ce{C3H6O}};
\node[above=2pt of chaine, font=\bfseries\color{popblue}] {la chaîne change};
\node[above=2pt of pos, font=\bfseries\color{popgreen}] {la place du groupe change};
\node[above=2pt of fonc, font=\bfseries\color{poporange}] {la famille change};
\end{tikzpicture}
\end{center}
\legende{Les trois types d'isomérie de constitution : même formule brute dans chaque colonne, mais enchaînement des atomes différent.}
\end{popfigure}

\courssub{Méthode systématique de recherche des isomères}

Face à une formule brute donnée, comment être sûr de n'oublier aucun isomère ? Il faut procéder \emph{par étapes ordonnées}, du changement le plus « doux » au plus « radical ».

\begin{methode}[Trouver tous les isomères d'une formule brute]
\begin{enumerate}
\item \textbf{Chaînes} : écrire d'abord la chaîne carbonée la plus longue possible, puis raccourcir progressivement la chaîne principale en ajoutant des ramifications (méthyle, puis éthyle si possible). Cela donne les isomères de \emph{chaîne}.
\item \textbf{Positions} : pour chaque chaîne obtenue, faire varier la \emph{position} du groupe caractéristique ou de la liaison multiple, en respectant la symétrie de la molécule (ne pas compter deux fois des positions équivalentes).
\item \textbf{Fonctions} : enfin, changer de \emph{fonction} chimique compatible avec la formule brute (alcool $\leftrightarrow$ éther ; aldéhyde $\leftrightarrow$ cétone ; acide $\leftrightarrow$ ester ; alcène $\leftrightarrow$ cycloalcane\ldots).
\item \textbf{Vérifier} : nommer chaque composé obtenu ; deux écritures portant le même nom sont le \emph{même} composé et ne comptent que pour un.
\end{enumerate}
\end{methode}

\begin{exoresolu}[Dénombrer les isomères de \ce{C4H10O}]
Énoncé. Déterminer le nombre et la nature de tous les isomères de constitution de formule brute \ce{C4H10O}.
\tcblower
\textbf{Solution détaillée.} La formule \ce{C4H10O} correspond à la formule générale \ce{C_nH_{2n+2}O} : composés \emph{saturés} à un oxygène, donc des \textbf{alcools} ou des \textbf{éthers}.

\emph{Étape 1 et 2 : les alcools.} Deux squelettes carbonés possibles (butane et 2-méthylpropane), et pour chacun on place le groupe \ce{-OH} :
\begin{enumerate}
\item butan-1-ol : \ce{CH3-CH2-CH2-CH2-OH} (primaire) ;
\item butan-2-ol : \ce{CH3-CH(OH)-CH2-CH3} (secondaire) ;
\item 2-méthylpropan-1-ol : \ce{(CH3)2CH-CH2-OH} (primaire) ;
\item 2-méthylpropan-2-ol : \ce{(CH3)3C-OH} (tertiaire).
\end{enumerate}
Soit \textbf{4 alcools}.

\emph{Étape 3 : les éthers.} On répartit les 4 carbones de part et d'autre de l'oxygène (1+3 ou 2+2) :
\begin{enumerate}
\setcounter{enumi}{4}
\item 1-méthoxypropane : \ce{CH3-O-CH2-CH2-CH3} ;
\item 2-méthoxypropane : \ce{CH3-O-CH(CH3)-CH3} ;
\item éthoxyéthane : \ce{CH3-CH2-O-CH2-CH3}.
\end{enumerate}
Soit \textbf{3 éthers}.

\emph{Conclusion} : \ce{C4H10O} possède \textbf{7 isomères de constitution} : 4 alcools et 3 éthers. Les alcools entre eux sont des isomères de chaîne ou de position ; chaque alcool est isomère de \emph{fonction} de chaque éther.
\end{exoresolu}

\begin{attention}[L'erreur la plus fréquente : le « faux isomère »]
Beaucoup d'élèves croient que deux écritures différentes du même zigzag représentent deux isomères. Par exemple :
\[ \chemfig{CH_3-CH_2-CH_2-CH_3} \qquad \text{et} \qquad \chemfig{CH_3-CH(-[2]CH_3)-[0]CH_3} \]
Non ! La seconde écriture, lue correctement, est aussi le butane : la « ramification » en haut prolonge simplement la chaîne. \textbf{Réflexe de survie} : nomme systématiquement chaque formule. Si les deux noms officiels coïncident, ce n'est pas un isomère, c'est le même composé dessiné autrement. On ne compte un isomère que si le \emph{nom} (ou l'enchaînement des atomes) est réellement différent.
\end{attention}

\courssub{Conséquences de l'isomérie de constitution}

Pourquoi ce découpage est-il si important ? Parce que \textbf{changer la constitution, c'est changer les propriétés}.

\begin{center}
\begin{tabularx}{\linewidth}{|l|X|X|}
\hline
\rowcolor{popblueL} \textbf{Type d'isomérie} & \textbf{Propriétés physiques} & \textbf{Propriétés chimiques} \\
\hline
Chaîne & différentes ($\theta_{\text{éb}}$ baisse avec la ramification) & voisines (même famille) \\
\hline
Position & différentes & proches mais nuances (classe de l'alcool, produits d'oxydation) \\
\hline
Fonction & très différentes & totalement différentes (familles distinctes, tests d'identification distincts) \\
\hline
\end{tabularx}
\end{center}

\begin{exemplebox}[Des tests d'identification pour distinguer les isomères]
Les isomères de fonction \ce{C3H6O} se distinguent expérimentalement :
\begin{itemize}
\item la \textbf{DNPH} donne un précipité jaune-orangé avec les deux (présence du carbonyle \ce{C=O}) ;
\item la \textbf{liqueur de Fehling} à chaud donne un précipité rouge brique de \ce{Cu2O} avec le propanal seulement :
\[ \ce{CH3-CH2-CHO + 2 Cu^2+ + 5 OH- -> CH3-CH2-COO- + Cu2O(v) + 3 H2O} \]
(réaction d'oxydoréduction en milieu basique, à chaud ; l'aldéhyde est oxydé en carboxylate, les ions \ce{Cu^2+} complexes sont réduits en oxyde cuivreux). La propanone ne réagit pas : le test est \textbf{négatif}.
\end{itemize}
\end{exemplebox}

\begin{savaistu}[Isomérie et industrie au Cameroun]
L'isomérie n'est pas qu'un jeu d'écriture : elle a des conséquences industrielles concrètes. Le butane et le 2-méthylpropane, isomères de \ce{C4H10}, n'ont pas les mêmes usages dans le gaz domestique commercialisé au Cameroun : le 2-méthylpropane, plus volatil, est aussi utilisé comme fluide frigorigène dans certains réfrigérateurs modernes. Et lorsque les savonneries artisanales transforment l'huile de palme en savon, la structure exacte des chaînes carbonées des acides gras conditionne la dureté et le pouvoir moussant du savon obtenu. Comprendre l'arrangement des atomes, c'est comprendre la matière qui nous entoure.
\end{savaistu}

\begin{aretenir}[L'essentiel de la section]
\begin{itemize}
\item Deux composés sont \cle{isomères} s'ils ont la même formule brute mais des formules développées différentes ; leurs propriétés physiques et/ou chimiques diffèrent.
\item Trois types d'\cle{isomérie de constitution} : de \textbf{chaîne} (squelette carboné différent), de \textbf{position} (place du groupe caractéristique différente), de \textbf{fonction} (familles chimiques différentes).
\item La ramification \textbf{abaisse} la température d'ébullition des isomères de chaîne.
\item Méthode de dénombrement : chaînes $\rightarrow$ positions $\rightarrow$ fonctions, puis vérification par la nomenclature.
\item \ce{C5H12} : 3 isomères ; \ce{C4H10O} : 7 isomères (4 alcools + 3 éthers).
\item Deux dessins différents peuvent représenter la \emph{même} molécule : toujours nommer avant de conclure.
\end{itemize}
\end{aretenir}

\begin{exercice}[Pour t'entraîner]
\begin{enumerate}
\item Écrire les formules semi-développées et nommer tous les isomères de constitution de formule brute \ce{C5H10} qui sont des alcènes (on ne demande pas de distinguer Z et E).
\item Le 2-méthylbutan-2-ol est-il un isomère de chaîne, de position ou de fonction du pentan-1-ol ? Justifier.
\item On réalise le test à la liqueur de Fehling sur deux isomères de formule \ce{C4H8O}. Le tube A donne un précipité rouge brique, le tube B non. Identifier la famille de chaque composé et proposer une formule semi-développée pour chacun.
\end{enumerate}
\end{exercice}

\begin{corrige}[Exercice « Pour t'entraîner »]
\begin{enumerate}
\item Les alcènes de formule \ce{C5H10} : pent-1-ène \ce{CH2=CH-CH2-CH2-CH3} ; pent-2-ène \ce{CH3-CH=CH-CH2-CH3} ; 2-méthylbut-1-ène \ce{CH2=C(CH3)-CH2-CH3} ; 3-méthylbut-1-ène \ce{CH2=CH-CH(CH3)2} ; 2-méthylbut-2-ène \ce{CH3-C(CH3)=CH-CH3}. Soit \textbf{5 isomères} (le pent-2-ène donnera en outre deux stéréo-isomères Z et E, voir section 4).
\item Même formule brute \ce{C5H12O} et même fonction alcool, mais le squelette est ramifié dans un cas, linéaire dans l'autre : ce sont des isomères de \textbf{chaîne} (la position du \ce{-OH} diffère aussi, mais le critère de chaîne prime car les squelettes diffèrent).
\item Le tube A contient un \textbf{aldéhyde} (réducteur) : par exemple le butanal \ce{CH3-CH2-CH2-CHO}. Le tube B contient une \textbf{cétone} : la butanone \ce{CH3-CO-CH2-CH3}.
\end{enumerate}
\end{corrige}

\coursec{Représentation spatiale et chiralité}

Dans la section précédente, nous avons découvert que des molécules de même formule brute peuvent différer par l'\emph{ordre} de leurs liaisons : c'est l'isomérie de constitution. Mais deux molécules peuvent aussi avoir \textbf{exactement les mêmes liaisons}, dans le même ordre, et pourtant ne pas être identiques ! C'est ce qui se passe quand on extrait le glucose du jus de canne ou quand on compare l'odeur du citron à celle de l'orange : les molécules responsables ont la même connectivité, mais diffèrent par leur \textbf{arrangement dans l'espace}. Bienvenue dans le monde fascinant de la stéréochimie, où la troisième dimension devient indispensable.

\courssub{Le carbone tétraédrique : pourquoi la 3\textsuperscript{e} dimension ?}

\begin{experience}[Le méthane, un carbone dans l'espace]
Considère la molécule de méthane \ce{CH4}. Si le carbone et ses quatre hydrogènes étaient dans un même plan, les quatre liaisons \ce{C-H} formeraient entre elles des angles de \ang{90}. Or les quatre doublets d'électrons de ces liaisons se repoussent mutuellement (théorie VSEPR) : la géométrie la plus stable est celle qui les écarte au maximum les uns des autres. Cette géométrie est le \cle{tétraèdre régulier} : le carbone occupe le centre et les quatre hydrogènes les quatre sommets. L'angle entre deux liaisons vaut alors \ang{109,5}.
\end{experience}

\begin{aretenir}[Le carbone tétraédrique]
Tout atome de carbone formant \textbf{quatre liaisons simples} est \cle{tétraédrique} : ses quatre voisins occupent les sommets d'un tétraèdre régulier dont il est le centre. L'angle entre deux liaisons vaut \ang{109,5}. Une formule développée plane ne peut donc pas rendre compte de la réalité : il faut une \textbf{représentation en trois dimensions}.
\end{aretenir}

\courssub{La représentation de Cram : dessiner la profondeur}

Pour dessiner une molécule tétraédrique sur une feuille de papier (qui est plane !), on utilise la \cle{représentation de Cram}, qui repose sur trois conventions simples.

\begin{definition}[Conventions de la représentation de Cram]
\begin{itemize}
\item un \textbf{trait plein épais} (triangle plein, ou \emph{wedge}) désigne une liaison située \emph{devant} le plan de la feuille, vers l'observateur ;
\item un \textbf{trait hachuré} (pointillés, ou \emph{dash}) désigne une liaison située \emph{derrière} le plan de la feuille ;
\item un \textbf{trait fin} désigne une liaison située \emph{dans} le plan de la feuille.
\end{itemize}
\end{definition}

\begin{attention}[Erreurs fréquentes avec Cram]
Deux pièges classiques : d'abord, le sommet du triangle plein doit partir de l'atome central et s'élargir vers le substituant (c'est le substituant qui « sort » vers toi). Ensuite, n'utilise jamais deux traits pleins ou deux traits hachurés sur le même carbone : la géométrie tétraédrique impose au maximum \textbf{une liaison devant et une liaison derrière}, les deux autres restant dans le plan.
\end{attention}

\begin{exemplebox}[Représentation de Cram du méthane et du chlorométhane]
\begin{itemize}
\item \textbf{Méthane \ce{CH4}} : on place le carbone au centre, deux liaisons \ce{C-H} dans le plan (traits fins), une en wedge, une en dash.
\item \textbf{Chlorométhane \ce{CH3Cl}} : le \ce{Cl} plus gros est souvent mis dans le plan pour la lisibilité, un \ce{H} en wedge, un \ce{H} en dash, le dernier \ce{H} dans le plan.
\end{itemize}
\[
\chemfig{H-[:30](-[2]H)(-[6]H)-[:-30]H} \qquad
\chemfig{Cl-[:30](-[2]H)(-[6]H)-[:-30]H}
\]
\emph{Note : les angles de 30\textdegree\ et -30\textdegree\ dans le code \texttt{chemfig} projettent le tétraèdre sur le plan de la feuille.}
\end{exemplebox}

\courssub{Le carbone asymétrique}

Examinons maintenant le butan-2-ol, \chemfig{CH_3-CH(-[2]OH)-CH_2-CH_3}. Le carbone numéro 2 porte quatre groupes : \ce{-OH}, \ce{-H}, \ce{-CH3} et \ce{-CH2-CH3}. Ces quatre substituants sont \textbf{tous différents}. Ce carbone est dit \emph{asymétrique}.

\begin{definition}[Carbone asymétrique]
Un \cle{carbone asymétrique} est un atome de carbone \textbf{tétraédrique} lié à \textbf{quatre atomes ou groupes d'atomes tous différents}. On le repère par un astérisque : C*.
\end{definition}

\begin{methode}[Repérer un carbone asymétrique]
\begin{enumerate}
\item Cherche les carbones porteurs de \textbf{quatre liaisons simples} (carbones tétraédriques).
\item Pour chacun, liste les quatre atomes ou groupes directement liés, \textbf{y compris les hydrogènes non écrits} dans la formule semi-développée.
\item Si les quatre substituants sont tous différents, le carbone est asymétrique : marque-le C*.
\item Un carbone engagé dans une \textbf{double ou triple liaison} n'est \emph{jamais} asymétrique (il n'est pas tétraédrique).
\end{enumerate}
\end{methode}

\begin{attention}[Piège du comptage des substituants]
L'erreur la plus fréquente consiste à oublier les \textbf{hydrogènes implicites}. Dans \chemfig{CH_3-CH(-[2]OH)-CH_3} (propan-2-ol), le carbone central semble porter trois groupes ; en réalité il porte aussi un hydrogène non écrit. Ses quatre substituants sont \ce{-OH}, \ce{-H}, \ce{-CH3}, \ce{-CH3} : deux sont identiques, donc ce carbone n'est \textbf{pas} asymétrique. Écris toujours les quatre substituants avant de conclure !
\end{attention}

\begin{popfigure}
\centering
\begin{tikzpicture}
\node at (0,0) {\chemfig{C(-[2]H)(<[5]NH_2)(<:[7]CH_3)-COOH}};
\node[popblue, font=\bfseries] at (0,1.8) {Alanine : C* central};
\node[popmuted, font=\small] at (0,-2.2) {4 substituants : \ce{-H}, \ce{-NH2}, \ce{-CH3}, \ce{-COOH}};
\node at (7,0) {\chemfig{CH_3-CH(-[2]OH)-CH_2-CH_3}};
\node[popblue, font=\bfseries] at (7,1.8) {Butan-2-ol : C2 est C*};
\node[popmuted, font=\small] at (7,-2.2) {4 substituants : \ce{-OH}, \ce{-H}, \ce{-CH3}, \ce{-C2H5}};
\end{tikzpicture}
\legende{Repérage du carbone asymétrique : l'alanine (représentée en Cram) et le butan-2-ol. Dans chaque cas, le carbone marqué porte quatre groupes tous différents.}
\end{popfigure}

\courssub{La chiralité : l'analogie des mains}

Regarde tes deux mains paumes vers toi. Elles se ressemblent trait pour trait, et pourtant : aucune rotation ne permet de superposer ta main gauche à ta main droite. Pose-les l'une sur l'autre paumes dans le même sens : les pouces pointent dans des directions opposées. Tes deux mains sont \emph{images l'une de l'autre dans un miroir} et \textbf{non superposables}.

\begin{definition}[Chiralité]
Une molécule est \cle{chirale} si elle n'est \textbf{pas superposable} à son image dans un miroir plan. À l'inverse, une molécule superposable à son image miroir est dite \cle{achirale}.
\end{definition}

\begin{popfigure}
\centering
\begin{tikzpicture}[scale=0.9]
% Main gauche (schématisée)
\draw[thick, popblue, fill=popblueL] (-4.2,0) .. controls (-4.4,1.2) and (-4.0,1.6) .. (-3.8,1.4)
 .. controls (-3.6,2.6) and (-3.2,2.6) .. (-3.2,1.4)
 .. controls (-3.0,2.8) and (-2.6,2.8) .. (-2.6,1.4)
 .. controls (-2.4,2.5) and (-2.0,2.5) .. (-2.0,1.3)
 .. controls (-1.6,1.9) and (-1.3,1.7) .. (-1.6,0.9)
 .. controls (-1.9,0.2) and (-2.2,-0.3) .. (-2.6,-0.3)
 .. controls (-3.4,-0.3) and (-4.0,-0.4) .. (-4.2,0);
\node[popblue, font=\bfseries] at (-3.0,-1.0) {Main gauche};
% Miroir
\draw[very thick, popmuted] (0,-1.6) -- (0,3.2);
\draw[popmuted, pattern=north east lines, pattern color=popmuted] (0,-1.6) -- (0.25,-1.6) -- (0.25,3.2) -- (0,3.2);
\node[popmuted, font=\small] at (0.7,3.5) {miroir};
% Main droite
\draw[thick, poporange, fill=poporangeL] (4.2,0) .. controls (4.4,1.2) and (4.0,1.6) .. (3.8,1.4)
 .. controls (3.6,2.6) and (3.2,2.6) .. (3.2,1.4)
 .. controls (3.0,2.8) and (2.6,2.8) .. (2.6,1.4)
 .. controls (2.4,2.5) and (2.0,2.5) .. (2.0,1.3)
 .. controls (1.6,1.9) and (1.3,1.7) .. (1.6,0.9)
 .. controls (1.9,0.2) and (2.2,-0.3) .. (2.6,-0.3)
 .. controls (3.4,-0.3) and (4.0,-0.4) .. (4.2,0);
\node[poporange, font=\bfseries] at (3.0,-1.0) {Main droite};
% Tentative de superposition
\node[popdark, font=\small, text width=4.5cm, align=center] at (0,-2.6) {Aucune rotation ne superpose les deux mains : elles sont \textbf{chirales}.};
\end{tikzpicture}
\legende{Analogie de la chiralité : les deux mains sont images l'une de l'autre dans un miroir et restent non superposables, quelle que soit la rotation effectuée.}
\end{popfigure}

\begin{savaistu}[D'où vient le mot « chiral » ?]
Le mot \emph{chiral} vient du grec \emph{kheir}, qui signifie « la main ». Il a été introduit par le physicien Lord Kelvin en 1894. Et la vie quotidienne confirme l'analogie : aucun gant droit ne va à la main gauche, aucune chaussure gauche au pied droit ! C'est Louis Pasteur qui, en 1848 à seulement 26 ans, sépara à la pince les cristaux de tartrate de sodium et d'ammonium en deux tas, images l'un de l'autre dans un miroir : la stéréochimie était née.
\end{savaistu}

\begin{propriete}[Lien entre carbone asymétrique et chiralité]
Une molécule possédant \textbf{un seul} carbone asymétrique est toujours chirale. Attention : la réciproque n'est vraie que pour ce cas simple ; nous verrons plus loin (section 4) que certaines molécules à deux C* peuvent être achirales (cas des composés \emph{méso}).
\end{propriete}

\courssub{Les énantiomères : des jumeaux non superposables}

\begin{definition}[Énantiomères]
Deux stéréo-isomères qui sont \textbf{images l'un de l'autre dans un miroir} et \textbf{non superposables} sont appelés \cle{énantiomères}. Une molécule à un seul carbone asymétrique existe donc sous la forme d'\textbf{un couple d'énantiomères} : il y a exactement \textbf{2 stéréo-isomères}.
\end{definition}

Prenons l'acide lactique (acide 2-hydroxypropanoïque), \chemfig{CH_3-CH(-[2]OH)-COOH}, présent dans le lait caillé, le bili-bili et dans nos muscles après un effort intense. Son carbone central porte \ce{-H}, \ce{-OH}, \ce{-CH3} et \ce{-COOH} : il est asymétrique. Il existe donc deux acides lactiques, images dans un miroir.

\begin{popfigure}
\centering
\begin{tikzpicture}
\node at (-3.2,0) {\chemfig{C(-[2]H)(<[5]CH_3)(<:[7]OH)-COOH}};
\node[popblue, font=\bfseries] at (-3.2,-2.2) {Énantiomère A};
% Miroir vertical
\draw[very thick, popmuted] (0,-2.0) -- (0,2.0);
\draw[pattern=north east lines, pattern color=popmuted] (0,-2.0) rectangle (0.22,2.0);
\node[popmuted, font=\small] at (0.75,2.35) {miroir};
\node at (3.2,0) {\chemfig{C(-[2]H)(<:[5]CH_3)(<[7]OH)-COOH}};
\node[poporange, font=\bfseries] at (3.2,-2.2) {Énantiomère B};
\node[popdark, font=\small, text width=9cm, align=center] at (0,-3.2) {Les deux formes de l'acide lactique : images dans un miroir, non superposables. Les liaisons wedge/dash sont inversées.};
\end{tikzpicture}
\legende{Représentation de Cram des deux énantiomères de l'acide lactique \ce{CH3-CHOH-COOH} : les liaisons pleines (wedge) et hachurées (dash) sont inversées par le miroir.}
\end{popfigure}

\begin{aretenir}[Propriétés des énantiomères]
Deux énantiomères ont :
\begin{itemize}
\item les \textbf{mêmes propriétés physiques} (températures de fusion et d'ébullition, densité, solubilité dans un solvant achiral) ;
\item les \textbf{mêmes propriétés chimiques} vis-à-vis de réactifs achiraux ;
\item des \textbf{propriétés différentes} vis-à-vis de la lumière polarisée (section suivante) et vis-à-vis de réactifs ou de récepteurs eux-mêmes chiraux (enzymes, récepteurs biologiques, solvants chiraux) : c'est pourquoi un médicament peut être actif sous une forme et toxique sous l'autre !
\end{itemize}
\end{aretenir}

\begin{definition}[Mélange racémique]
Un \cle{mélange racémique} est un mélange \textbf{équimolaire} (50 \% -- 50 \%) des deux énantiomères d'une même molécule chirale. On le note souvent $(\pm)$ ou \emph{rac}. Ses effets optiques s'annulent : un mélange racémique est optiquement \emph{inactif}.
\end{definition}

\begin{exoresolu}[Compter les stéréo-isomères du pentan-2-ol]
\textbf{Énoncé.} On considère le pentan-2-ol de formule semi-développée \chemfig{CH_3-CH(-[2]OH)-CH_2-CH_2-CH_3}. (1) Cette molécule possède-t-elle un carbone asymétrique ? (2) Combien de stéréo-isomères existe-t-il ? (3) On prépare un mélange contenant \SI{0,30}{mol} de chaque énantiomère : comment appelle-t-on ce mélange et quelle est sa quantité totale de matière ?
\tcblower
\textbf{Solution.} (1) Le carbone n\textsuperscript{o}2 est tétraédrique (quatre liaisons simples). Ses quatre substituants sont : \ce{-OH}, \ce{-H} (hydrogène implicite !), \ce{-CH3} et \ce{-CH2-CH2-CH3}. Ils sont \textbf{tous différents} : le carbone n\textsuperscript{o}2 est asymétrique (C*). (2) Une molécule à un seul C* existe sous la forme d'un couple d'énantiomères : il y a donc \textbf{2 stéréo-isomères}. (3) Le mélange contient des quantités égales des deux énantiomères : c'est un \textbf{mélange racémique}, optiquement inactif. Quantité totale : $n = 0{,}30 + 0{,}30 = \SI{0,60}{mol}$.
\end{exoresolu}

\begin{exercice}[À toi de jouer]
Parmi les molécules suivantes, indique celles qui possèdent un carbone asymétrique et repère-le par un astérisque : (a) le 2-bromobutane \chemfig{CH_3-CH(-[2]Br)-CH_2-CH_3} ; (b) le 2-méthylpropane \chemfig{CH(-[2]CH_3)(-[4]CH_3)(-[6]CH_3)-H} ; (c) l'acide 2-aminopropanoïque (alanine) \chemfig{CH_3-CH(-[2]NH_2)-COOH} ; (d) l'éthanal \chemfig{CH_3-CHO}.
\end{exercice}

\begin{corrige}[Exercice 1]
(a) Le C2 porte \ce{-Br}, \ce{-H}, \ce{-CH3}, \ce{-C2H5} : quatre groupes différents, il est asymétrique (C*). (b) Le carbone central porte \textbf{trois} groupes \ce{-CH3} identiques : pas de C*. (c) Le C2 porte \ce{-NH2}, \ce{-H}, \ce{-CH3}, \ce{-COOH} : asymétrique (C*). (d) Le carbone du groupe \ce{-CHO} porte une \textbf{double liaison} \ce{C=O} : il n'est pas tétraédrique, donc jamais asymétrique.
\end{corrige}

\begin{aretenir}[L'essentiel de la section]
\begin{itemize}
\item Le carbone à quatre liaisons simples est \cle{tétraédrique} (angle \ang{109,5}) : il faut le représenter en 3D par la \cle{représentation de Cram} (trait plein wedge devant, trait hachuré dash derrière, trait fin dans le plan).
\item Un \cle{carbone asymétrique} (C*) porte quatre substituants tous différents ; pense aux hydrogènes implicites !
\item Une molécule \cle{chirale} n'est pas superposable à son image miroir ; une molécule à un C* est toujours chirale.
\item Deux énantiomères sont images dans un miroir, non superposables ; un mélange 50--50 est un \cle{mélange racémique} (optiquement inactif).
\end{itemize}
\end{aretenir}

Maintenant que tu sais reconnaître et représenter ces « jumeaux moléculaires », une question se pose : comment les distinguer expérimentalement, puisqu'ils ont les mêmes propriétés physiques et chimiques usuelles ? La réponse tient en un mot : la \emph{lumière polarisée}. C'est l'objet de la section suivante, consacrée à l'activité optique et au pouvoir rotatoire.

\coursec{Activité optique et pouvoir rotatoire}

Dans la section précédente, nous avons découvert que certaines molécules sont \cle{chirales} : comme nos deux mains, elles ne sont pas superposables à leur image dans un miroir, et elles existent alors sous la forme de deux \cle{énantiomères}. Mais voici une question pratique : comment distinguer expérimentalement deux énantiomères, puisqu'ils ont la même masse molaire, la même température d'ébullition, la même solubilité ? La réponse est aussi élégante qu'étonnante : on utilise la \emph{lumière}. Les molécules chirales possèdent un pouvoir remarquable : elles font \emph{tourner} le plan de vibration de la lumière polarisée. C'est ce phénomène, appelé \cle{activité optique}, que nous allons étudier maintenant, et qui permet de mesurer la concentration du sucre dans les sirops de la SOSUCAM comme dans les laboratoires du monde entier.

\courssub{La lumière polarisée}

\textbf{De la lumière ordinaire à la lumière polarisée.}\par
La lumière est une onde électromagnétique : elle se propage en ligne droite, mais son champ électrique \emph{vibre} perpendiculairement à la direction de propagation. Dans la lumière ordinaire (soleil, lampe), cette vibration s'effectue dans \emph{tous les plans} contenant le rayon lumineux, de façon aléatoire. Imagine une corde que l'on agite dans tous les sens : c'est la lumière naturelle.

Un \cle{polariseur} est un filtre optique (cristal de tourmaline, lame polarisante) qui ne laisse passer que les vibrations situées dans \emph{un seul plan}. Après traversée du polariseur, la lumière est dite \cle{polarisée rectilignement} : son champ électrique vibre dans un plan unique, appelé \emph{plan de polarisation}.

\begin{definition}[Lumière polarisée]
Une lumière est dite \cle{polarisée} (rectilignement) lorsque les vibrations de son champ électrique s'effectuent dans un seul plan, le \emph{plan de polarisation}. On l'obtient en faisant traverser la lumière ordinaire par un \cle{polariseur}.
\end{definition}

\courssub{Le polarimètre : observer la rotation du plan de polarisation}

Pour mettre en évidence l'effet des molécules chirales sur la lumière polarisée, on utilise un appareil imaginé au XIX\textsuperscript{e} siècle : le \cle{polarimètre}.

\begin{experience}[Le polarimètre]
\textbf{Dispositif.} Le polarimètre comprend, alignés sur un axe optique :
\begin{itemize}
\item une \textbf{source lumineuse} (historiquement une lampe à vapeur de \ce{Na}, qui émet une lumière monochromatique jaune, la raie D de longueur d'onde $\lambda = \SI{589}{nm}$) ;
\item un \textbf{polariseur} fixe, qui polarise la lumière ;
\item une \textbf{cuve} (tube) contenant la solution à étudier ;
\item un \textbf{analyseur}, c'est-à-dire un second polariseur \emph{mobile}, que l'on peut faire tourner ;
\item un \textbf{observateur} (ou un capteur) qui mesure l'intensité lumineuse transmise.
\end{itemize}

\textbf{Protocole et observations.}
\begin{enumerate}
\item Cuve vide (ou remplie d'eau pure) : on tourne l'analyseur jusqu'à \emph{l'extinction} (obscurité totale), lorsque son axe est perpendiculaire à celui du polariseur. On note la position de référence.
\item On remplit la cuve d'une solution de glucose (ou d'huile essentielle de citron). \textbf{Observation :} la lumière réapparaît ! L'extinction n'a plus lieu.
\item Pour retrouver l'extinction, il faut \emph{tourner l'analyseur} d'un certain angle $\alpha$.
\end{enumerate}

\textbf{Interprétation.} La solution a fait \emph{pivoter le plan de polarisation} de la lumière d'un angle $\alpha$ pendant la traversée de la cuve. L'angle dont on doit tourner l'analyseur mesure exactement cette rotation.
\end{experience}

\begin{popfigure}
\begin{center}
\begin{tikzpicture}[scale=0.92, >=Stealth]
% Axe optique
\draw[popmuted, dashed, thick] (-0.8,0) -- (13.4,0);
% Source
\draw[fill=popgoldL, draw=popgold] (-0.6,-0.7) rectangle (0.6,0.7);
\node[below, popdark, font=\small] at (0,-0.75) {Source (\ce{Na})};
\draw[popgold, ->] (0.6,0.25) -- (1.5,0.25);
\draw[popgold, ->] (0.6,-0.25) -- (1.5,-0.25);
% Polariseur
\draw[fill=popblueL, draw=popblue, thick] (1.7,-1.0) rectangle (2.1,1.0);
\node[below, popdark, font=\small] at (1.9,-1.05) {Polariseur};
\draw[popblue, ->] (1.9,-0.8) -- (1.9,0.8);
% Lumière polarisée avant cuve : double flèche verticale
\draw[popink, <->, very thick] (3.0,-0.55) -- (3.0,0.55);
% Cuve
\draw[fill=popgreenL, draw=popgreen, thick] (4.0,-0.55) rectangle (8.0,0.55);
\node[popdark, font=\small] at (6.0,0) {solution chirale};
\draw[<->, popdark] (4.0,-0.85) -- (8.0,-0.85) node[midway, below, font=\small] {$\ell$ (dm)};
% Lumière après cuve : double flèche inclinée
\draw[poporange, <->, very thick] (8.6,-0.28) -- ++(50:1.1);
\draw[poporange, <->, very thick] (8.6,-0.28) -- ++(230:1.1);
% Angle alpha
\draw[poporange, ->] (8.6,-0.28) ++(0:1.1) arc (0:50:1.1) node[midway, right, font=\small] {$\alpha$};
% Analyseur
\draw[fill=popblueL, draw=popblue, thick, rotate around={50:(10.3,0)}] (10.1,-1.0) rectangle (10.5,1.0);
\node[below, popdark, font=\small] at (10.6,-1.05) {Analyseur};
\draw[popink, ->] (10.3,1.15) arc (90:50:1.15);
% Observateur
\draw[popdark] (12.3,0.35) circle (0.28);
\draw[popdark, fill=popdark] (12.3,0.35) circle (0.08);
\node[below, popdark, font=\small] at (12.3,-0.4) {Observateur};
\end{tikzpicture}
\end{center}
\legende{Schéma de principe du polarimètre : la lumière polarisée verticalement voit son plan de polarisation tourner d'un angle $\alpha$ en traversant la solution chirale ; on mesure $\alpha$ en tournant l'analyseur.}
\end{popfigure}

\courssub{Substances optiquement actives, dextrogynes et lévogyres}

\begin{definition}[Substance optiquement active]
Une substance est dite \cle{optiquement active} si elle fait tourner le plan de polarisation de la lumière polarisée qui la traverse. L'angle de rotation est noté $\alpha$ et s'exprime en degrés ($\si{\degree}$).
\end{definition}

Quelles substances possèdent cette propriété ? L'expérience le montre sans ambiguïté : ce sont les substances constituées de \emph{molécules chirales}, c'est-à-dire les énantiomères purs (ou les mélanges où l'un des deux énantiomères est en excès). Une molécule achirale (éthanol, acétone, eau) ne dévie jamais le plan de polarisation.

\begin{popfigure}
\begin{center}
\begin{tikzpicture}[scale=0.95, >=Stealth]
% Avant la cuve
\draw[popmuted, dashed] (0,0) -- (2.2,0);
\draw[popink, <->, very thick] (1.1,-0.8) -- (1.1,0.8);
\draw[popink, dashed] (1.1,-0.9) -- (1.1,0.9);
\node[below, popdark, font=\small, align=center] at (1.1,-1.0) {Avant la cuve :\\plan vertical};
% Cuve
\draw[fill=popgreenL, draw=popgreen, thick] (2.6,-0.6) rectangle (5.6,0.6);
\node[popdark, font=\small] at (4.1,0) {glucose (+)};
% Après la cuve
\draw[popmuted, dashed] (5.6,0) -- (7.8,0);
\draw[poporange, <->, very thick, rotate around={40:(6.7,0)}] (6.7,-0.8) -- (6.7,0.8);
\draw[poporange, dashed, rotate around={40:(6.7,0)}] (6.7,-0.9) -- (6.7,0.9);
% Angle alpha centré sur le faisceau (6.7,0)
\draw[poporange, ->] (6.7,0) ++(0:0.9) arc (0:40:0.9) node[midway, right, font=\small] {$\alpha>0$};
\node[below, popdark, font=\small, align=center] at (6.7,-1.0) {Après la cuve :\\plan tourné de $\alpha$};
\end{tikzpicture}
\end{center}
\legende{Traversée d'une solution de glucose (dextrogyre) par la lumière polarisée : le plan de polarisation tourne d'un angle $\alpha$ vers la droite pour l'observateur qui reçoit la lumière.}
\end{popfigure}

\textbf{Dextrogyre ou lévogyre ?}\par
Place-toi à la place de l'observateur qui \emph{reçoit} la lumière (regard dirigé vers la source) :
\begin{itemize}
\item si le plan de polarisation a tourné vers la \emph{droite} (sens des aiguilles d'une montre), la substance est dite \cle{dextrogyre} ; on note sa rotation $\alpha > 0$, signe $(+)$ ;
\item si le plan a tourné vers la \emph{gauche} (sens inverse des aiguilles d'une montre), la substance est dite \cle{lévogyre} ; on note $\alpha < 0$, signe $(-)$.
\end{itemize}

\begin{propriete}[Propriété fondamentale des énantiomères]
Deux \cle{énantiomères} font tourner le plan de polarisation du \emph{même angle} mais en \emph{sens opposés} : si l'un est dextrogyre ($+\alpha$), l'autre est lévogyre ($-\alpha$). C'est la seule propriété physique qui les distingue (avec leur comportement vis-à-vis d'autres objets chiraux).
\end{propriete}

\textbf{Conséquence capitale : le mélange racémique.} Un \cle{mélange racémique} contient les deux énantiomères en quantités égales ($50\,\%-50\,\%$). Chaque molécule dextrogyre fait tourner le plan de $+\alpha$, chaque molécule lévogyre de $-\alpha$ : les deux effets se \emph{compensent exactement}. Le mélange racémique est donc \cle{optiquement inactif} : $\alpha = 0$. Attention : cela ne signifie pas que les molécules sont achirales, mais que leurs effets s'annulent globalement !

\courssub{La loi de Biot : quantifier l'activité optique}

L'angle de rotation dépend logiquement de trois facteurs : la nature de la substance, la \emph{quantité de matière traversée} (concentration et longueur de la cuve), ainsi que la température et la longueur d'onde de la lumière. Le physicien Jean-Baptiste Biot a résumé cela en une loi simple.

\begin{aretenir}[Loi de Biot]
Pour une solution d'une substance optiquement active :
\[ \alpha = [\alpha] \times \ell \times c \]
\begin{itemize}
\item $\alpha$ : angle de rotation du plan de polarisation, en degrés ($\si{\degree}$) ;
\item $[\alpha]$ : \cle{pouvoir rotatoire spécifique} de la substance, constante caractéristique (à température et longueur d'onde données), en $\si{\degree}\cdot\si{\milli\liter}\cdot\si{\gram^{-1}}\cdot\si{\deci\meter^{-1}}$ ;
\item $\ell$ : longueur de la cuve traversée, en \emph{décimètres} (dm) ;
\item $c$ : concentration massique de la solution, en $\si{\gram\per\milli\liter}$ (ou $\si{\gram\per\centi\meter\cubed}$).
\end{itemize}
Le pouvoir rotatoire spécifique se note souvent $[\alpha]_D^{20}$ : mesuré à \SI{20}{\celsius} avec la raie D du sodium.
\end{aretenir}

\begin{exemplebox}[Rotation d'une solution de glucose]
Une solution de glucose de concentration $c = \SI{0,1}{g.mL^{-1}}$ est placée dans une cuve de longueur $\ell = \SI{2}{dm}$. Le pouvoir rotatoire spécifique du glucose est $[\alpha]_D^{20} = \SI{+52,5}{\degree.mL.g^{-1}.dm^{-1}}$.
\[ \alpha = [\alpha] \times \ell \times c = 52,5 \times 2 \times 0,1 = \SI{+10,5}{\degree} \]
Le plan de polarisation tourne de $\SI{10,5}{\degree}$ vers la droite : le glucose est \emph{dextrogyre} (c'est d'ailleurs l'origine de son autre nom : la \emph{dextrose}).
\end{exemplebox}

\begin{methode}[Exploiter la loi de Biot]
\begin{enumerate}
\item Identifier la grandeur inconnue (souvent la concentration $c$).
\item Convertir les unités : $\ell$ \emph{impérativement en dm} ($\SI{10}{cm} = \SI{1}{dm}$), $c$ en $g\cdot mL^{-1}$.
\item Appliquer la loi : $c = \dfrac{\alpha}{[\alpha]\times \ell}$.
\item Conclure en précisant le signe de $\alpha$ (dextrogyre ou lévogyre).
\end{enumerate}
\end{methode}

\begin{exoresolu}[Dosage polarimétrique d'un sirop de sucre]
\textbf{Énoncé.} Une sucrerie contrôle un sirop de saccharose. On dilue le sirop puis on place la solution dans une cuve de $\ell = \SI{20,0}{cm}$. L'analyseur doit tourner de $\alpha = \SI{+13,3}{\degree}$ pour retrouver l'extinction. On donne $[\alpha]_D^{20}(\text{saccharose}) = \SI{+66,5}{\degree.mL.g^{-1}.dm^{-1}}$. Calculer la concentration massique $c$ de la solution dosée.
\tcblower
\textbf{Solution.} Conversion de la longueur : $\ell = \SI{20,0}{cm} = \SI{2,00}{dm}$.\par
D'après la loi de Biot :
\[ c = \frac{\alpha}{[\alpha]\times \ell} = \frac{13,3}{66,5 \times 2,00} = \SI{0,100}{g.mL^{-1}} \]
La solution contient donc $\SI{0,100}{g}$ de saccharose par millilitre, soit $\SI{100}{g.L^{-1}}$. Le signe positif confirme que le saccharose est dextrogyre. Cette méthode, appelée \cle{saccharimétrie}, est rapide, non destructive et ne nécessite aucun réactif : c'est pourquoi elle est utilisée en routine dans l'industrie sucrière.
\end{exoresolu}

\begin{attention}[Erreurs fréquentes]
\begin{itemize}
\item \textbf{Ne confonds jamais le signe optique $(+)$ ou $(-)$ avec la configuration absolue (R ou S) !} Le signe $(+)$ ou $(-)$ est un résultat \emph{expérimental} mesuré au polarimètre : il \emph{ne se devine pas} en regardant la formule de la molécule. Un énantiomère (R) peut être dextrogyre ou lévogyre selon la molécule.
\item \textbf{Unités de la loi de Biot :} la longueur de cuve doit être exprimée en \emph{dm}, pas en cm ! Oublier cette conversion introduit une erreur d'un facteur 10.
\item \textbf{Un angle nul ne signifie pas « molécule achirale » :} cela peut aussi être un mélange racémique (compensation) ou une concentration nulle.
\item Le pouvoir rotatoire dépend de la température et de la longueur d'onde : toujours préciser les conditions ($[\alpha]_D^{20}$).
\end{itemize}
\end{attention}

\begin{savaistu}[La saccharimétrie et la filière sucre au Cameroun]
La polarimétrie appliquée au dosage des sucres porte un nom : la \cle{saccharimétrie}. Elle est utilisée quotidiennement dans les sucreries pour contrôler la concentration des sirops et la pureté du sucre cristallisé. Au Cameroun, la SOSUCAM (Société Sucrière du Cameroun), implantée à Nkoteng et Mbandjock dans le Centre, exploite des milliers d'hectares de canne à sucre : chaque étape de la production, du jus de canne pressé au sucre raffiné, est suivie par des mesures polarimétriques. Historiquement, c'est d'ailleurs en étudiant les sels d'acide tartrique au polarimètre puis au microscope que Louis \textsc{Pasteur}, en 1848, sépara pour la première fois deux énantiomères à la main, fondant ainsi la stéréochimie moléculaire !
\end{savaistu}

\begin{exercice}[Pouvoir rotatoire de l'huile essentielle de citron]
Le limonène extrait des écorces de citron est dextrogyre : $[\alpha]_D^{20} = \SI{+125}{\degree.mL.g^{-1}.dm^{-1}}$. Une solution de concentration $c = \SI{0,04}{g.mL^{-1}}$ est placée dans une cuve de $\SI{1,0}{dm}$.
\begin{enumerate}
\item Calculer l'angle de rotation $\alpha$ observé.
\item On prélève le même volume d'une solution contenant un mélange racémique de limonène à la même concentration totale. Quel angle observe-t-on ? Justifier.
\end{enumerate}
\end{exercice}

\begin{corrige}[Exercice : pouvoir rotatoire du limonène]
\begin{enumerate}
\item Loi de Biot : $\alpha = [\alpha]\times \ell \times c = 125 \times 1,0 \times 0,04 = \SI{+5,0}{\degree}$. Le plan tourne de $\SI{5,0}{\degree}$ vers la droite.
\item Pour le mélange racémique, les deux énantiomères sont présents en quantités égales ; leurs rotations ($+\alpha$ et $-\alpha$) se compensent exactement : on observe $\alpha = \SI{0}{\degree}$. La solution est optiquement inactive bien que ses molécules soient chirales.
\end{enumerate}
\end{corrige}

En résumé, l'activité optique est la \emph{signature expérimentale} de la chiralité : elle permet de distinguer les énantiomères, de vérifier si un échantillon est racémique, et de doser des substances chirales grâce à la loi de Biot. Mais tous les stéréo-isomères ne sont pas des énantiomères : dans la section suivante, nous rencontrerons les \emph{diastéréoisomères}, dont les fameux isomères Z et E des alcènes.

\coursec{Diastéréoisomères : Z/E et molécules à deux C*}

Dans la section précédente, nous avons découvert les \cle{énantiomères} : ces couples de molécules images l'une de l'autre dans un miroir, non superposables, qui font tourner le plan de la lumière polarisée en sens opposés. Mais tous les stéréo-isomères ne sont pas des images miroir ! Imagine deux paires de chaussures : une paire gauche/droite (énantiomères), et deux chaussures de pointures différentes (elles ne sont ni identiques, ni images miroir). Cette seconde situation correspond exactement aux \cle{diastéréoisomères}, que nous allons étudier maintenant à travers deux grandes familles : l'isomérie \cle{Z/E} des alcènes et les molécules possédant \emph{plusieurs} carbones asymétriques.

\courssub{Définition des diastéréoisomères}

\begin{definition}[Diastéréoisomères]
Des \cle{diastéréoisomères} sont des \textbf{stéréo-isomères} (même formule brute, même enchaînement d'atomes, mais arrangement spatial différent) qui \textbf{ne sont pas des images l'un de l'autre dans un miroir}.
\end{definition}

L'idée intuitive est simple : les énantiomères forment des \emph{couples} (main gauche / main droite), tandis que les diastéréoisomères sont des stéréo-isomères « tous les autres ». Contrairement aux énantiomères, qui ont des propriétés physiques identiques (même température d'ébullition, même densité) et ne se distinguent que par leur pouvoir rotatoire, les diastéréoisomères ont des \textbf{propriétés physiques et chimiques différentes} : températures de fusion et d'ébullition, solubilité, réactivité... On peut donc les séparer par des méthodes classiques (distillation, cristallisation, chromatographie), ce qui est impossible pour un mélange racémique.

\begin{aretenir}[Le test en deux questions]
Face à deux stéréo-isomères, pose-toi deux questions :
\begin{enumerate}
\item Sont-ils images l'un de l'autre dans un miroir ? Si \textbf{non} : ce sont des \cle{diastéréoisomères}.
\item Si oui, sont-ils superposables ? Si \textbf{non} : ce sont des \cle{énantiomères}.
\end{enumerate}
\end{aretenir}

\courssub{L'isomérie Z/E des alcènes}

\textbf{Pourquoi la double liaison bloque-t-elle la rotation ?}\par
Autour d'une liaison simple \ce{C-C}, les groupes peuvent tourner librement (nous le verrons avec les conformations). En revanche, une double liaison \ce{C=C} est constituée d'une liaison $\sigma$ \emph{et} d'une liaison $\pi$, formée par le recouvrement latéral de deux orbitales $p$. Tourner autour de la double liaison imposerait de rompre cette liaison $\pi$, ce qui coûte une énergie considérable (environ \SI{270}{kJ.mol^{-1}}). Résultat : à température ordinaire, la molécule est « figée » dans le plan de la double liaison. Les groupes fixés de part et d'autre ne peuvent pas échanger leurs positions : deux arrangements spatiaux distincts peuvent alors exister.

\textbf{La condition d'existence.}\par
Pour que deux arrangements soient \emph{réellement différents}, il faut que \textbf{chaque} carbone de la double liaison porte \textbf{deux groupes différents}. On note la double liaison \ce{Cab=Ccd} : il faut $a \neq b$ \emph{et} $c \neq d$.

\begin{attention}[Piège classique du Bac]
Affirmer que « toute double liaison \ce{C=C} donne lieu à l'isomérie Z/E » est \textbf{faux}. Contre-exemples :
\begin{itemize}
\item Le but-1-ène \ce{CH2=CH-CH2-CH3} : le carbone 1 porte \emph{deux} atomes d'hydrogène identiques $\Rightarrow$ \textbf{pas} d'isomérie Z/E.
\item Le 2-méthylprop-1-ène \ce{CH2=C(CH3)2} : le carbone 1 porte deux H identiques $\Rightarrow$ pas d'isomérie Z/E.
\end{itemize}
Vérifie toujours, sur \emph{chacun} des deux carbones de la double liaison, que les deux substituants sont différents !
\end{attention}

\textbf{La convention Z/E (règles CIP simplifiées).}\par
Sur chaque carbone de la double liaison, on détermine le \cle{groupe prioritaire} en comparant les \textbf{numéros atomiques} des atomes directement liés au carbone de la double liaison : l'atome de numéro atomique le plus élevé donne le groupe prioritaire.
\begin{itemize}
\item Exemples d'ordre de priorité croissant : \ce{H} (Z=1) $<$ \ce{C} (Z=6) $<$ \ce{N} (Z=7) $<$ \ce{O} (Z=8) $<$ \ce{Cl} (Z=17) $<$ \ce{Br} (Z=35) $<$ \ce{I} (Z=53).
\item En cas d'égalité sur le premier atome, on compare les atomes suivants (liste triée par numéro atomique décroissant).
\end{itemize}
On compare alors les positions relatives des deux groupes prioritaires (un sur chaque carbone) :
\begin{itemize}
\item s'ils sont du \textbf{même côté} de la double liaison : stéréo-isomère \cle{Z} (de l'allemand \emph{zusammen}, « ensemble ») ;
\item s'ils sont de \textbf{part et d'autre} : stéréo-isomère \cle{E} (de l'allemand \emph{entgegen}, « opposé »).
\end{itemize}

\begin{methode}[Décider si un alcène est Z ou E]
\begin{enumerate}
\item Repérer la double liaison \ce{C=C} et vérifier que chaque carbone porte deux groupes différents (sinon : pas d'isomérie Z/E, on s'arrête là).
\item Sur chaque carbone, identifier le groupe prioritaire selon la règle du numéro atomique (ex: \ce{CH3} prioritaire sur \ce{H} ; \ce{Cl} prioritaire sur \ce{CH3}).
\item Observer les positions relatives des deux groupes prioritaires : même côté $\Rightarrow$ \textbf{Z} ; côtés opposés $\Rightarrow$ \textbf{E}.
\item Écrire le nom complet : (Z)-but-2-ène ou (E)-but-2-ène.
\end{enumerate}
\end{methode}

\begin{exemplebox}[Les deux but-2-ènes]
Le but-2-ène \ce{CH3-CH=CH-CH3} vérifie la condition : chaque carbone de la double liaison porte un \ce{H} et un \ce{CH3}. Le carbone (Z=6) est prioritaire sur l'hydrogène (Z=1). Il existe donc deux stéréo-isomères :
\begin{itemize}
\item \textbf{(Z)-but-2-ène} : les deux groupes \ce{CH3} du même côté (cis) ;
\item \textbf{(E)-but-2-ène} : les deux groupes \ce{CH3} de part et d'autre (trans).
\end{itemize}
Ces deux molécules ne sont \emph{pas} images miroir l'une de l'autre : ce sont bien des \cle{diastéréoisomères}. Leurs propriétés physiques diffèrent nettement :
\begin{center}
\begin{tabularx}{\linewidth}{|l|X|X|}
\hline
\rowcolor{popblueL} \textbf{Propriété} & \textbf{(Z)-but-2-ène} & \textbf{(E)-but-2-ène} \\
\hline
Température d'ébullition & \SI{3,7}{\celsius} & \SI{0,9}{\celsius} \\
\hline
Température de fusion & \SI{-138,9}{\celsius} & \SI{-105,5}{\celsius} \\
\hline
\end{tabularx}
\end{center}
La molécule (E), plus symétrique, s'empile mieux dans le cristal : sa température de fusion est nettement plus élevée. La molécule (Z), polaire (moment dipolaire non nul), a une ébullition plus haute.
\end{exemplebox}

\begin{popfigure}
\begin{center}
\begin{tikzpicture}
% (Z)-but-2-ene
\node at (-4,3) {\textbf{(Z)-but-2-ène}};
\node at (-4,0) {\chemfig{C(-[2]CH_3)(-[6]H)=C(-[2]CH_3)(-[6]H)}};
\draw[poporange,very thick,->] (-5.5,1.2) -- (-5.5,2.2) node[midway,left,popdark]{\scriptsize \ce{CH3} (prio)};
\draw[poporange,very thick,->] (-2.5,1.2) -- (-2.5,2.2) node[midway,right,popdark]{\scriptsize \ce{CH3} (prio)};
\node[popdark] at (-4,-1.2) {\scriptsize Groupes prioritaires du \textbf{même côté}};

% (E)-but-2-ene
\begin{scope}[xshift=9cm]
\node at (0,3) {\textbf{(E)-but-2-ène}};
\node at (0,0) {\chemfig{C(-[2]CH_3)(-[6]H)=C(-[2]H)(-[6]CH_3)}};
\draw[popgreen,very thick,->] (-1.5,1.2) -- (-1.5,2.2) node[midway,left,popdark]{\scriptsize \ce{CH3} (prio)};
\draw[popgreen,very thick,->] (1.5,-1.2) -- (1.5,-2.2) node[midway,right,popdark]{\scriptsize \ce{CH3} (prio)};
\node[popdark] at (0,-1.2) {\scriptsize Groupes prioritaires de \textbf{part et d'autre}};
\end{scope}
\end{tikzpicture}
\end{center}
\legende{Les deux diastéréoisomères du but-2-ène. Liaison double horizontale. Liaisons pleines : dans le plan. Liaisons \ce{CH3} vers le haut (wedge) ou vers le bas (dash) pour montrer la géométrie. (Z) : les deux \ce{CH3} prioritaires du même côté. (E) : côtés opposés.}
\end{popfigure}

\begin{savaistu}[Voir, c'est isomériser !]
La vision repose sur une isomérisation Z/E ! Dans tes yeux, les cellules de la rétine contiennent un pigment, la rhodopsine, associé à une molécule appelée \emph{rétinal}. Au repos, le rétinal est sous forme \textbf{(Z)} (plus précisément 11-cis). Quand un photon lumineux le frappe, il bascule en quelques centaines de femtosecondes vers la forme \textbf{(E)} (tout-trans). Ce changement de géométrie déforme la protéine, ce qui déclenche l'influx nerveux envoyé au cerveau : tu viens de « voir ». Sans isomérie Z/E, pas de vision !
\end{savaistu}

\courssub{Molécules à plusieurs carbones asymétriques}

\textbf{Combien de stéréo-isomères ?}\par
Chaque carbone asymétrique possède deux configurations spatiales possibles (configuration $R$ ou $S$, que nous verrons en projection de Fischer). Si une molécule contient $n$ carbones asymétriques \emph{différents} (portant des ensembles de substituants distincts), les configurations se combinent indépendamment : le nombre maximal de stéréo-isomères est
\[ N_{\max} = 2^{n} \qquad (n = \text{nombre de carbones asymétriques}). \]

\begin{exemplebox}[Exemple chiffré : deux C* différents]
Une molécule possédant $n = 2$ carbones asymétriques différents (par exemple l'acide 3-bromo-2-chlorobutanoïque \ce{CH3-CHCl-CHBr-COOH}) possède au maximum
\[ N_{\max} = 2^{2} = 4 \ \text{stéréo-isomères}. \]
Ces quatre stéréo-isomères s'organisent en \textbf{deux couples d'énantiomères} : (1 et 2) images miroir, (3 et 4) images miroir. En revanche, 1 et 3 (ou 1 et 4, 2 et 3, 2 et 4) ne sont pas images miroir : ce sont des \cle{diastéréoisomères}. Avec $n = 3$ carbones asymétriques différents, on aurait jusqu'à $2^{3} = 8$ stéréo-isomères, soit 4 couples d'énantiomères.
\end{exemplebox}

\textbf{L'exemple célèbre de l'acide tartrique.}\par
L'acide tartrique, \ce{HOOC-CHOH-CHOH-COOH}, est l'acide du raisin et du vin (il forme les cristaux de « tartre » au fond des cuves et des bouteilles). Il possède deux carbones asymétriques... mais \emph{portant exactement les mêmes groupes} (\ce{H}, \ce{OH}, \ce{COOH}, \ce{CHOHCOOH}). La formule $2^{n}$ donne un maximum de 4, mais on n'observe que \textbf{3 stéréo-isomères} :
\begin{itemize}
\item deux stéréo-isomères (1) et (2), images miroir non superposables : un couple d'\textbf{énantiomères} (l'un dextrogyre, l'autre lévogyre) ;
\item un troisième stéréo-isomère (3), dit \emph{méso}, qui possède un plan de symétrie interne : il est superposable à son image miroir, donc \textbf{achiral} et optiquement inactif.
\end{itemize}
Les relations : (1) et (2) sont énantiomères ; (1) et (3), (2) et (3) sont des \textbf{diastéréoisomères}.

\begin{attention}[La formule $2^{n}$ donne un maximum]
Le nombre $2^{n}$ est un \textbf{maximum}. Lorsque la molécule possède des éléments de symétrie (comme l'acide tartrique, dont les deux C* portent les mêmes groupes), certains stéréo-isomères « attendus » se confondent : le nombre réel est inférieur à $2^{n}$. Le cas de la forme \emph{méso} est un \textbf{complément culturel hors programme}, mais il explique pourquoi l'acide tartrique n'a que 3 stéréo-isomères au lieu de 4 : retiens simplement que $2^{n}$ est la borne supérieure.
\end{attention}

\begin{popfigure}
\begin{center}
\begin{tikzpicture}[scale=0.9]
% Isomere 1 (2R,3R) ou (2S,3S) - Enantiomere 1
\node at (-5,3.5) {\textbf{(1)}};
\node at (-5,1.5) {\chemfig{HOOC-[:-30]C(-[2]H)(<[6]OH)-[:30]C(-[2]OH)(<:[6]H)-[:-30]COOH}};
\draw[popmuted,very thick,dashed] (-3,-0.5) -- (-3,4);
\node[popmuted] at (-3,4.3) {\scriptsize Miroir};

% Isomere 2 (2S,3S) ou (2R,3R) - Enantiomere 2
\node at (0.5,3.5) {\textbf{(2)}};
\node at (0.5,1.5) {\chemfig{HOOC-[:-30]C(-[2]OH)(<[6]H)-[:30]C(-[2]H)(<:[6]OH)-[:-30]COOH}};

% Isomere 3 Meso
\node at (6,3.5) {\textbf{(3) Forme \emph{méso}}};
\node at (6,1.5) {\chemfig{HOOC-[:-30]C(-[2]H)(<[6]OH)-[:30]C(-[2]H)(<:[6]OH)-[:-30]COOH}};
\draw[popgold,thick,dashed] (4.5,1.5) -- (7.5,1.5);
\node[popgold] at (6,0.5) {\scriptsize Plan de symétrie};
\end{tikzpicture}
\end{center}
\legende{Les trois stéréo-isomères de l'acide tartrique \ce{HOOC-CHOH-CHOH-COOH} en représentation de Cram (chaîne carbonée en zigzag). (1) et (2) sont images miroir non superposables : énantiomères. (3) possède un plan de symétrie horizontal : il est achiral (forme méso) et diastéréoisomère de (1) et de (2).}
\end{popfigure}

\begin{exoresolu}[Compter et classer les stéréo-isomères]
\textbf{Énoncé.} On considère la molécule A de formule semi-développée \ce{CH3-CHOH-CHBr-CH3}.
\begin{enumerate}
\item Combien A possède-t-elle de carbones asymétriques ? Quel est le nombre maximal de stéréo-isomères ?
\item Comment sont reliés entre eux ces stéréo-isomères ?
\item Le pent-2-ène \ce{CH3-CH=CH-CH2-CH3} présente-t-il l'isomérie Z/E ? Justifier.
\end{enumerate}
\tcblower
\textbf{Solution détaillée.}
\begin{enumerate}
\item Le carbone 2 porte \ce{H}, \ce{OH}, \ce{CH3} et \ce{CHBr-CH3} : quatre groupes différents $\Rightarrow$ C*. Le carbone 3 porte \ce{H}, \ce{Br}, \ce{CH3} et \ce{CHOH-CH3} : quatre groupes différents $\Rightarrow$ C*. Les carbones 1 et 4 portent trois hydrogènes : non asymétriques. Donc $n = 2$ carbones asymétriques, et comme ils portent des ensembles de groupes différents (pas de symétrie possible) :
\[ N_{\max} = 2^{2} = 4 \ \text{stéréo-isomères}. \]
\item Ces 4 stéréo-isomères forment \textbf{deux couples d'énantiomères} (images miroir non superposables). Deux stéréo-isomères pris dans deux couples différents sont des \textbf{diastéréoisomères} : ils ont des propriétés physiques différentes (points de fusion, solubilités, pouvoirs rotatoires spécifiques non opposés).
\item Oui : sur le carbone 2 de la double liaison, on trouve \ce{H} (Z=1) et \ce{CH3} (atome C, Z=6) $\Rightarrow$ différents ; sur le carbone 3, \ce{H} (Z=1) et \ce{CH2-CH3} (atome C, Z=6) $\Rightarrow$ différents. La condition est remplie : il existe un (Z)-pent-2-ène (groupes carbonés prioritaires du même côté) et un (E)-pent-2-ène (groupes carbonés prioritaires de part et d'autre).
\end{enumerate}
\end{exoresolu}

\courssub{Bilan : l'arbre de classification des isomères}

Récapitulons tout ce que tu as appris depuis le début de la leçon. Deux molécules de même formule brute sont des \cle{isomères}. Deux cas se présentent : soit elles diffèrent par l'\emph{enchaînement des atomes} (isomérie de constitution : chaîne, position, fonction), soit seulement par leur \emph{arrangement dans l'espace} (stéréo-isomérie). Parmi les stéréo-isomères, on distingue les \emph{énantiomères} (images miroir non superposables) et les \emph{diastéréoisomères} (tous les autres : isomérie Z/E, molécules à plusieurs C*, et les conformations que nous étudierons bientôt).

\begin{popfigure}
\begin{center}
\begin{tikzpicture}[
every node/.style={draw,rounded corners=2pt,align=center,font=\small,inner sep=4pt},
niv1/.style={fill=popblueL},
niv2/.style={fill=poporangeL},
niv3/.style={fill=popgreenL},
edge from parent/.style={draw=popmuted,thick},
level 1/.style={sibling distance=7.5cm,level distance=1.6cm},
level 2/.style={sibling distance=4.0cm,level distance=1.6cm},
level 3/.style={sibling distance=2.8cm,level distance=1.5cm}]
\node[niv1,fill=popgoldL] {\textbf{ISOMÈRES}\\ même formule brute}
  child { node[niv1] {Isomérie de\\ \textbf{constitution}}
    child { node[niv3] {de chaîne\\ \scriptsize butane / 2-méthylpropane} }
    child { node[niv3] {de position\\ \scriptsize butan-1-ol / butan-2-ol} }
    child { node[niv3] {de fonction\\ \scriptsize éthanol / méthoxyméthane} } }
  child { node[niv2] {\textbf{Stéréo-isomérie}\\ même enchaînement}
    child { node[niv3] {\textbf{Énantiomères}\\ images miroir\\ non superposables}
      child { node[niv3,fill=poppinkL] {1 C* : acide lactique} }
      child { node[niv3,fill=poppinkL] {n C* différents} } }
    child { node[niv3] {\textbf{Diastéréoisomères}\\ non images miroir}
      child { node[niv3,fill=poppinkL] {Z / E\\ (alcènes rigides)} }
      child { node[niv3,fill=poppinkL] {n C* (dont méso)} }
      child { node[niv3,fill=poppinkL] {Conformations\\ (Newman, cyclohexane)} } } };
\end{tikzpicture}
\end{center}
\legende{Arbre de classification des différents types d'isomérie. Les énantiomères et les diastéréoisomères sont les deux grandes familles de stéréo-isomères. Les conformations (Section 6) sont des diastéréoisomères interconvertibles rapidement.}
\end{popfigure}

\begin{aretenir}[L'essentiel de la section]
\begin{itemize}
\item Des \cle{diastéréoisomères} sont des stéréo-isomères qui ne sont \textbf{pas} images miroir l'un de l'autre ; leurs propriétés physiques et chimiques \textbf{diffèrent} (séparables par méthodes physiques classiques).
\item L'isomérie \cle{Z/E} existe pour un alcène \ce{Cab=Ccd} si $a \neq b$ \textbf{et} $c \neq d$. Priorité = numéro atomique de l'atome voisin. \textbf{Z} : priorités même côté ; \textbf{E} : priorités côtés opposés.
\item Une molécule à $n$ carbones asymétriques possède \textbf{au maximum} $2^{n}$ stéréo-isomères ; la symétrie interne (cas méso) peut réduire ce nombre.
\item Deux stéréo-isomères d'une molécule à plusieurs C* sont soit énantiomères (images miroir), soit diastéréoisomères.
\end{itemize}
\end{aretenir}

\begin{exercice}[À toi de jouer]
\begin{enumerate}
\item Parmi les alcènes suivants, lesquels présentent l'isomérie Z/E ? Justifier par la condition sur les substituants : (a) but-1-ène ; (b) hex-3-ène ; (c) 2-méthylbut-2-ène ; (d) 1,2-dichloroéthène \ce{ClCH=CHCl}.
\item La molécule de cholestérol possède 8 carbones asymétriques. Calculer le nombre maximal de stéréo-isomères théoriquement possibles. Commenter l'enjeu pour la synthèse en pharmacie.
\item Représenter en perspective (Cram) le (Z)-but-2-ène et le (E)-but-2-ène (liaison double horizontale, substituants en wedge/dash), puis préciser leur relation d'isomérie.
\end{enumerate}
\end{exercice}

\begin{corrige}[Exercice : À toi de jouer]
\begin{enumerate}
\item (a) But-1-ène \ce{CH2=CH-CH2-CH3} : le carbone 1 porte deux H identiques $\Rightarrow$ \textbf{non}. (b) Hex-3-ène \ce{CH3CH2-CH=CH-CH2CH3} : chaque carbone de la double liaison porte H et \ce{C2H5} (différents) $\Rightarrow$ \textbf{oui} (formes Z et E). (c) 2-Méthylbut-2-ène \ce{CH3-C(CH3)=CH-CH3} : le carbone 2 porte deux \ce{CH3} identiques $\Rightarrow$ \textbf{non}. (d) \ce{ClCH=CHCl} : chaque carbone porte H (Z=1) et Cl (Z=17) $\Rightarrow$ \textbf{oui} : (Z) avec les deux Cl du même côté, (E) côtés opposés.
\item $N_{\max} = 2^{8} = 256$ stéréo-isomères possibles, dont un seul est le cholestérol naturel biologiquement actif. Cela montre pourquoi la synthèse stéréosélective (contrôle de chaque C*) est un défi majeur de l'industrie pharmaceutique : synthétiser le bon isomère parmi 256 !
\item Représentations Cram : liaison \ce{C=C} horizontale. (Z) : deux \ce{CH3} en wedge (ou deux en dash). (E) : un \ce{CH3} en wedge, l'autre en dash. Ce sont des \textbf{diastéréoisomères} : stéréo-isomères non images miroir, propriétés physiques différentes.
\end{enumerate}
\end{corrige}

\coursec{Projection de Fischer}

Dans la section précédente, nous avons appris à distinguer les \cle{diastéréoisomères} des \cle{énantiomères}, en manipulant des représentations de Cram où les liaisons sont dessinées en perspective (coins pleins vers l'avant, hachures vers l'arrière). Cette représentation est fidèle, mais elle devient vite fastidieuse lorsque la molécule possède plusieurs carbones asymétriques : imaginez dessiner en perspective le glucose, qui en compte quatre ! Les chimistes avaient besoin d'une écriture \emph{plane}, rapide et sans ambiguïté. C'est exactement ce qu'a inventé le chimiste allemand Emil Fischer à la fin du XIX\textsuperscript{e} siècle : la \cle{projection de Fischer}, que tu retrouveras partout en biochimie, du manuel de Terminale jusqu'aux laboratoires de recherche.

\courssub{Principe et conventions de la projection de Fischer}

\textbf{Une idée simple : aplatir la molécule sans trahir sa géométrie.}

Prenons une molécule chirale comme l'alanine, \ce{CH3-CH(NH2)-COOH}, l'un des acides aminés qui constituent les protéines de notre alimentation (viande, poisson, niébé). Son carbone asymétrique porte quatre groupes différents : \ce{-COOH}, \ce{-NH2}, \ce{-CH3} et \ce{-H}. En représentation de Cram, deux de ces groupes sont dans le plan de la feuille, un sort vers nous, un rentre derrière. Fischer a eu l'idée suivante : orientons la molécule de sorte que la chaîne carbonée soit \emph{verticale} et \emph{derrière} le plan, les deux autres substituants étant \emph{horizontaux} et \emph{devant} le plan, puis écrasons le tout sur la feuille. On obtient une croix : le carbone asymétrique est à l'intersection des deux traits, et on ne l'écrit même plus.

\begin{definition}[Projection de Fischer]
La \cle{projection de Fischer} est une représentation \emph{plane} d'une molécule chirale dans laquelle :
\begin{itemize}
\item le \cle{carbone asymétrique} se trouve à l'\textbf{intersection} de deux traits perpendiculaires (il n'est pas écrit) ;
\item les liaisons \textbf{horizontales} sortent du plan de la feuille, \textbf{vers l'observateur} ;
\item les liaisons \textbf{verticales} rentrent \textbf{derrière} le plan de la feuille.
\end{itemize}
Par convention, la chaîne carbonée principale est placée \textbf{verticalement}, avec le groupe le plus \textbf{oxydé} (celui dont le carbone porte le plus d'atomes d'oxygène ou le degré d'oxydation le plus élevé) \textbf{en haut}.
\end{definition}

\begin{attention}[La convention est la clé de tout]
La croix de Fischer n'est \emph{pas} un simple dessin plat : chaque trait cache une information spatiale. Un élève qui oublie que « horizontal = vers soi, vertical = vers l'arrière » lira la projection à l'envers et confondra deux énantiomères. Retiens l'image suivante : la molécule est comme une personne qui \emph{te fait face}, les bras écartés vers toi (liaisons horizontales), la tête et les pieds rejetés en arrière (liaisons verticales).
\end{attention}

\textbf{Exemples fondateurs.} Trois molécules classiques à mémoriser :

\begin{itemize}
\item le \textbf{glycéraldéhyde} \ce{HOCH2-CHOH-CHO} : le groupe le plus oxydé est l'aldéhyde \ce{-CHO}, placé en haut ; le \ce{-CH2OH} est en bas ; \ce{-OH} et \ce{-H} occupent les positions horizontales ;
\item l'\textbf{alanine} \ce{CH3-CH(NH2)-COOH} : l'acide carboxylique \ce{-COOH} en haut, le méthyle \ce{-CH3} en bas, \ce{-NH2} et \ce{-H} à l'horizontale ;
\item l'\textbf{acide lactique} \ce{CH3-CHOH-COOH} (présent dans le lait caillé et dans nos muscles après l'effort) : \ce{-COOH} en haut, \ce{-CH3} en bas, \ce{-OH} et \ce{-H} à l'horizontale.
\end{itemize}

\courssub{Passage de la représentation de Cram à la projection de Fischer}

\begin{methode}[Convertir une représentation de Cram en projection de Fischer]
\begin{enumerate}
\item \textbf{Identifier} le carbone asymétrique et ses quatre substituants.
\item \textbf{Orienter} mentalement la molécule : faire pivoter la molécule (dans l'espace, c'est autorisé !) de façon que la chaîne carbonée principale soit verticale, groupe le plus oxydé en haut, et que les deux liaisons de cette chaîne pointent \emph{vers l'arrière} (liaisons hachurées en Cram).
\item \textbf{Vérifier} que les deux substituants restants sont alors horizontaux et pointent \emph{vers l'avant} (liaisons en coin plein en Cram).
\item \textbf{Écraser} la molécule sur le plan : dessiner la croix, placer les quatre groupes, sans écrire le carbone central.
\item \textbf{Contrôler} le résultat : horizontal = vers soi, vertical = vers l'arrière.
\end{enumerate}
\end{methode}

\begin{popfigure}
\begin{center}
\begin{tikzpicture}[scale=1.1, every node/.style={font=\small}]
% --- Cram de l'alanine (Conformation éclipsée propice à Fischer) ---
% Carbone central
\coordinate (C) at (0,0);
% Liaisons verticales : VERS L'ARRIERE (hachurées)
% COOH en haut
\draw[popdark, thick, decorate, decoration={zigzag, segment length=3, amplitude=0.8, pre length=2pt, post length=2pt}] (C) -- ++(90:1.6) node[above, popdark] {\ce{COOH}};
% CH3 en bas
\draw[popdark, thick, decorate, decoration={zigzag, segment length=3, amplitude=0.8, pre length=2pt, post length=2pt}] (C) -- ++(-90:1.6) node[below, popdark] {\ce{CH3}};
% Liaisons horizontales : VERS L'AVANT (coins pleins)
% NH2 à gauche
\fill[poporange!80!black] (C) -- ++(150:0.3) -- ++(150:1.0) -- ++(-30:0.3) -- cycle;
\node[poporange!80!black, left] at (-1.4,0.3) {\ce{H2N}};
% H à droite
\fill[popblue!80!black] (C) -- ++(30:0.3) -- ++(30:1.0) -- ++(150:0.3) -- cycle;
\node[popblue!80!black, right] at (1.4,0.3) {\ce{H}};
% Légendes Cram
\node[font=\bfseries\small, popblue] at (0,2.4) {Représentation de Cram (conformation éclipsée)};
\node[font=\footnotesize\itshape, poporange!80!black] at (-2.0,-0.8) {Liaisons horizontales :\newline vers l'avant (coin plein)};
\node[font=\footnotesize\itshape, popdark] at (1.8,-1.5) {Liaisons verticales :\newline vers l'arrière (hachurées)};

% --- Flèche de conversion ---
\draw[-{Stealth[length=5mm]}, very thick, popgreen] (3.2,0) -- (4.8,0) node[midway, above, font=\footnotesize\itshape, align=center]{On écrase\\sur le plan};

% --- Fischer de l'alanine ---
\begin{scope}[shift={(8,0)}]
\coordinate (Cf) at (0,0);
% Croix
\draw[thick, popdark] (0,1.2) -- (0,-1.2); % Vertical
\draw[thick, popdark] (-1.2,0) -- (1.2,0); % Horizontal
% Groupes
\node[above] at (0,1.2) {\ce{COOH}};
\node[below] at (0,-1.2) {\ce{CH3}};
\node[left] at (-1.2,0) {\ce{H2N}};
\node[right] at (1.2,0) {\ce{H}};
% Indication centre
\fill[popgreen!30] (0,0) circle (0.25);
\node[font=\footnotesize, popgreen!60!black] at (0,0) {\ce{C*}};
% Titre
\node[font=\bfseries\small, popblue] at (0,2.4) {Projection de Fischer};
\node[font=\footnotesize\itshape, align=center, popmuted] at (0,-2.2) {Horizontal : vers soi\\Vertical : vers l'arrière};
\end{scope}
\end{tikzpicture}
\end{center}
\legende{De la représentation de Cram de l'alanine (conformation éclipsée : chaîne verticale derrière, substituants horizontaux devant) à sa projection de Fischer.}
\end{popfigure}

\begin{exemplebox}[L'acide lactique en projection de Fischer]
L'acide lactique \ce{CH3-CHOH-COOH} possède un carbone asymétrique porteur de \ce{-COOH}, \ce{-OH}, \ce{-CH3} et \ce{-H}. Le groupe le plus oxydé est \ce{-COOH} : il va en haut. Le \ce{-CH3} va en bas. Restent \ce{-OH} et \ce{-H} à l'horizontale. Il existe donc \textbf{deux} projections possibles : \ce{-OH} à droite et \ce{-H} à gauche, ou l'inverse. Ces deux projections sont images l'une de l'autre dans un miroir et représentent les \textbf{deux énantiomères} de l'acide lactique : l'acide lactique produit par nos muscles (lévogyre) et celui des ferments lactiques du lait caillé ou du bili-bili (dextrogyre) ne font pas tourner la lumière polarisée dans le même sens !
\end{exemplebox}

\courssub{Projection de Fischer et molécules à plusieurs carbones asymétriques}

La puissance de la projection de Fischer apparaît pleinement pour les molécules à plusieurs carbones asymétriques (glucides, acides aminés à chaîne longue). On aligne alors \textbf{tous} les carbones de la chaîne principale sur l'axe vertical, le plus oxydé en haut. Chaque carbone asymétrique devient une croix. Les liaisons horizontales de \emph{chaque} croix sortent vers l'observateur ; les liaisons verticales (chaîne) restent derrière.

\begin{exemplebox}[Érythrose et Thréose : diastéréoisomères en Fischer]
Les aldoses à 4 carbones (tétroses) possèdent \textbf{deux} carbones asymétriques (C2 et C3). On place le \ce{-CHO} en haut.
\begin{itemize}
\item \textbf{D-Érythrose} : \ce{-OH} sur C2 à \textbf{droite}, \ce{-OH} sur C3 à \textbf{droite}.
\item \textbf{D-Thréose} : \ce{-OH} sur C2 à \textbf{droite}, \ce{-OH} sur C3 à \textbf{gauche}.
\end{itemize}
Ces deux molécules ne sont pas images miroir : ce sont des \textbf{diastéréoisomères}. La projection de Fischer permet de voir immédiatement la différence de configuration sur l'un des centres chiraux.
\end{exemplebox}

\courssub{Lecture d'une projection de Fischer : reconnaître les énantiomères}

Une propriété remarquable rend la projection de Fischer très pratique : \textbf{deux projections de Fischer images l'une de l'autre dans un miroir vertical représentent deux énantiomères}. En effet, le miroir échange la droite et la gauche, c'est-à-dire les deux groupes horizontaux de \emph{chaque} carbone asymétrique, ce qui inverse la configuration de tous les centres chiraux.

\begin{popfigure}
\begin{center}
\begin{tikzpicture}[scale=1.0, every node/.style={font=\small}]
% --- Fischer gauche : D-glycéraldéhyde ---
\node[font=\bfseries\small, popblue] at (-3,2.4) {(+)-D-glycéraldéhyde};
\node at (-3,1.3) {\ce{CHO}};
\node at (-3,-1.3) {\ce{CH2OH}};
\node at (-4.2,0) {\ce{H}};
\node at (-1.8,0) {\ce{OH}};
\draw[thick, popdark] (-3,0.95) -- (-3,-0.95);
\draw[thick, popdark] (-3.85,0) -- (-2.15,0);
% --- Miroir ---
\draw[thick, poppurple, decorate, decoration={zigzag, segment length=6, amplitude=1.2}] (0,-1.9) -- (0,1.9);
\node[font=\footnotesize\itshape, poppurple] at (0,2.3) {Miroir vertical};
% --- Fischer droite : L-glycéraldéhyde ---
\node[font=\bfseries\small, poporange] at (3,2.4) {(--)-L-glycéraldéhyde};
\node at (3,1.3) {\ce{CHO}};
\node at (3,-1.3) {\ce{CH2OH}};
\node at (4.2,0) {\ce{H}};
\node at (1.8,0) {\ce{OH}};
\draw[thick, popdark] (3,0.95) -- (3,-0.95);
\draw[thick, popdark] (2.15,0) -- (3.85,0);
% Flèches
\draw[{Stealth[length=3mm]}-{Stealth[length=3mm]}, popmuted] (-1.6,0.6) -- (1.6,0.6) node[midway, above, font=\footnotesize\itshape]{Images miroir};
\node[font=\footnotesize\itshape, align=center, popmuted] at (0,-2.5) {Non superposables : deux énantiomères};
\end{tikzpicture}
\end{center}
\legende{Les deux énantiomères du glycéraldéhyde \ce{HOCH2-CHOH-CHO} en projection de Fischer : le passage de l'un à l'autre se fait par échange des groupes horizontaux \ce{-H} et \ce{-OH} sur le carbone asymétrique.}
\end{popfigure}

\begin{aretenir}[Règle d'or de lecture (un carbone asymétrique)]
Pour une molécule à \textbf{un} carbone asymétrique, échanger les deux groupes \textbf{horizontaux} (ou les deux groupes verticaux) dans une projection de Fischer revient à passer d'un énantiomère à l'autre. Échanger un groupe horizontal avec un groupe vertical change aussi la configuration : la nouvelle projection représente l'énantiomère.
\end{aretenir}

\courssub{Règles de manipulation d'une projection de Fischer}

Voici le point le plus délicat — et le plus souvent piégé au Baccalauréat. Une projection de Fischer n'est pas un dessin ordinaire : on ne peut pas la manipuler n'importe comment, car toute manipulation doit \emph{conserver la convention} « horizontal vers soi, vertical vers l'arrière ».

\textbf{Rotation de 180° dans le plan : autorisée.} Tourne ta feuille d'un demi-tour. Les groupes qui étaient horizontaux restent horizontaux (donc toujours vers toi), les verticaux restent verticaux (toujours derrière). La convention est respectée : c'est \textbf{la même molécule}, avec la même configuration absolue.

\textbf{Rotation de 90° dans le plan : interdite.} Après un quart de tour, les groupes initialement verticaux (qui devaient être derrière) deviennent horizontaux (donc dessinés comme étant devant), et inversement. La convention est violée : le dessin obtenu représente \textbf{l'énantiomère}, pas la même molécule.

\textbf{« Retourner » la projection comme une crêpe (retournement hors plan) : interdit.} Si tu retournes la feuille pour regarder la projection par derrière, les liaisons qui sortaient vers toi rentrent désormais : tu obtiens l'énantiomère.

\textbf{Échanges de substituants :} Un échange unique (horizontal/horizontal, vertical/vertical ou horizontal/vertical) inverse la configuration (énantiomère). \textbf{Deux échanges successifs} (nombre pair) ramènent à la configuration initiale (même molécule).

\begin{methode}[Vérifier que deux projections de Fischer représentent la même molécule]
\begin{enumerate}
\item Ne jamais sortir la projection du plan de la feuille (pas de rotation 3D, pas de retournement).
\item N'appliquer \textbf{que des rotations de 180°} dans le plan (autant de fois que nécessaire).
\item Éventuellement, effectuer un \textbf{nombre pair} d'échanges de substituants (deux échanges successifs ramènent à la même configuration).
\item Si l'on peut superposer les deux projections par ces seules opérations : c'est la \textbf{même molécule}. Sinon, ce sont deux stéréo-isomères (énantiomères s'il n'y a qu'un \ce{C*} ; diastéréoisomères s'il y en a plusieurs).
\end{enumerate}
\end{methode}

\begin{attention}[Erreurs fréquentes à bannir]
\begin{itemize}
\item Tourner la projection de \textbf{90°} et conclure « c'est la même molécule » : \textbf{faux}, on a obtenu l'énantiomère.
\item « Retourner » la projection comme une crêpe (retournement hors plan) : même piège, on inverse la configuration.
\item Écrire le carbone asymétrique (\ce{C} ou \ce{C*}) à l'intersection des traits : par convention, il est \emph{sous-entendu}, on ne l'écrit pas.
\item Placer n'importe quel groupe en haut : le groupe le plus \textbf{oxydé} doit être en haut de la chaîne verticale (exemple : \ce{-CHO} avant \ce{-CH2OH}, \ce{-COOH} avant \ce{-CH3}, \ce{-COOH} avant \ce{-CHO}).
\item Confondre projection de Fischer et représentation de Cram : en Fischer, \textbf{tous} les horizontaux sortent, \textbf{tous} les verticaux rentrent. En Cram, on mélange pleins, hachurés et traits pleins dans le plan.
\end{itemize}
\end{attention}

\begin{exoresolu}[Deux projections de l'alanine : même molécule ou énantiomères ?]
\textbf{Énoncé.} On considère trois projections de Fischer de l'alanine.\\
Projection A : \ce{COOH} en haut, \ce{CH3} en bas, \ce{NH2} à gauche, \ce{H} à droite.\\
Projection B : \ce{CH3} en haut, \ce{COOH} en bas, \ce{NH2} à droite, \ce{H} à gauche.\\
Projection C : \ce{COOH} en haut, \ce{CH3} en bas, \ce{NH2} à droite, \ce{H} à gauche.\\
Comparer A et B, puis A et C.
\tcblower
\textbf{Solution détaillée.}
\emph{Comparaison de A et B.} Faisons tourner la projection A de 180° dans le plan de la feuille : le \ce{COOH} passe en bas, le \ce{CH3} en haut, le \ce{NH2} (qui était à gauche) passe à droite, et le \ce{H} passe à gauche. On obtient exactement la projection B. Or une rotation de 180° dans le plan est \textbf{autorisée} : A et B représentent donc \textbf{la même molécule}, le même énantiomère de l'alanine (configuration L usuelle).

\emph{Comparaison de A et C.} A et C ont les mêmes groupes verticaux (\ce{COOH} en haut, \ce{CH3} en bas) mais les groupes horizontaux \ce{NH2} et \ce{H} sont échangés. Un seul échange inverse la configuration du carbone asymétrique : A et C sont \textbf{images miroir non superposables}, ce sont deux \textbf{énantiomères}. Aucune rotation de 180° ne permet de passer de A à C (la rotation de 180° de A donne B, qui a \ce{CH3} en haut, donc différent de C).

\emph{Conclusion.} A et B : même molécule. A et C : énantiomères.
\end{exoresolu}

\begin{exercice}[À toi de jouer]
On donne la projection de Fischer de l'acide lactique : \ce{COOH} en haut, \ce{CH3} en bas, \ce{OH} à droite, \ce{H} à gauche (projection 1). On considère la projection 2 : \ce{CH3} en haut, \ce{COOH} en bas, \ce{OH} à gauche, \ce{H} à droite ; et la projection 3 : \ce{COOH} en haut, \ce{CH3} en bas, \ce{OH} à gauche, \ce{H} à droite. La projection 2 représente-t-elle la même molécule que la 1 ? Et la projection 3 ? Justifier rigoureusement.
\end{exercice}

\begin{corrige}[Exercice : acide lactique]
Rotation de 180° de la projection 1 dans le plan : \ce{CH3} passe en haut, \ce{COOH} en bas, \ce{OH} (à droite) passe à gauche, \ce{H} (à gauche) passe à droite. On obtient exactement la projection 2. La rotation de 180° étant autorisée, les projections 1 et 2 représentent \textbf{la même molécule} (même configuration absolue). La projection 3 diffère de la 1 par un seul échange des groupes horizontaux (\ce{OH} et \ce{H}) : c'est \textbf{l'énantiomère} de la projection 1 (configuration opposée).
\end{corrige}

\begin{savaistu}[Emil Fischer, le père de la stéréochimie des sucres]
Emil Fischer (1852--1919), prix Nobel de chimie en 1902, a introduit cette projection pour classer les \textbf{sucres} (glucose, fructose, galactose\ldots) dont les molécules possèdent plusieurs carbones asymétriques : dessiner toutes leurs perspectives aurait été un cauchemar ! Grâce à sa convention, il a pu établir la famille des sucres dits « D » (le \ce{-OH} du \textbf{dernier} carbone asymétrique — le plus éloigné du groupe carbonyl — à \textbf{droite}) et « L » (à \textbf{gauche}), classification encore utilisée aujourd'hui. La projection de Fischer reste l'outil standard en biochimie pour représenter les \textbf{glucides} et les \textbf{acides aminés} : tous les acides aminés $\alpha$-aminés naturels des protéines — comme l'alanine du niébé ou du poisson braisé — appartiennent à la série \textbf{L} (groupe \ce{-NH2} à gauche en projection de Fischer avec \ce{-COOH} en haut), un fait remarquable que les scientifiques expliquent encore mal !
\end{savaistu}

\begin{aretenir}[L'essentiel de la section]
\begin{itemize}
\item La projection de Fischer est une représentation \textbf{plane} d'une molécule chirale : le \ce{C*} est à l'intersection des traits, les liaisons \textbf{horizontales viennent vers l'observateur}, les \textbf{verticales partent derrière}.
\item La chaîne principale est verticale, le groupe le plus \textbf{oxydé en haut}. Pour plusieurs \ce{C*}, la chaîne s'aligne verticalement.
\item Deux projections images miroir (échange des horizontaux sur \emph{tous} les \ce{C*}) = deux \textbf{énantiomères}.
\item Manipulation autorisée : \textbf{rotation de 180°} dans le plan (ou un nombre pair d'échanges). Interdites : rotation de 90°, retournement de la feuille — on obtiendrait l'énantiomère (ou un diastéréoisomère si plusieurs \ce{C*}).
\end{itemize}
\end{aretenir}

\coursec{Conformations : Newman, éthane et cyclohexane}

Dans la section précédente, nous avons appris à représenter des molécules chirales en projection de Fischer : il s'agit de \emph{configurations} fixes, qui ne peuvent changer qu'en rompant des liaisons covalentes. Mais une molécule organique n'est pas un édifice rigide : autour de chaque liaison simple \ce{C-C}, les groupes peuvent \emph{tourner} librement, un peu comme les pales d'un ventilateur autour de leur axe. La molécule adopte alors, instantanément, des formes spatiales différentes sans jamais changer de structure. Ce sont les \cle{conformations}, dernier niveau de la stéréochimie au programme de Terminale. Tu les rencontreras partout : dans le glucose de ton sang, dans les huiles de palme, dans les carburants de la SONARA.

\courssub{Stéréo-isomères de conformation : une idée simple}

\begin{definition}[Conformation]
Des \cle{conformations} (ou stéréo-isomères de conformation) sont des arrangements spatiaux différents d'une \emph{même} molécule, qui s'interconvertissent par \textbf{simple rotation autour d'une liaison simple} \ce{C-C}, \emph{sans rupture d'aucune liaison}.
\end{definition}

L'intuition est la suivante : prends deux cure-dents plantés dans une boule de pâte à modeler (la liaison \ce{C-C}). Tu peux faire tourner un cure-dent par rapport à l'autre autour de l'axe commun : la molécule « change de forme », mais rien n'a été cassé ni reconstruit. Contrairement aux énantiomères et aux diastéréoisomères (isomères de \emph{configuration}), les conformations ne sont \textbf{pas isolables} : à température ordinaire, elles se transforment les unes dans les autres des milliards de fois par seconde.

Pour décrire une conformation, on utilise l'\cle{angle dièdre} $\theta$ : c'est l'angle entre deux liaisons situées sur les deux carbones voisins, mesuré en regardant le long de l'axe \ce{C-C}.

\courssub{La représentation de Newman}

Comment dessiner ce que l'on voit quand on regarde une molécule « dans l'axe » d'une liaison \ce{C-C} ? Le chimiste américain Melvin Newman a proposé en 1952 une convention simple et universelle.

\begin{methode}[Construire une représentation de Newman]
\begin{enumerate}
\item Placer l'\oe il \textbf{dans l'axe de la liaison} \ce{C-C} étudiée, comme si tu regardais dans un tuyau.
\item Le carbone \textbf{avant} (le plus proche de l'\oe il) est représenté par un \textbf{point} ; ses trois substituants partent du point, à \SI{120}{\degree} les uns des autres.
\item Le carbone \textbf{arrière} (masqué par le premier) est représenté par un \textbf{cercle} ; ses trois substituants partent du \emph{bord} du cercle, également à \SI{120}{\degree}.
\item Orienter les substituants arrière selon l'angle dièdre choisi : face à ceux de l'avant (éclipsée, $\theta = \SI{0}{\degree}$) ou dans les intervalles (décalée, $\theta = \SI{60}{\degree}$).
\end{enumerate}
\end{methode}

\begin{popfigure}
\begin{center}
\begin{tikzpicture}[scale=1.1, line cap=round, line join=round]
% ===== Newman éclipsée =====
\begin{scope}
% Cercle arrière
\draw[popline, thick] (0,0) circle (1.0);
% Liaisons avant (point central)
\fill[popink] (0,0) circle (0.07);
\draw[popblue, very thick] (0,0) -- (90:1.0) node[above, popdark, font=\small] {\ce{H}};
\draw[popblue, very thick] (0,0) -- (210:1.0) node[below left, popdark, font=\small] {\ce{H}};
\draw[popblue, very thick] (0,0) -- (330:1.0) node[below right, popdark, font=\small] {\ce{H}};
% Liaisons arrière (éclipsées : mêmes angles, dessinées depuis le bord du cercle vers l'extérieur)
\draw[poporange, very thick] (90:1.0) -- (90:1.55) node[above, popdark, font=\small] {\ce{H}};
\draw[poporange, very thick] (210:1.0) -- (210:1.55) node[below left, popdark, font=\small] {\ce{H}};
\draw[poporange, very thick] (330:1.0) -- (330:1.55) node[below right, popdark, font=\small] {\ce{H}};
\node[popdark, font=\small] at (0,-1.8) {\textbf{Éclipsée} : $\theta = \SI{0}{\degree}$};
\node[popmuted, font=\small] at (0,-2.2) {Énergie maximale (gêne maximale)};
\end{scope}
% ===== Newman décalée =====
\begin{scope}[xshift=7cm]
\draw[popline, thick] (0,0) circle (1.0);
\fill[popink] (0,0) circle (0.07);
\draw[popblue, very thick] (0,0) -- (90:1.0) node[above, popdark, font=\small] {\ce{H}};
\draw[popblue, very thick] (0,0) -- (210:1.0) node[below left, popdark, font=\small] {\ce{H}};
\draw[popblue, very thick] (0,0) -- (330:1.0) node[below right, popdark, font=\small] {\ce{H}};
% Liaisons arrière décalées de 60° (angles : 30, 150, 270)
\draw[poporange, very thick] (30:1.0) -- (30:1.55) node[above right, popdark, font=\small] {\ce{H}};
\draw[poporange, very thick] (150:1.0) -- (150:1.55) node[above left, popdark, font=\small] {\ce{H}};
\draw[poporange, very thick] (270:1.0) -- (270:1.55) node[below, popdark, font=\small] {\ce{H}};
\node[popdark, font=\small] at (0,-1.8) {\textbf{Décalée} : $\theta = \SI{60}{\degree}$};
\node[popmuted, font=\small] at (0,-2.2) {Énergie minimale (gêne minimale)};
\end{scope}
\end{tikzpicture}
\end{center}
\legende{Représentations de Newman de l'éthane \ce{CH3-CH3} vue selon l'axe \ce{C-C}. En bleu : liaisons du carbone avant (point) ; en orange : liaisons du carbone arrière (bord du cercle).}
\end{popfigure}

\courssub{Les conformations de l'éthane : éclipsée et décalée}

L'éthane \ce{CH3-CH3} est la molécule la plus simple possédant une liaison \ce{C-C}. En faisant tourner un groupe \ce{CH3} par rapport à l'autre, on distingue deux conformations extrêmes.

\textbf{La conformation éclipsée} ($\theta = \SI{0}{\degree}$) : chaque hydrogène du carbone avant se trouve exactement \emph{devant} un hydrogène du carbone arrière. Les doublets d'électrons des liaisons \ce{C-H} se font face et se repoussent : c'est la \cle{gêne maximale}, donc l'\textbf{énergie maximale}. C'est la conformation la plus instable.

\textbf{La conformation décalée} ($\theta = \SI{60}{\degree}$) : chaque hydrogène de l'arrière se loge \emph{entre} deux hydrogènes de l'avant. Les répulsions sont minimales : c'est la conformation la plus \textbf{stable}, d'énergie minimale.

Entre les deux, l'énergie varie continuellement. La différence d'énergie entre la conformation éclipsée et la conformation décalée s'appelle la \cle{barrière de rotation} : elle vaut environ \SI{12}{kJ.mol^{-1}} pour l'éthane. C'est très faible : l'agitation thermique à température ordinaire suffit largement à la franchir. La rotation autour de la liaison \ce{C-C} est donc \textbf{quasi libre} à \SI{25}{\celsius}, et l'éthane passe sans cesse d'une conformation à l'autre.

\begin{popfigure}
\begin{center}
\begin{tikzpicture}
\begin{axis}[
  width=0.9\linewidth, height=5.5cm,
  xlabel={Angle dièdre $\theta$ (degrés)},
  ylabel={Énergie potentielle (unités arbitraires)},
  xmin=0, xmax=360,
  xtick={0,60,120,180,240,300,360},
  ytick=\empty,
  axis lines=left,
  every axis x label/.style={at={(ticklabel* cs:1.02)}, anchor=west},
  every axis y label/.style={at={(ticklabel* cs:1.02)}, anchor=south},
  clip=false
]
\addplot[popcurve, domain=0:360, samples=300] {1-cos(3*x)};
\node[popmuted, font=\small] at (axis cs:30,1.8) {Éclipsée};
\node[popmuted, font=\small] at (axis cs:90,0.2) {Décalée};
\node[popmuted, font=\small] at (axis cs:150,1.8) {Éclipsée};
\node[popmuted, font=\small] at (axis cs:210,0.2) {Décalée};
\node[popmuted, font=\small] at (axis cs:270,1.8) {Éclipsée};
\node[popmuted, font=\small] at (axis cs:330,0.2) {Décalée};
\draw[popgreen, <->, thick] (axis cs:355,0.05) -- (axis cs:355,1.95) node[midway, right, font=\small, popdark] {$\approx \SI{12}{kJ.mol^{-1}}$};
\end{axis}
\end{tikzpicture}
\end{center}
\legende{Énergie potentielle de l'éthane en fonction de l'angle de rotation autour de la liaison \ce{C-C} : trois maxima (conformations éclipsées) et trois minima (conformations décalées) par tour complet.}
\end{popfigure}

\begin{attention}[Conformation n'est pas configuration]
Erreur classique au Bac : affirmer que la forme éclipsée et la forme décalée de l'éthane sont « deux isomères différents que l'on peut séparer ». C'est \textbf{faux}. Les conformations s'interconvertissent librement par simple rotation, sans rupture de liaison : on ne peut \emph{pas} les isoler dans deux flacons distincts. Seuls les isomères de \emph{configuration} (énantiomères, diastéréoisomères, isomères $Z$/$E$) sont séparables, car leur interconversion exigerait de casser des liaisons.
\end{attention}

\courssub{Complément utile : le butane vu selon la liaison C2–C3}

Le raisonnement s'étend au butane \ce{CH3-CH2-CH2-CH3}, observé dans l'axe de la liaison centrale \ce{C2-C3}. Chaque carbone porte alors un groupe méthyle \ce{CH3} et deux hydrogènes. La présence de groupes volumineux (\ce{CH3}) complexifie le profil énergétique : trois conformations décalées (minima) et trois éclipsées (maxima) se succèdent.

\begin{popfigure}
\begin{center}
\begin{tikzpicture}[scale=1.0, line cap=round, line join=round]
% --- Anti ---
\begin{scope}
\draw[popline, thick] (0,0) circle (1.0);
\fill[popink] (0,0) circle (0.07);
% Avant : CH3 en haut (90), H à 210, H à 330
\draw[popblue, very thick] (0,0) -- (90:1.0) node[above, popdark, font=\small] {\ce{CH3}};
\draw[popblue, very thick] (0,0) -- (210:1.0) node[below left, popdark, font=\small] {\ce{H}};
\draw[popblue, very thick] (0,0) -- (330:1.0) node[below right, popdark, font=\small] {\ce{H}};
% Arrière : CH3 en bas (270 = 90+180), H à 30, H à 150
\draw[poporange, very thick] (270:1.0) -- (270:1.55) node[below, popdark, font=\small] {\ce{CH3}};
\draw[poporange, very thick] (30:1.0) -- (30:1.55) node[above right, popdark, font=\small] {\ce{H}};
\draw[poporange, very thick] (150:1.0) -- (150:1.55) node[above left, popdark, font=\small] {\ce{H}};
\node[popdark, font=\small] at (0,-1.8) {\textbf{Anti} ($\theta=\SI{180}{\degree}$)};
\node[popgreen, font=\small] at (0,-2.2) {La plus stable};
\end{scope}
% --- Gauche ---
\begin{scope}[xshift=7cm]
\draw[popline, thick] (0,0) circle (1.0);
\fill[popink] (0,0) circle (0.07);
% Avant : CH3 en haut (90), H à 210, H à 330
\draw[popblue, very thick] (0,0) -- (90:1.0) node[above, popdark, font=\small] {\ce{CH3}};
\draw[popblue, very thick] (0,0) -- (210:1.0) node[below left, popdark, font=\small] {\ce{H}};
\draw[popblue, very thick] (0,0) -- (330:1.0) node[below right, popdark, font=\small] {\ce{H}};
% Arrière : CH3 à 30 (gauche), H à 150, H à 270
\draw[poporange, very thick] (30:1.0) -- (30:1.55) node[above right, popdark, font=\small] {\ce{CH3}};
\draw[poporange, very thick] (150:1.0) -- (150:1.55) node[above left, popdark, font=\small] {\ce{H}};
\draw[poporange, very thick] (270:1.0) -- (270:1.55) node[below, popdark, font=\small] {\ce{H}};
\node[popdark, font=\small] at (0,-1.8) {\textbf{Gauche} ($\theta=\SI{60}{\degree}$)};
\node[poporange, font=\small] at (0,-2.2) {Moins stable ($\Delta E \approx \SI{3,8}{kJ.mol^{-1}}$)};
\end{scope}
\end{tikzpicture}
\end{center}
\legende{Conformations décalées du butane selon l'axe \ce{C2-C3}. Anti : les deux \ce{CH3} opposés (stabilité maximale). Gauche : les deux \ce{CH3} voisins (gêne stérique).}
\end{popfigure}

Deux conformations décalées se distinguent particulièrement :
\begin{itemize}
\item la conformation \cle{anti} ($\theta = \SI{180}{\degree}$) : les deux groupes \ce{CH3} sont à l'opposé l'un de l'autre ; la gêne stérique est minimale, c'est la conformation \textbf{la plus stable} (minimum global) ;
\item la conformation \cle{gauche} ($\theta = \SI{60}{\degree}$ ou \SI{300}{\degree}) : les deux \ce{CH3} sont voisins ; la gêne stérique est plus forte, l'énergie est plus élevée d'environ \SI{3,8}{kJ.mol^{-1}} par rapport à l'anti (minimum local).
\end{itemize}

La barrière de rotation la plus élevée du butane (conformation éclipsée avec les deux \ce{CH3} face à face, $\theta = \SI{0}{\degree}$) vaut environ \SI{19}{kJ.mol^{-1}} par rapport à l'anti. Elle reste faible : à température ordinaire, le butane explore toutes ses conformations, en passant le plus clair de son temps dans la forme anti.

\courssub{Le cyclohexane : pourquoi l'hexagone plan est impossible}

Le cyclohexane \ce{C6H12} est un cycle à six atomes de carbone. Si l'on dessinait ce cycle comme un hexagone \emph{plan}, les angles \ce{C-C-C} vaudraient \SI{120}{\degree}. Or un carbone tétraédrique ($sp^3$) exige des angles voisins de \SI{109,5}{\degree}. L'hexagone plan subirait donc une énorme \cle{tension d'angle} (\SI{10,5}{\degree} par angle), sans compter que toutes les liaisons \ce{C-H} seraient éclipsées (gêne de Pitzer). Le cyclohexane plan n'existe pas : le cycle se \emph{plisse} pour retrouver des angles de \SI{109,5}{\degree} et décaler les liaisons.

Deux conformations plissées sont à connaître :

\begin{itemize}
\item la conformation \cle{chaise} : quatre carbones sont coplanaires (le « siège »), un cinquième est au-dessus du plan (le « dossier »), le sixième en dessous (le « pied »). Tous les angles valent \SI{109,5}{\degree} et toutes les liaisons \ce{C-H} sont décalées : c'est la conformation \textbf{stable} (minimum global) ;
\item la conformation \cle{bateau} : deux carbones opposés sont relevés du même côté, comme la proue et la poupe d'une pirogue. Des liaisons \ce{C-H} y sont éclipsées et deux hydrogènes « d'étambot » (sur les carbones relevés) se gênent fortement : c'est une conformation \textbf{instable} (maximum local ou point de passage vers le twist-bateau).
\end{itemize}

\begin{exemplebox}[Écart de stabilité chaise / bateau]
La conformation chaise du cyclohexane est plus stable que la conformation bateau d'environ \SI{25}{kJ.mol^{-1}} à \SI{30}{kJ.mol^{-1}}. À température ordinaire, l'immense majorité des molécules de cyclohexane (plus de \SI{99,9}{\percent}) adopte donc la forme chaise. Retenir l'ordre de grandeur : $\Delta E \approx \SI{28}{kJ.mol^{-1}}$.
\end{exemplebox}

Dans la conformation chaise, les douze liaisons \ce{C-H} ne sont pas équivalentes ; on distingue deux familles de positions :

\begin{itemize}
\item les liaisons \cle{axiales} (notées $a$) : verticales, parallèles à l'axe du cycle, dirigées alternativement vers le haut et vers le bas (6 liaisons) ;
\item les liaisons \cle{équatoriales} (notées $e$) : dirigées vers l'extérieur du cycle, dans le plan moyen, comme les rayons d'une roue (6 liaisons).
\end{itemize}

Un substituant volumineux (par exemple \ce{CH3}, \ce{OH}, \ce{Cl}) préfère toujours la position \emph{équatoriale}, où il dispose de plus d'espace et gêne moins ses voisins (interactions 1,3-diaxiales évitées).

\begin{popfigure}
\begin{center}
\begin{tikzpicture}[scale=1.1, line join=round, line cap=round]
% ===== CHAISE (Vue de côté standard) =====
\begin{scope}
% Coordonnées du cycle chaise (C1 haut, C2 bas, C3 haut, C4 bas, C5 haut, C6 bas)
\coordinate (C1) at (0.0, 1.0);
\coordinate (C2) at (1.0, 0.0);
\coordinate (C3) at (2.0, 1.0);
\coordinate (C4) at (3.0, 0.0);
\coordinate (C5) at (2.0, -1.0);
\coordinate (C6) at (1.0, -1.0);
% Trait du cycle
\draw[popblue, very thick] (C1)--(C2)--(C3)--(C4)--(C5)--(C6)--cycle;
% Liaisons Axiales (verticales)
\draw[poporange, thick] (C1) -- ++(0, 0.7) node[above, poporange, font=\small] {$a$};
\draw[poporange, thick] (C2) -- ++(0, -0.7) node[below, poporange, font=\small] {$a$};
\draw[poporange, thick] (C3) -- ++(0, 0.7) node[above, poporange, font=\small] {$a$};
\draw[poporange, thick] (C4) -- ++(0, -0.7) node[below, poporange, font=\small] {$a$};
\draw[poporange, thick] (C5) -- ++(0, 0.7) node[above, poporange, font=\small] {$a$};
\draw[poporange, thick] (C6) -- ++(0, -0.7) node[below, poporange, font=\small] {$a$};
% Liaisons Équatoriales (vers l'extérieur, ~horizontal)
\draw[popgreen, thick] (C1) -- ++(-0.7, -0.4) node[left, popgreen, font=\small] {$e$};
\draw[popgreen, thick] (C2) -- ++(-0.7, 0.4) node[left, popgreen, font=\small] {$e$};
\draw[popgreen, thick] (C3) -- ++(0.7, -0.4) node[right, popgreen, font=\small] {$e$};
\draw[popgreen, thick] (C4) -- ++(0.7, 0.4) node[right, popgreen, font=\small] {$e$};
\draw[popgreen, thick] (C5) -- ++(0.7, 0.4) node[right, popgreen, font=\small] {$e$};
\draw[popgreen, thick] (C6) -- ++(-0.7, 0.4) node[left, popgreen, font=\small] {$e$};
\node[popdark, font=\small] at (1.5, -1.8) {\textbf{Chaise} (stable)};
\end{scope}
% ===== BATEAU =====
\begin{scope}[xshift=8cm]
\coordinate (B1) at (0.0, 1.0); % Proue haute
\coordinate (B2) at (1.0, 0.2);
\coordinate (B3) at (2.0, 0.2);
\coordinate (B4) at (3.0, 1.0); % Poupe haute
\coordinate (B5) at (2.0, -0.6);
\coordinate (B6) at (1.0, -0.6);
\draw[poppurple, very thick] (B1)--(B2)--(B3)--(B4)--(B5)--(B6)--cycle;
% Hydrogènes d'étambot (flagpole) sur B1 et B4
\draw[poppurple, thick] (B1) -- ++(0, 0.6);
\draw[poppurple, thick] (B4) -- ++(0, 0.6);
% Gêne entre H d'étambot
\draw[popline, dashed, thick] (0,1.6) -- (3,1.6) node[midway, above, font=\small, popmuted] {Gêne forte};
% Quelques liaisons équatoriales/axiales suggérées
\foreach \p/\a/\e in {B2/0/-0.5, B3/0/-0.5, B5/0/0.5, B6/0/0.5}{
  \draw[poporange, thin] (\p) -- ++(0,\a*0.5);
  \draw[popgreen, thin] (\p) -- ++(\e*0.7,0);
}
\node[popdark, font=\small] at (1.5, -1.4) {\textbf{Bateau} (instable)};
\end{scope}
\end{tikzpicture}
\end{center}
\legende{Conformations du cyclohexane. Chaise : liaisons axiales $a$ (orange, verticales alternées) et équatoriales $e$ (vert, vers l'extérieur). Bateau : les deux pointes relevées rapprochent les hydrogènes d'étambot (gêne) et créent des interactions éclipsées.}
\end{popfigure}

\begin{savaistu}[La chaise du cyclohexane est partout dans le vivant]
La forme chaise minimise \emph{à la fois} les tensions d'angle et les répulsions entre liaisons : c'est la conformation « parfaite » d'un cycle à six atomes. La nature l'a adoptée massivement : la molécule de \cle{glucose}, le sucre qui alimente tes cellules, existe presque exclusivement sous forme de cycle à six atomes (pyranose) en conformation chaise, avec la majorité de ses groupes \ce{-OH} en position équatoriale. On retrouve la même architecture dans de nombreux glucides (amidon, cellulose), stéroïdes (cholestérol) et principes actifs de la pharmacopée traditionnelle (écorces de quinquina, racines de moringa). Comprendre la chaise du cyclohexane, c'est comprendre la forme des molécules de la vie.
\end{savaistu}

\begin{exoresolu}[Conformations du butane]
On étudie le butane \ce{CH3-CH2-CH2-CH3} en regardant selon l'axe de la liaison \ce{C2-C3}.
\begin{enumerate}
\item Représenter en projection de Newman la conformation anti.
\item Comparer sa stabilité à celle de la conformation éclipsée correspondante ($\theta = \SI{0}{\degree}$).
\end{enumerate}
\tcblower
\textbf{1.} L'\oe il est placé dans l'axe \ce{C2-C3}. Le carbone avant (C2, point) porte \ce{CH3}, \ce{H}, \ce{H} à \SI{120}{\degree} ; le carbone arrière (C3, cercle) porte également \ce{CH3}, \ce{H}, \ce{H}. Dans la conformation \emph{anti}, l'angle dièdre entre les deux groupes \ce{CH3} vaut $\theta = \SI{180}{\degree}$ : le \ce{CH3} arrière est dessiné à l'opposé du \ce{CH3} avant (en bas si l'avant est en haut), les hydrogènes occupant les positions intermédiaires décalées (\SI{60}{\degree} et \SI{300}{\degree}).

\textbf{2.} La conformation anti est une conformation \emph{décalée} : toutes les liaisons sont décalées deux à deux et, de plus, les deux groupes volumineux \ce{CH3} sont aussi éloignés que possible (\SI{180}{\degree}). La conformation éclipsée correspondante ($\theta = \SI{0}{\degree}$) cumule les répulsions entre doublets de liaison \ce{C-H}/\ce{C-H} et place les deux groupes \ce{CH3} face à face (gêne stérique maximale \ce{CH3}/\ce{CH3}). L'écart d'énergie entre cette forme éclipsée (maximum global) et la forme anti (minimum global) est d'environ \SI{19}{kJ.mol^{-1}}. La conformation anti est donc nettement la plus stable ; à température ordinaire, le butane y séjourne la majeure partie du temps, tout en tournant librement autour de \ce{C2-C3}.
\end{exoresolu}

\begin{aretenir}[L'essentiel sur les conformations]
\begin{itemize}
\item Des conformations sont des formes d'une même molécule qui s'interconvertissent par \textbf{rotation autour d'une liaison simple} ; elles ne sont \textbf{pas isolables}.
\item Représentation de Newman : point = carbone avant, cercle = carbone arrière, substituants à \SI{120}{\degree}.
\item Éthane : éclipsée ($\theta = \SI{0}{\degree}$, énergie maximale) et décalée ($\theta = \SI{60}{\degree}$, énergie minimale) ; barrière de rotation $\approx \SI{12}{kJ.mol^{-1}}$, rotation libre à température ordinaire.
\item Butane selon \ce{C2-C3} : anti (\SI{180}{\degree}, minimum global, la plus stable) et gauche (\SI{60}{\degree}, minimum local, $\Delta E \approx \SI{3,8}{kJ.mol^{-1}}$).
\item Cyclohexane : jamais plan (tension d'angle + éclipsements) ; conformation \textbf{chaise} stable (angles \SI{109,5}{\degree}, tout décalé, positions axiales $a$ et équatoriales $e$) ; conformation \textbf{bateau} instable (éclipsements + gêne étambot, $\Delta E \approx \SI{25}{}$ à \SI{30}{kJ.mol^{-1}}).
\item Un substituant volumineux sur un cyclohexane cherche la position équatoriale.
\end{itemize}
\end{aretenir}

\coursec{Importance de la stéréo-isomérie}

Dans la section précédente, nous avons vu qu'une molécule n'est pas un objet figé : par rotation autour des liaisons simples, elle adopte des \cle{conformations} différentes (éclipsée, décalée, chaise, bateau). Mais ces conformations restent \emph{interconvertibles} : il s'agit toujours de la même molécule. Nous abordons maintenant la question qui donne tout son sens à la stéréochimie : lorsque deux stéréo-isomères sont de \emph{vraies molécules distinctes} — comme deux énantiomères —, cela a-t-il des conséquences concrètes ? La réponse est immense : c'est une question de vie ou de mort, de médicament ou de poison, d'odeur de menthe ou d'odeur de carvi. Cette section te montrera pourquoi la stéréochimie est devenue l'un des piliers de l'industrie pharmaceutique moderne.

\courssub{Pourquoi le vivant est-il sensible à la stéréochimie ?}

\textbf{L'intuition : la main dans le gant.}\par
Imagine que tu cherches à enfiler un gant droit. Ta main droite s'y glisse parfaitement ; ta main gauche, pourtant constituée des mêmes « éléments » (cinq doigts, une paume), n'y entre pas correctement. Le gant est un objet \cle{chiral}, et il distingue sans effort deux objets images l'un de l'autre dans un miroir.

Or le monde vivant est exactement comme ce gant : il est \emph{chiral}. Les récepteurs de nos cellules, les enzymes qui catalysent nos réactions, les récepteurs olfactifs de notre nez sont des protéines, c'est-à-dire de longues chaînes d'\cle{acides aminés}. Et voici le fait fondamental :

\begin{propriete}[La chiralité du vivant]
Les acides aminés naturels qui constituent les protéines sont (presque) tous de la même famille stéréochimique : la \cle{série L}. Il en résulte que toutes les macromolécules biologiques (enzymes, récepteurs, anticorps) sont chirales et possèdent des sites actifs à la géométrie tridimensionnelle bien définie. Un tel site ne « reconnaît » en général qu'\emph{un seul} des deux énantiomères d'une molécule, comme une serrure ne reconnaît qu'une seule clé.
\end{propriete}

C'est le fameux modèle de la \emph{clé et de la serrure}, imaginé par Emil Fischer (le même que la projection de Fischer !) à la fin du XIX\textsuperscript{e} siècle.

\begin{popfigure}
\begin{center}
\begin{tikzpicture}[scale=1.0]
% Serrure enzymatique (site actif) - Haut
\draw[very thick, popdark, fill=popblueL] (0,0) -- (0,2.4) -- (0.7,2.4) -- (0.7,1.5) -- (1.5,1.5) -- (1.5,2.4) -- (2.4,2.4) -- (2.4,0) -- cycle;
\node[popdark] at (1.2,-0.4) {\small \textbf{Site actif de l'enzyme (chiral)}};
% Clé R qui s'insère
\begin{scope}[xshift=-3.4cm, yshift=1.05cm]
\draw[very thick, popgreen, fill=popgreenL] (0,0) rectangle (0.7,0.9);
\draw[very thick, popgreen, fill=popgreenL] (0.7,0.28) rectangle (1.7,0.62);
\node[popgreen!60!black] at (0.35,1.15) {\small \textbf{Énantiomère actif}};
\draw[-{Stealth}, very thick, popgreen!60!black] (1.9,0.45) -- (2.9,0.45);
\node[popgreen!60!black] at (2.4,0.75) {\small s'insère};
\end{scope}
% Clé S qui ne s'insère pas
\begin{scope}[xshift=-3.4cm, yshift=-1.9cm]
\draw[very thick, poporange, fill=poporangeL] (0,0) rectangle (0.7,0.9);
\draw[very thick, poporange, fill=poporangeL] (0.7,0.05) rectangle (1.7,0.39);
\node[poporange!80!black] at (0.35,1.15) {\small \textbf{Énantiomère inactif}};
\draw[-{Stealth}, very thick, poporange!80!black] (1.9,0.45) -- (2.7,0.45);
\draw[very thick, poppink] (2.75,0.2) -- (3.05,0.7);
\draw[very thick, poppink] (3.05,0.2) -- (2.75,0.7);
\node[poppink] at (2.4,0.85) {\small bloqué};
\end{scope}
% Même serrure en bas
\draw[very thick, popdark, fill=popblueL] (0,-2.4) -- (0,-0.85) -- (0.7,-0.85) -- (0.7,-1.75) -- (1.5,-1.75) -- (1.5,-0.85) -- (2.4,-0.85) -- (2.4,-2.4) -- cycle;
\end{tikzpicture}
\end{center}
\legende{Modèle clé--serrure : le site actif chiral de l'enzyme n'accepte que l'énantiomère dont la géométrie est complémentaire. L'autre énantiomère, image dans un miroir, ne peut pas se fixer correctement.}
\end{popfigure}

\begin{attention}[Les énantiomères n'ont pas TOUTES leurs propriétés différentes]
Erreur très fréquente au Baccalauréat : croire que deux énantiomères diffèrent par toutes leurs propriétés. C'est faux ! Deux énantiomères ont des \textbf{propriétés physiques identiques} : même température de fusion, même température d'ébullition, même densité, même solubilité dans un solvant achiral, même spectre d'absorption. Ils ne se distinguent que par :
\begin{itemize}
\item le \cle{pouvoir rotatoire} : ils dévient la lumière polarisée de la même valeur absolue mais en sens opposés (dextrogyre / lévogyre) ;
\item leur \cle{activité biologique} : ils interagissent différemment avec les récepteurs chiraux du vivant (goût, odeur, effet thérapeutique, toxicité).
\end{itemize}
C'est précisément parce que leurs propriétés physiques sont identiques que leur \emph{séparation} est si difficile en industrie : impossible de les séparer par simple distillation ou cristallisation classique !
\end{attention}

\courssub{Le drame du thalidomide : la leçon la plus tragique de la stéréochimie}

\textbf{L'histoire.}\par
À la fin des années 1950, un laboratoire allemand commercialise le \cle{thalidomide}, un sédatif très efficace contre l'insomnie et les nausées des femmes enceintes. Le médicament est vendu sous forme de \cle{mélange racémique}, c'est-à-dire contenant les deux énantiomères en proportions égales. Personne, à l'époque, ne s'en inquiète : on pense que deux énantiomères, ayant les mêmes propriétés physiques, ont le même effet biologique.

C'était une erreur dramatique. L'énantiomère $(R)$ possède bien l'effet \emph{sédatif} recherché, mais l'énantiomère $(S)$ est \emph{tératogène} : il perturbe le développement du fœtus. Entre 1957 et 1961, environ \num{10000} enfants naissent dans le monde avec des malformations graves (membres atrophiés ou absents), avant que le médicament ne soit retiré du marché. Pire encore : même si l'on avait administré le seul énantiomère $(R)$, l'organisme le convertit partiellement en $(S)$ \emph{in vivo}. Le drame était donc presque inévitable avec cette molécule.

\begin{definition}[Médicament racémique]
Un médicament \cle{racémique} est un médicament qui contient les deux énantiomères d'une molécule chirale en proportions égales (mélange racémique, optiquement inactif). Très souvent, \emph{un seul} des deux énantiomères est thérapeutiquement actif ; l'autre est inactif, voire toxique. Un médicament ne contenant que l'énantiomère actif est dit \emph{énantiopure}.
\end{definition}

\begin{savaistu}[Une révolution réglementaire née d'un drame]
Après la catastrophe du thalidomide, les agences du médicament (FDA aux États-Unis, puis leurs homologues européennes) ont profondément durci leurs exigences. Depuis les années 1990, toute molécule chirale candidate à la mise sur le marché doit faire l'objet d'une \textbf{étude séparée de chaque énantiomère} : activité, toxicité, devenir dans l'organisme. Cette réglementation a transformé l'industrie pharmaceutique : la maîtrise de la stéréochimie est devenue un enjeu économique et sanitaire majeur. Le prix Nobel de chimie 2001 a d'ailleurs récompensé Knowles, Noyori et Sharpless pour leurs travaux sur la \emph{catalyse asymétrique}, qui permet de fabriquer directement l'énantiomère désiré.
\end{savaistu}

\courssub{L'ibuprofène : un racémique que le corps « répare »}

L'\cle{ibuprofène} est l'un des antalgiques et anti-inflammatoires les plus consommés au monde — on le trouve dans toutes les pharmacies de Yaoundé ou de Douala. Sa molécule possède un carbone asymétrique, donc deux énantiomères $(R)$ et $(S)$.

\begin{itemize}
\item Seul l'énantiomère $(S)$ inhibe efficacement l'enzyme \emph{cyclooxygénase} (COX), responsable de la production des médiateurs de la douleur et de l'inflammation.
\item L'énantiomère $(R)$ est quasiment inactif sur cette cible.
\end{itemize}

\begin{exemplebox}[Une activité 100 fois supérieure]
Des études pharmacologiques montrent que l'énantiomère $(S)$ de l'ibuprofène est environ \textbf{100 fois plus actif} que l'énantiomère $(R)$ sur l'inhibition de la cyclooxygénase. Autrement dit, pour obtenir le même effet antalgique avec l'énantiomère $(R)$ pur, il faudrait une dose environ 100 fois plus élevée !
\end{exemplebox}

Alors pourquoi l'ibuprofène est-il vendu sous forme racémique ? Parce que l'organisme possède une enzyme (une isomérase) qui \emph{convertit partiellement} l'énantiomère $(R)$ inactif en énantiomère $(S)$ actif. Environ 60 \% de la dose de $(R)$ est ainsi « recyclée » en $(S)$. Le mélange racémique, beaucoup moins coûteux à produire, reste donc efficace. Certains fabricants commercialisent néanmoins le \emph{dexibuprofène}, qui ne contient que l'énantiomère $(S)$ : à effet égal, la dose est plus faible et l'estomac moins sollicité.

\begin{exoresolu}[Dose efficace d'ibuprofène]
Un comprimé d'ibuprofène racémique contient \SI{400}{mg} de principe actif, soit \SI{200}{mg} de chaque énantiomère. On admet que l'organisme convertit 60 \% de l'énantiomère $(R)$ en $(S)$, et que l'activité du $(R)$ restant est négligeable.
\begin{enumerate}
\item Calculer la masse d'énantiomère $(S)$ réellement disponible dans l'organisme après conversion.
\item Quelle masse de dexibuprofène (énantiomère $(S)$ pur) faudrait-il pour obtenir le même effet ?
\end{enumerate}
\tcblower
\textbf{1.} Masse de $(R)$ convertie en $(S)$ :
\[ m_{\text{convertie}} = 0,60 \times 200 = \SI{120}{mg} \]
Masse totale de $(S)$ disponible :
\[ m_S = 200 + 120 = \SI{320}{mg} \]
\textbf{2.} Comme seul l'énantiomère $(S)$ est actif, il suffirait d'un comprimé de dexibuprofène de \SI{320}{mg} pour obtenir le même effet thérapeutique qu'un comprimé racémique de \SI{400}{mg}. On comprend l'intérêt commercial du médicament énantiopure : dose réduite de 20 \%, donc moins d'effets secondaires.
\end{exoresolu}

\courssub{Les arômes : quand le nez fait de la stéréochimie}

Nos récepteurs olfactifs sont des protéines chirales : ils distinguent les énantiomères comme le gant distingue les mains. Le résultat est étonnant : \emph{deux énantiomères peuvent avoir des odeurs totalement différentes}.

\textbf{La carvone.}\par
La \cle{carvone} est une molécule à un carbone asymétrique, présente dans plusieurs huiles essentielles :
\begin{itemize}
\item l'énantiomère $(R)$ de la carvone a une odeur de \textbf{menthe} (c'est le principal constituant odorant de l'huile de menthe verte) ;
\item l'énantiomère $(S)$ a une odeur de \textbf{carvi} (une épice au parfum anisé, proche de l'aneth).
\end{itemize}
Même formule brute, même formule développée, même température d'ébullition... et pourtant deux odeurs que personne ne confondrait !

\begin{popfigure}
\begin{center}
\setchemfig{atom sep=1.8em}
% Enantiomère (R) - Carvone : wedge sur le groupe isopropényle (liaison sortante)
\chemfig{[:-30]*6(=O-(-CH_3)=--CH(-[1]>[:0]CH_2=C(-[0]CH_3)-)-=)}
\hspace{3em}
% Enantiomère (S) - Carvone : dash sur le groupe isopropényle (liaison rentrante)
\chemfig{[:-30]*6(=O-(-CH_3)=--CH(-[1]<[:0]CH_2=C(-[0]CH_3)-)-=)}
\end{center}
\legende{Les deux énantiomères de la carvone : à gauche, l'énantiomère $(R)$ (liaison wedge vers l'isopropényle) à l'odeur de \textbf{menthe} ; à droite, l'énantiomère $(S)$ (liaison dash vers l'isopropényle) à l'odeur de \textbf{carvi}. Le carbone asymétrique est le C5 du cycle (portant le H implicite et le groupe isopropényle).}
\end{popfigure}

\textbf{Le limonène.}\par
Même phénomène pour le \cle{limonène}, molécule des écorces d'agrumes : l'énantiomère $(R)$ sent l'\textbf{orange}, l'énantiomère $(S)$ sent le \textbf{citron} (ou la térébenthine). La prochaine fois que tu presses un citron dans ta cuisine, souviens-toi que ton nez est en train de réaliser une analyse stéréochimique d'une précision remarquable, qu'aucun appareil de laboratoire simple ne saurait reproduire !

\begin{popfigure}
\begin{center}
\setchemfig{atom sep=1.8em}
% Enantiomère (R) - Limonène : wedge sur l'isopropényle en C4
\chemfig{[:-30]*6((-CH_3)=--CH(-[1]>[:0]CH_2=C(-[0]CH_3)-)--)}
\hspace{3em}
% Enantiomère (S) - Limonène : dash sur l'isopropényle en C4
\chemfig{[:-30]*6((-CH_3)=--CH(-[1]<[:0]CH_2=C(-[0]CH_3)-)--)}
\end{center}
\legende{Les deux énantiomères du limonène : à gauche, l'énantiomère $(R)$ (odeur d'\textbf{orange}, dextrogyre) ; à droite, l'énantiomère $(S)$ (odeur de \textbf{citron}, lévogyre). Le carbone asymétrique est le C4 du cycle.}
\end{popfigure}

\begin{methode}[Raisonner face à un problème de stéréochimie et d'activité biologique]
Devant un exercice ou une situation concrète, adopte la démarche suivante :
\begin{enumerate}
\item \textbf{Repérer} les carbones asymétriques (carbone tétragonal portant quatre groupes différents) ou l'axe de stéréo-isomérie (double liaison $Z/E$).
\item \textbf{Nommer} les stéréo-isomères possibles : $2^n$ au maximum pour $n$ carbones asymétriques.
\item \textbf{Classer} les paires d'isomères : images dans un miroir non superposables $\Rightarrow$ énantiomères ; sinon $\Rightarrow$ diastéréoisomères.
\item \textbf{Prévoir} les conséquences : propriétés physiques identiques pour les énantiomères, mais pouvoirs rotatoires opposés et activités biologiques potentiellement très différentes.
\end{enumerate}
\end{methode}

\courssub{Conséquences industrielles : synthèse asymétrique et séparation des énantiomères}

Tout ce qui précède se traduit par un défi industriel considérable : comment produire \emph{un seul} énantiomère ?

\textbf{Première stratégie : la séparation (dédoublement).}\par
On synthétise le mélange racémique, puis on sépare les deux énantiomères. Mais attention au piège signalé plus haut : leurs propriétés physiques étant identiques, les techniques classiques (distillation, cristallisation simple) échouent. On utilise alors des méthodes \emph{chirales} :
\begin{itemize}
\item faire réagir le racémique avec un réactif énantiopur pour former deux \emph{diastéréoisomères}, qui eux ont des propriétés physiques différentes et peuvent être séparés (méthode historique de Pasteur avec les sels d'acide tartrique) ;
\item la \emph{chromatographie chirale}, dont la colonne contient un support chiral qui retient un énantiomère plus longtemps que l'autre ;
\item la \emph{résolution enzymatique} : une enzyme ne transforme qu'un seul énantiomère, ce qui permet de l'éliminer sélectivement.
\end{itemize}

\textbf{Deuxième stratégie : la synthèse asymétrique.}\par
Plus élégante et plus économique, elle consiste à fabriquer directement l'énantiomère désiré, en utilisant un \emph{catalyseur chiral} qui oriente la réaction vers un seul stéréo-isomère. Schématiquement, au lieu de :
\[ \ce{\text{substrat achiral} ->[\text{catalyseur classique}] (R) + (S) \quad (50/50)} \]
on réalise :
\[ \ce{\text{substrat achiral} ->[\text{catalyseur chiral}] (S) \quad (\text{quasi exclusivement})} \]
C'est le principe des travaux récompensés par le prix Nobel 2001. Aujourd'hui, une grande partie des médicaments chiraux (certains antibiotiques, antihypertenseurs, antiviraux) est produite par synthèse asymétrique ou par voie enzymatique.

\begin{attention}[Pièges à éviter dans cette section]
\begin{itemize}
\item Ne dis jamais qu'un mélange racémique est « inactif » : il est \emph{optiquement} inactif (les pouvoirs rotatoires se compensent), mais il peut être \emph{biologiquement} actif (ibuprofène !).
\item Ne confonds pas « énantiomère inactif » et « énantiomère toxique » : dans le cas de l'ibuprofène, le $(R)$ est simplement peu actif ; dans le cas du thalidomide, le $(S)$ est dangereux. Chaque cas est particulier.
\item La conversion d'un énantiomère en l'autre dans l'organisme (comme pour l'ibuprofène ou le thalidomide) montre qu'administrer un énantiomère pur ne garantit pas toujours l'absence de l'autre \emph{in vivo}.
\end{itemize}
\end{attention}

\begin{aretenir}[L'essentiel de la section]
\begin{itemize}
\item Le vivant est chiral : acides aminés de série L, enzymes et récepteurs chiraux $\Rightarrow$ les énantiomères sont reconnus différemment (modèle clé--serrure).
\item Deux énantiomères ont les \textbf{mêmes propriétés physiques} mais des \textbf{pouvoirs rotatoires opposés} et des \textbf{activités biologiques différentes}.
\item Thalidomide : un énantiomère sédatif, l'autre tératogène $\Rightarrow$ durcissement mondial des contrôles pharmaceutiques.
\item Ibuprofène : seul le $(S)$ est actif (environ 100 fois plus que le $(R)$) ; le racémique est vendu car le corps convertit partiellement $(R)$ en $(S)$.
\item Carvone (menthe / carvi) et limonène (orange / citron) : nos récepteurs olfactifs sont chiraux.
\item Enjeu industriel majeur : séparer les énantiomères (dédoublement) ou fabriquer directement le bon (synthèse asymétrique, catalyseurs chiraux).
\end{itemize}
\end{aretenir}

\begin{exercice}[Pour t'entraîner]
L'énantiomère $(R)$ du limonène, d'odeur d'orange, est dextrogyre : son pouvoir rotatoire spécifique vaut $[\alpha] = +125$ (unités usuelles du polarimètre).
\begin{enumerate}
\item Que vaut le pouvoir rotatoire spécifique de l'énantiomère $(S)$, d'odeur de citron ? Justifier.
\item On mesure la rotation de la lumière polarisée par un échantillon de limonène extrait d'écorces d'orange et on trouve une déviation nulle. Que peut-on conclure sur la composition de l'échantillon ?
\item Ces deux énantiomères peuvent-ils être séparés par distillation fractionnée ? Justifier.
\end{enumerate}
\end{exercice}

\begin{corrige}[Exercice : limonène]
\textbf{1.} Deux énantiomères ont des pouvoirs rotatoires de même valeur absolue et de signes opposés : l'énantiomère $(S)$ a donc $[\alpha] = -125$ (lévogyre).

\textbf{2.} Une déviation nulle signifie que l'échantillon est optiquement inactif : il s'agit soit d'un mélange racémique (50 \% de $(R)$ et 50 \% de $(S)$, les effets se compensent exactement), soit éventuellement d'un mélange dont les proportions conduiraient à une compensation — mais le cas naturel et attendu ici est le mélange racémique. L'odeur de l'échantillon serait d'ailleurs intermédiaire entre l'orange et le citron.

\textbf{3.} Non : les deux énantiomères ont exactement la même température d'ébullition (propriété physique identique), ils s'évaporent donc ensemble. Il faut une méthode chirale : chromatographie chirale, dédoublement par formation de diastéréoisomères, ou résolution enzymatique.
\end{corrige}

\begin{savaistu}[Et au quotidien ?]
La stéréochimie est partout : l'aspartame (édulcorant des boissons « light ») est sucré, alors que son énantiomère est amer ; le glucose naturel $(D)$ est assimilé par notre organisme, mais son image miroir $(L)$ ne l'est pas — on pourrait manger du sucre $(L)$ sans grossir, s'il n'était pas si cher à fabriquer ! Même les arômes des jus de fruits que l'on déguste à l'heure de la récréation (bissap, ananas, mangue) doivent leur parfum à des molécules dont la géométrie dans l'espace fait toute la différence. La prochaine fois qu'un parfum te plaît, remercie la chiralité !
\end{savaistu}

\newpage

\coursec{Travaux pratiques}

\courssub{Objectifs du TP}

À l'issue de cette séance de travaux pratiques, tu seras capable de :

\begin{itemize}
\item mettre en évidence expérimentalement l'\cle{activité optique} d'une solution de glucose à l'aide d'un polarimètre artisanal ;
\item mesurer l'\cle{angle de rotation} $\alpha$ du plan de polarisation de la lumière pour deux concentrations différentes ;
\item vérifier qualitativement la \cle{loi de Biot} : $\alpha$ est proportionnel à la concentration $c$ de la solution ;
\item construire, à l'aide de modèles moléculaires, les deux \cle{énantiomères} du butan-2-ol, montrer qu'ils ne sont pas superposables et les représenter en projection de Fischer ;
\item matérialiser les \cle{conformations} éclipsée et décalée de l'éthane et les représenter en projection de Newman ;
\item réaliser les conformations \cle{chaise} et \cle{bateau} du cyclohexane et comparer leur stabilité.
\end{itemize}

Ce TP illustre directement la famille de situations \og choix d'un stéréo-isomère, comparaison de deux énantiomères \fg{} : tu vas \emph{voir} de tes propres yeux qu'une molécule chirale dévie la lumière polarisée, et \emph{toucher} de tes propres mains ce qu'est la chiralité.

\courssub{Matériel et produits}

\begin{center}
\rowcolors{2}{popblueL}{white}
\begin{tabularx}{\linewidth}{|l|X|}
\hline
\rowcolor{popblue!40} \textbf{Matériel / produit} & \textbf{Alternative locale si indisponible} \\
\hline
Deux filtres polarisants (polariseur et analyseur) & Verres de lunettes 3D de cinéma récupérés, ou deux verres de lunettes de soleil polarisées \\
\hline
Lampe torche ou lampe de poche à LED & Lampe de téléphone portable, faisceau bien dirigé \\
\hline
Tube en verre ou cuvette transparente de longueur $\ell = \SI{10}{cm}$ & Tube à essai large, éprouvette, ou cristallisoir en verre propre (mesurer la longueur traversée par la lumière) \\
\hline
Glucose en poudre (D-glucose) & Sucre de table (saccharose) du marché, bien dissous et filtré ; noter que le pouvoir rotatoire diffère \\
\hline
Balance au centième de gramme & Balance de ménagère ou de boutique ; à défaut, doser avec une cuillère à café étalonnée ($\approx \SI{5}{g}$) \\
\hline
Fioles jaugées de \SI{100}{mL}, béchers, agitateur en verre & Bouteilles en plastique propres graduées au marqueur à l'aide d'une éprouvette \\
\hline
Rapporteur, feuilles blanches, marqueur, papier calque & Rapporteur d'écolier classique \\
\hline
Modèles moléculaires (boules et tiges) & Boules de pâte à modeler de couleurs différentes (noir = C, blanc = H, rouge = O) et allumettes ou cure-dents comme liaisons \\
\hline
\end{tabularx}
\end{center}

\courssub{Sécurité}

\begin{attention}[Consignes de sécurité]
\begin{itemize}
\item \textbf{Pictogramme \og matériel en verre \fg{}} : le tube en verre et les fioles sont fragiles. Manipule-les au-dessus de la paillasse, jamais au-dessus du sol. En cas de casse, ne ramasse jamais les éclats à mains nues : utilise une pelle et une brosse.
\item \textbf{Pictogramme \og rayonnement lumineux \fg{}} : ne dirige jamais le faisceau de la lampe vers les yeux d'un camarade, même si la lumière d'une torche est peu intense.
\item Le glucose et le saccharose sont des produits \textbf{non dangereux} (comestibles), mais en chimie on adopte le réflexe professionnel : \emph{on ne goûte jamais un produit du laboratoire}, même du sucre.
\item Porte une blouse ou un tablier, attache les cheveux longs, et garde la paillasse dégagée.
\item La pâte à modeler ne se mange pas ; lave-toi les mains à la fin de la séance.
\end{itemize}
\end{attention}

\courssub{Schéma du dispositif}

\begin{popfigure}
\begin{center}
\begin{tikzpicture}[scale=1.0, >=Stealth]
% source lumineuse
\draw[fill=popgold!40, draw=popdark] (0,0.4) rectangle (0.8,1.6);
\draw[popdark] (0.4,1.6) -- (0.4,1.9);
\node[below, align=center, font=\small] at (0.4,0.4) {Source\\lumineuse};
% rayons avant polariseur
\draw[->, poporange, thick] (0.8,0.8) -- (1.7,0.8);
\draw[->, poporange, thick] (0.8,1.2) -- (1.7,1.2);
\node[above, font=\small, text=poporange] at (1.3,1.35) {lumière naturelle};
% polariseur
\draw[fill=popblue!50, draw=popdark] (1.8,0.2) rectangle (2.1,1.8);
\node[below, align=center, font=\small] at (1.95,0.2) {Polariseur\\(fixe)};
% lumière polarisée (flèches verticales)
\draw[->, poppurple, thick] (2.1,1.0) -- (3.4,1.0);
\foreach \x in {2.4, 2.8, 3.2} {
    \draw[poppurple, thick] (\x,0.7) -- (\x,1.3);
}
\node[above, font=\small, text=poppurple] at (2.75,1.4) {lumière polarisée plane};
% cuve
\draw[fill=popgreenL, draw=popdark, thick] (3.5,0.3) rectangle (6.5,1.7);
\draw[<->, popdark] (3.5,2.0) -- (6.5,2.0) node[midway, above, font=\small] {$\ell$};
\node[below, align=center, font=\small] at (5.0,0.3) {Cuve : solution de glucose\\(substance active)};
% hélice pour rotation
\draw[poppurple, thick, decorate, decoration={coil, aspect=0.4, segment length=6pt, amplitude=4pt}] (3.6,1.0) -- (6.4,1.0);
% analyseur
\draw[fill=popblue!50, draw=popdark] (6.8,0.2) rectangle (7.1,1.8);
\node[below, align=center, font=\small] at (6.95,0.2) {Analyseur\\(mobile)};
\draw[->, poppink, thick] (6.95,2.1) arc[start angle=90, end angle=30, radius=0.7];
\node[above right, font=\small, text=poppink] at (7.3,2.3) {rotation d'angle $\alpha$};
% oeil
\draw[->, poppurple, thick] (7.1,1.0) -- (8.0,1.0);
\draw[popdark] (8.3,1.0) ellipse (0.28 and 0.4);
\node[below, font=\small] at (8.3,0.55) {Observateur};
\end{tikzpicture}
\end{center}
\legende{Schéma du polarimètre artisanal : la lumière traverse le polariseur, puis la cuve de solution de glucose qui fait tourner le plan de polarisation d'un angle $\alpha$ ; on le mesure en tournant l'analyseur jusqu'à l'extinction.}
\end{popfigure}

\courssub{Protocole}

\textbf{Partie 1 : polarimétrie artisanale (activité optique et loi de Biot)}

\begin{enumerate}
\item \textbf{Préparation des solutions.} Prépare \SI{100,0}{mL} d'une solution $S_1$ en dissolvant $m_1 = \SI{20,0}{g}$ de glucose (D-glucose) dans de l'eau distillée (ou de l'eau bouillie refroidie). Filtre si la solution est trouble : toute suspension diffuse la lumière et gêne l'observation. Prépare ensuite la solution $S_2$ en diluant $S_1$ deux fois : prélève \SI{50,0}{mL} de $S_1$ (pipette jaugée ou seringue graduée) et complète à \SI{100,0}{mL} avec de l'eau. On obtient $c_1 = \SI{200}{g.L^{-1}}$ ($= \SI{0,200}{g.mL^{-1}}$) et $c_2 = \SI{100}{g.L^{-1}}$ ($= \SI{0,100}{g.mL^{-1}}$).
\item \textbf{Réglage du zéro (blanc).} Place la cuve remplie d'\emph{eau pure} entre les deux filtres. Tourne l'analyseur jusqu'à obtenir l'\cle{extinction} (obscurité maximale du champ lumineux). Repère la position de l'analyseur sur le rapporteur fixé derrière lui : c'est la position $0\degree$ (zéro de l'appareil).
\item \textbf{Mesure avec $S_1$.} Remplace l'eau par la solution $S_1$ : le champ s'éclaire, car le plan de polarisation a tourné. Tourne lentement l'analyseur dans le sens qui rétablit l'extinction. Note l'angle $\alpha_1$ dont tu as dû tourner l'analyseur, ainsi que le \textbf{sens} de rotation (vers la droite de l'observateur = sens horaire : substance \cle{dextrogyre}, notée $+$ ; vers la gauche = sens trigonométrique : substance \cle{lévogyre}, notée $-$).
\item \textbf{Mesure avec $S_2$.} Recommence avec la solution $S_2$ et note $\alpha_2$ (même sens de rotation).
\item \textbf{Reproductibilité.} Refais chaque mesure trois fois (vidange, rinçage à l'eau, rinçage avec la solution, remplissage) et calcule la moyenne $\bar{\alpha}_1$ et $\bar{\alpha}_2$.
\end{enumerate}

\textbf{Partie 2 : modèles moléculaires (chiralité, conformations, projection de Fischer)}

\begin{enumerate}
\item \textbf{Énantiomères du butan-2-ol (perspective de Cram).} Construis une molécule de butan-2-ol \chemfig{CH_3-CH(-[2]OH)-CH_2-CH_3} : le carbone n°2 porte quatre groupes différents (\ce{H}, \ce{OH}, \ce{CH3}, \ce{CH2CH3}) : c'est un \cle{carbone asymétrique} (noté C*). Construis ensuite son image dans un miroir (inverse deux liaisons ou retourne le modèle). Essaie de superposer les deux modèles en les tournant dans tous les sens, \emph{sans casser aucune liaison}.
\item \textbf{Projection de Fischer du butan-2-ol.} Sur une feuille, trace deux croix (représentant le C*). Rappel : les liaisons horizontales sortent du plan (vers toi), les liaisons verticales rentrent dans le plan. Place les 4 groupes sur le premier modèle (réel) : \ce{CH3} en haut, \ce{CH2CH3} en bas, \ce{H} à gauche, \ce{OH} à droite (par exemple). Reporte cette configuration sur la première croix. Fais de même pour l'image miroir sur la deuxième croix. Vérifie qu'on ne peut pas passer de l'une à l'autre par une rotation de $180\degree$ dans le plan de la feuille.
\item \textbf{Conformations de l'éthane (projection de Newman).} Construis une molécule d'éthane \ce{CH3-CH3}. Fais tourner un groupe méthyle autour de la liaison \ce{C-C} : passe de la conformation \cle{éclipsée} (les hydrogènes se font face, répulsion maximale) à la conformation \cle{décalée} (rotation de $60\degree$, répulsion minimale). Dessine pour chacune la représentation de Newman correspondante (cercle = carbone avant, traits = liaisons du carbone arrière).
\item \textbf{Conformations du cyclohexane (chaise et bateau).} Avec le kit (ou 6 boules noires + 12 blanches + liaisons), construis le cycle en conformation \cle{chaise} (alternance haut/bas des carbones, toutes les liaisons C-H décalées). Puis déforme-le pour obtenir la conformation \cle{bateau} (deux carbones opposés relevés, interactions de drapeau). Compare la facilité de construction et la place occupée.
\end{enumerate}

\courssub{Observations et résultats attendus}

\begin{itemize}
\item Avec l'eau pure, l'extinction est obtenue sans rotation supplémentaire : l'eau est \textbf{optiquement inactive}.
\item Avec les solutions de glucose, il faut tourner l'analyseur \textbf{vers la droite} (sens horaire vu de l'observateur) : le D-glucose est \cle{dextrogyre} ($\alpha > 0$). (Le saccharose est aussi dextrogyre, $[\alpha]_D = +66,5\degree$).
\item L'angle mesuré avec $S_1$ est environ le \textbf{double} de celui mesuré avec $S_2$.
\item Partie 2 : les deux modèles du butan-2-ol \textbf{ne sont pas superposables}, quelle que soit la rotation effectuée : ce sont deux énantiomères. Leurs projections de Fischer sont images l'une de l'autre par rapport à un axe vertical (échange gauche/droite).
\item Pour l'éthane, la conformation éclipsée est plus \og tendue \fg{} (gêne stérique entre les hydrogènes) que la conformation décalée (plus stable).
\item Pour le cyclohexane, la conformation chaise est rigide et sans tension (angles \SI{109,5}{\degree}, liaisons décalées) ; la conformation bateau présente des interactions de drapeau (H-H trop proches) et des angles déformés : elle est moins stable d'environ \SI{23}{kJ.mol^{-1}}.
\end{itemize}

\begin{center}
\rowcolors{2}{popgreenL}{white}
\begin{tabularx}{\linewidth}{|l|X|X|X|}
\hline
\rowcolor{popgreen!40} \textbf{Solution} & \textbf{Concentration $c$ (\si{g.L^{-1}})} & \textbf{Concentration $c$ (\si{g.mL^{-1}})} & \textbf{Angle $\alpha$ mesuré (moyenne)} \\
\hline
Eau pure & 0 & 0,000 & $0,0\degree$ \\
\hline
$S_1$ (glucose) & 200 & 0,200 & $+10,5\degree$ \\
\hline
$S_2$ (glucose dilué) & 100 & 0,100 & $+5,3\degree$ \\
\hline
\end{tabularx}
\end{center}

\courssub{Exploitation}

\begin{exoresolu}[Vérification de la loi de Biot et dosage]
\textbf{Énoncé.} (1) Rappeler la loi de Biot en précisant la signification et les unités usuelles de chaque grandeur. (2) À l'aide des mesures du tableau, montrer que $\alpha$ est proportionnel à $c$. (3) En déduire le pouvoir rotatoire spécifique $[\alpha]_D^{20}$ du glucose (longueur d'onde raie D du sodium, $T = \SI{20}{\celsius}$). (4) Une solution inconnue de glucose donne une rotation $\alpha = +7,9\degree$ dans la même cuve. Quelle est sa concentration en \si{g.L^{-1}} ?
\tcblower
\textbf{Solution.}
\begin{enumerate}
\item La loi de Biot s'écrit :
\[ \alpha = [\alpha]_\lambda^T \times \ell \times c \]
où :
\begin{itemize}
\item $\alpha$ : angle de rotation du plan de polarisation (en degrés, $\degree$) ; $+$ pour dextrogyre, $-$ pour lévogyre.
\item $[\alpha]_\lambda^T$ : \cle{pouvoir rotatoire spécifique} (en $\degree \cdot \si{dm^{-1}} \cdot \si{mL.g^{-1}}$ ou $\degree \cdot \si{dm^{-1}} \cdot \si{cm^3.g^{-1}}$). C'est une constante caractéristique de la substance chirale, pour une longueur d'onde $\lambda$ (souvent raie D, $\SI{589}{nm}$) et une température $T$ (souvent \SI{20}{\celsius}).
\item $\ell$ : longueur de la cuve traversée par la lumière (en décimètres, \si{dm}).
\item $c$ : concentration massique de la substance active (en \si{g.mL^{-1}}, soit \si{g.cm^{-3}}). \textbf{Attention :} l'unité légale du SI pour $c$ dans cette formule est \si{g.mL^{-1}} (et non \si{g.L^{-1}}).
\end{itemize}
\item Vérification de la proportionnalité. Calculons le rapport $\dfrac{\alpha}{c}$ pour chaque solution (avec $c$ en \si{g.mL^{-1}}) :
\[ \frac{\alpha_1}{c_1} = \frac{10,5}{0,200} = \SI{52,5}{\degree.mL.g^{-1}} \qquad ; \qquad \frac{\alpha_2}{c_2} = \frac{5,3}{0,100} = \SI{53,0}{\degree.mL.g^{-1}} \]
Les deux rapports sont égaux aux erreurs expérimentales près ($\approx 1\%$) : $\alpha$ est bien \textbf{proportionnel} à $c$, ce qui vérifie la loi de Biot.
\item Détermination de $[\alpha]_D^{20}$. Avec $\ell = \SI{10,0}{cm} = \SI{1,00}{dm}$ :
\[ [\alpha]_D^{20} = \frac{\alpha}{\ell \times c} = \frac{10,5}{1,00 \times 0,200} = \SI{+52,5}{\degree.dm^{-1}.mL.g^{-1}} \]
La valeur de référence pour le D-glucose (équilibre $\alpha/\beta$) est $[\alpha]_D^{20} = \SI{+52,7}{\degree.dm^{-1}.mL.g^{-1}}$. L'accord est remarquable pour un dispositif artisanal !
\item Dosage de la solution inconnue. D'après la loi de Biot :
\[ c_{\text{inconnue}} = \frac{\alpha}{[\alpha] \times \ell} = \frac{7,9}{52,5 \times 1,00} \approx \SI{0,150}{g.mL^{-1}} = \SI{150}{g.L^{-1}} \]
La solution inconnue a une concentration voisine de \SI{150}{g.L^{-1}}.
\end{enumerate}
\end{exoresolu}

\begin{exoresolu}[Chiralité du butan-2-ol et projection de Fischer]
\textbf{Énoncé.} (1) Pourquoi le carbone n°2 du butan-2-ol est-il dit asymétrique ? (2) Pourquoi les deux modèles construits ne sont-ils pas superposables ? Comment appelle-t-on ces deux molécules ? (3) Le butan-1-ol \chemfig{CH_3-CH_2-CH_2-CH_2-OH} possède-t-il des énantiomères ? Justifier. (4) Donne la configuration absolue $(R)$ ou $(S)$ pour les deux projections de Fischer tracées (règle de Cahn-Ingold-Prelog : priorité OH > CH2CH3 > CH3 > H).
\tcblower
\textbf{Solution.}
\begin{enumerate}
\item Le carbone n°2 est lié à \textbf{quatre groupes tous différents} : \ce{H}, \ce{OH}, \ce{CH3} (méthyle) et \ce{CH2-CH3} (éthyle). C'est la définition d'un \cle{carbone asymétrique} (ou stéréocentre), noté C*.
\item Les deux modèles sont images l'un de l'autre dans un miroir plan. Aucune rotation dans l'espace (opération propre) ne permet de faire coïncider simultanément les quatre groupes : ils sont \textbf{non superposables}. Ce sont deux \cle{énantiomères}, et la molécule est dite \cle{chirale} (du grec \emph{kheir}, main), comme nos deux mains.
\item Dans le butan-1-ol, aucun carbone ne porte quatre groupes différents : le carbone 1 porte deux atomes d'hydrogène (\ce{CH2OH}), les carbones 2 et 3 sont des \ce{CH2}. La molécule possède un plan de symétrie : elle est \textbf{achirale}, superposable à son image miroir, et ne possède pas d'énantiomère. Elle est donc optiquement inactive.
\item Règle CIP : priorités : 1=\ce{OH}, 2=\ce{CH2CH3}, 3=\ce{CH3}, 4=\ce{H}.
\begin{itemize}
\item Projection 1 (H à gauche, OH à droite, CH3 en haut, CH2CH3 en bas) : le groupe de plus basse priorité (H) est sur liaison horizontale (sortant). Le sens 1$\to$2$\to$3 est horaire. Comme H sort, la configuration réelle est l'inverse : \textbf{$(S)$}.
\item Projection 2 (image miroir : H à droite, OH à gauche) : le sens 1$\to$2$\to$3 est trigonométrique. H sort $\rightarrow$ configuration \textbf{$(R)$}.
\end{itemize}
Les deux énantiomères sont donc le \textbf{$(R)$-butan-2-ol} et le \textbf{$(S)$-butan-2-ol}.
\end{enumerate}
\end{exoresolu}

\begin{exoresolu}[Isomérie Z/E : but-2-ène]
\textbf{Énoncé.} Dessine les représentations de Cram des deux stéréoisomères du but-2-ène (\ce{CH3-CH=CH-CH3}). Nomme-les selon la nomenclature Z/E. Sont-ils des énantiomères ou des diastéréoisomères ? Sont-ils chiraux ?
\tcblower
\textbf{Solution.}
Le but-2-ène présente une double liaison \ce{C=C} avec sur chaque carbone deux groupes différents (\ce{H} et \ce{CH3}). Il existe deux diastéréoisomères :
\begin{itemize}
\item Isomère \textbf{Z} (\emph{zusammen}, ensemble) : les deux groupes de priorité forte (\ce{CH3}) du même côté de la double liaison. Représentation de Cram : \chemfig{CH_3-/[0]=_/[6]CH_3} (schématique). Molécule \textbf{achirale} (plan de symétrie moléculaire).
\item Isomère \textbf{E} (\emph{entgegen}, opposé) : les deux groupes \ce{CH3} de part et d'autre. Représentation de Cram : \chemfig{CH_3-/[0]=_[6]CH_3}. Molécule \textbf{achirale} (centre de symétrie).
\end{itemize}
Ces deux isomères sont des \textbf{diastéréoisomères} (stéréoisomères non images l'un de l'autre dans un miroir). Ils ont des propriétés physiques différentes (PF : Z \SI{-139}{\celsius}, E \SI{-106}{\celsius}).
\end{exoresolu}

\courssub{Conclusion}

Ce TP t'a permis de \emph{concrétiser} les notions abstraites du cours. D'une part, la \textbf{chiralité moléculaire} se manifeste par une propriété macroscopique mesurable : une solution de glucose fait tourner le plan de polarisation de la lumière, et l'angle de rotation obéit à la \textbf{loi de Biot} $\alpha = [\alpha]\,\ell\,c$, ce qui permet de doser des substances actives (c'est le principe du \cle{saccharimètre} utilisé dans les sucreries comme la SOSUCAM). D'autre part, les modèles moléculaires montrent que :
\begin{itemize}
\item Deux énantiomères, comme le $(R)$- et le $(S)$-butan-2-ol, sont des molécules \textbf{distinctes et non superposables}, images l'une de l'autre dans un miroir ; leurs projections de Fischer s'en déduisent directement.
\item L'isomérie Z/E (but-2-ène) donne des diastéréoisomères achiraux aux propriétés physiques distinctes.
\item Les molécules flexibles adoptent différentes \textbf{conformations} par rotation autour des liaisons simples : l'éthane alterne entre forme éclipsée (instable) et décalée (stable) ; le cyclohexane privilégie la conformation \textbf{chaise} (sans tension) à la conformation \textbf{bateau} (interactions de drapeau).
\end{itemize}
Tu comprends désormais pourquoi la stéréochimie est cruciale en pharmacie : deux énantiomères peuvent avoir des effets biologiques radicalement différents, comme le rappelle le drame historique de la \cle{thalidomide} (un énantiomère sédatif, l'autre tératogène) ou l'\cle{ibuprofène} (seul l'énantiomère $(S)$ est actif anti-inflammatoire, l'autre est inoffensif mais inactive). En agroalimentaire, l'arôme de la \cle{carvone} change selon la configuration : $(R)$ = menthe, $(S)$ = carvi. La stéréochimie n'est pas une abstraction, c'est la clé de la reconnaissance moléculaire dans le vivant.

\newpage

\coursec{Méthodes et exercices résolus}

\courssub{Fiche méthode 1 : Différencier deux isomères de constitution}

\begin{methode}[Reconnaître le type d'isomérie de constitution]
Deux composés sont des \cle{isomères de constitution} s'ils ont la \textbf{même formule brute} mais des \textbf{enchaînements d'atomes différents}. Pour identifier le type :
\begin{enumerate}
\item Écrire la \textbf{formule brute} de chaque composé : si elles diffèrent, ce ne sont pas des isomères.
\item Écrire la \textbf{formule semi-développée} de chacun et repérer les \textbf{groupes caractéristiques} (fonctions).
\item Comparer :
\begin{itemize}
\item \textbf{Fonctions différentes} $\Rightarrow$ isomérie de \cle{fonction} (ex. alcool \ce{$R$-OH} / éther \ce{$R$-O-$R'$}).
\item \textbf{Même fonction, position différente} sur la chaîne $\Rightarrow$ isomérie de \cle{position} (ex. butan-1-ol / butan-2-ol).
\item \textbf{Même fonction, même position, squelette carboné différent} (ramifications) $\Rightarrow$ isomérie de \cle{chaîne} (ex. butane / 2-méthylpropane).
\end{itemize}
\item Conclure en une phrase : « A et B sont des isomères de \ldots car \ldots ».
\end{enumerate}
\textbf{Réflexe Bac} : toujours vérifier d'abord la formule brute, puis la fonction, puis la position, puis la chaîne — dans cet ordre.
\end{methode}

\begin{exoresolu}[Isomères de formule brute \ce{C4H10O}]
On considère les quatre composés suivants :
A : \ce{CH3-CH2-CH2-CH2-OH} ; B : \ce{CH3-CH2-CH(OH)-CH3} ; C : \ce{(CH3)2CH-CH2-OH} ; D : \ce{CH3-CH2-O-CH2-CH3}.
\begin{enumerate}
\item Montrer que ces composés sont isomères.
\item Préciser, pour chaque couple (A,B), (A,C) et (A,D), le type d'isomérie.
\end{enumerate}
\tcblower
\textbf{1. Vérification de la formule brute.}

Comptons les atomes de chaque composé :
\begin{center}
\begin{tabularx}{\linewidth}{|l|X|c|}
\hline
\rowcolor{popblueL} Composé & Formule semi-développée & Formule brute \\
\hline
A & \ce{CH3-CH2-CH2-CH2-OH} & \ce{C4H10O} \\
\hline
B & \ce{CH3-CH2-CH(OH)-CH3} & \ce{C4H10O} \\
\hline
C & \ce{(CH3)2CH-CH2-OH} & \ce{C4H10O} \\
\hline
D & \ce{CH3-CH2-O-CH2-CH3} & \ce{C4H10O} \\
\hline
\end{tabularx}
\end{center}
Les quatre composés ont la \textbf{même formule brute} \ce{C4H10O} mais des enchaînements d'atomes différents : ce sont bien des \textbf{isomères}.

\textbf{2. Nature des isoméries.}

\textbf{Couple (A, B)} : A est le butan-1-ol, B est le butan-2-ol. Même fonction alcool, même chaîne à 4 carbones, mais le groupe \ce{-OH} est porté par le carbone 1 dans A et par le carbone 2 dans B : \textbf{isomérie de position}.

\textbf{Couple (A, C)} : A est le butan-1-ol (chaîne linéaire), C est le 2-méthylpropan-1-ol (chaîne ramifiée). Même fonction alcool en position 1, mais squelettes carbonés différents : \textbf{isomérie de chaîne}.

\textbf{Couple (A, D)} : A est un alcool (groupe \ce{-OH}), D est l'éthoxyéthane, un éther-oxyde (groupe \ce{-O-}). Fonctions chimiques différentes : \textbf{isomérie de fonction}.

\textbf{Conclusion} : A et B sont isomères de position ; A et C isomères de chaîne ; A et D isomères de fonction.
\end{exoresolu}

\courssub{Fiche méthode 2 : Représenter une molécule en perspective (Cram)}

\begin{methode}[Construire une représentation de Cram]
La représentation de \cle{Cram} traduit la géométrie tétraédrique d'un carbone $sp^3$ :
\begin{enumerate}
\item Repérer le \textbf{carbone central} et ses 4 substituants.
\item Placer \textbf{deux liaisons dans le plan} de la feuille (traits pleins fins), avec un angle d'environ $109^\circ$ représenté de façon aplatie.
\item Placer \textbf{un substituant devant le plan} : liaison en \textbf{trait plein épais} (coin).
\item Placer \textbf{un substituant derrière le plan} : liaison en \textbf{hachures ou pointillés}.
\item Vérifier : 4 liaisons au total, jamais 3 liaisons dans le plan (cela déformerait la géométrie).
\end{enumerate}
\textbf{Réflexe} : pour passer d'un énantiomère à l'autre, il suffit d'\textbf{échanger deux substituants} (par exemple celui de devant avec celui de derrière).
\end{methode}

\begin{exoresolu}[Représentation de Cram de l'acide lactique]
L'acide lactique, présent dans le lait caillé et les muscles après un effort, a pour formule semi-développée \ce{CH3-CH(OH)-COOH}.
\begin{enumerate}
\item Repérer le carbone asymétrique.
\item Donner une représentation de Cram de cette molécule.
\end{enumerate}
\tcblower
\textbf{1. Carbone asymétrique.}

Un \cle{carbone asymétrique} est un carbone tétraédrique porteur de \textbf{quatre substituants différents}. Ici, le carbone central porte : \ce{-H}, \ce{-OH}, \ce{-CH3} et \ce{-COOH} : quatre groupes différents. C'est bien un carbone asymétrique, noté C*.

\textbf{2. Représentation de Cram.}

On place \ce{CH3} et \ce{COOH} dans le plan, \ce{OH} devant (trait plein) et \ce{H} derrière (hachures) :
\begin{center}
\chemfig{C(-[2]H)(<[5]CH_3)(<:[7]OH)-COOH}
\end{center}
Toute autre disposition respectant la règle (2 liaisons dans le plan, 1 devant, 1 derrière) est acceptée. En échangeant \ce{OH} et \ce{H}, on obtient l'énantiomère image.

\textbf{Conclusion} : l'acide lactique possède un carbone asymétrique (le C2) ; il existe donc sous forme de deux énantiomères représentables en perspective de Cram.
\end{exoresolu}

\courssub{Fiche méthode 3 : Représenter les énantiomères (Cram et Fischer)}

\begin{methode}[Représenter deux énantiomères]
\textbf{En perspective de Cram :}
\begin{enumerate}
\item Représenter le premier énantiomère selon la fiche 2.
\item Obtenir le second en \textbf{échangeant deux substituants} (ou en dessinant l'image dans un miroir plan).
\end{enumerate}
\textbf{En projection de Fischer :}
\begin{enumerate}
\item Placer le carbone asymétrique au \textbf{centre d'une croix}.
\item Les liaisons \textbf{horizontales} sortent du plan (vers l'observateur) ; les liaisons \textbf{verticales} rentrent dans le plan.
\item Placer la chaîne principale sur la \textbf{verticale}, le groupe le plus oxydé (ex. \ce{COOH}) \textbf{en haut}.
\item L'énantiomère s'obtient en \textbf{permutant les deux groupes horizontaux} (ou les deux verticaux).
\end{enumerate}
\textbf{Attention} : une rotation de $90^\circ$ d'une projection de Fischer \textbf{change} la configuration ; seule une rotation de $180^\circ$ dans le plan la conserve.
\end{methode}

\begin{exoresolu}[Énantiomères de l'acide lactique en Fischer]
Représenter en projection de Fischer les deux énantiomères de l'acide lactique \ce{CH3-CH(OH)-COOH}.
\tcblower
La chaîne carbonée \ce{COOH-CH3} est placée sur la verticale, avec le groupe le plus oxydé \ce{COOH} en haut. Les groupes \ce{OH} et \ce{H} occupent les positions horizontales (vers l'avant).

\textbf{Énantiomère 1} (\ce{OH} à droite) :
\begin{center}
\chemfig{C(-[2]COOH)(-[6]CH_3)(<:[:150]H)(<[:30]OH)}
\end{center}
ou, en croix de Fischer conventionnelle :
\begin{center}
\begin{tabular}{c}
\ce{COOH} \\
$|$ \\
\ce{H} --- C* --- \ce{OH} \\
$|$ \\
\ce{CH3}
\end{tabular}
\end{center}

\textbf{Énantiomère 2} : on permute \ce{H} et \ce{OH} :
\begin{center}
\begin{tabular}{c}
\ce{COOH} \\
$|$ \\
\ce{HO} --- C* --- \ce{H} \\
$|$ \\
\ce{CH3}
\end{tabular}
\end{center}

Les deux représentations sont \textbf{images l'une de l'autre dans un miroir} et \textbf{non superposables} : ce sont bien deux énantiomères. Un mélange équimolaire des deux constitue un \cle{mélange racémique}, optiquement inactif.

\textbf{Conclusion} : en Fischer, chaîne verticale (groupe oxydé en haut), substituants horizontaux vers l'avant ; la permutation de deux groupes donne l'énantiomère.
\end{exoresolu}

\courssub{Fiche méthode 4 : Discerner énantiomères et diastéréoisomères}

\begin{methode}[Classer une paire de stéréo-isomères]
\begin{enumerate}
\item Vérifier que les deux composés ont la \textbf{même formule brute} et la \textbf{même constitution} (même enchaînement d'atomes) : sinon, ce sont des isomères de constitution.
\item Repérer les \textbf{éléments de stéréoisomérie} : carbones asymétriques C* et doubles liaisons \ce{C=C} porteuses de groupes différents.
\item Comparer les configurations :
\begin{itemize}
\item Configurations \textbf{inversées sur tous} les éléments stéréogènes, molécules images dans un miroir \textbf{non superposables} $\Rightarrow$ \cle{énantiomères}.
\item Configurations inversées sur \textbf{certains seulement} des éléments (ou images superposables) $\Rightarrow$ \cle{diastéréoisomères}.
\end{itemize}
\item Cas particulier : pour une double liaison \ce{C=C}, les isomères \cle{Z} (groupes prioritaires du même côté) et \cle{E} (de part et d'autre) sont des diastéréoisomères.
\end{enumerate}
\textbf{Réflexe} : énantiomères = propriétés physiques identiques (sauf pouvoir rotatoire, opposé) ; diastéréoisomères = propriétés physiques différentes (températures de fusion, ébullition, solubilité).
\end{methode}

\begin{exoresolu}[Acide 2,3-dihydroxybutanedioïque (acide tartrique)]
L'acide tartrique, présent dans le vin et le jus de palme fermenté, a pour formule \ce{HOOC-CH(OH)-CH(OH)-COOH}. On considère les trois stéréo-isomères suivants, décrits par la position du groupe \ce{OH} sur chaque carbone asymétrique (C2 et C3) en projection de Fischer :
\begin{itemize}
\item Isomère I : \ce{OH} à droite sur C2, \ce{OH} à droite sur C3 ;
\item Isomère II : \ce{OH} à gauche sur C2, \ce{OH} à gauche sur C3 ;
\item Isomère III : \ce{OH} à droite sur C2, \ce{OH} à gauche sur C3.
\end{itemize}
\begin{enumerate}
\item Justifier l'existence de deux carbones asymétriques.
\item Préciser la relation de stéréo-isomérie entre I et II, puis entre I et III.
\end{enumerate}
\tcblower
\textbf{1. Carbones asymétriques.}

Le carbone C2 porte : \ce{-H}, \ce{-OH}, \ce{-COOH} et \ce{-CH(OH)COOH} : quatre groupes différents. Il en va de même pour C3 par symétrie de la formule. La molécule possède donc \textbf{deux carbones asymétriques} C2* et C3*.

\textbf{2. Relations de stéréo-isomérie.}

\textbf{Couple (I, II)} : la configuration est inversée sur C2 (droite $\to$ gauche) \textbf{et} sur C3 (droite $\to$ gauche). Les deux molécules sont images l'une de l'autre dans un miroir et non superposables : I et II sont des \textbf{énantiomères}.

\textbf{Couple (I, III)} : la configuration est inversée sur C3 seulement, mais \textbf{pas} sur C2. I et III ne sont pas images dans un miroir : ce sont des \textbf{diastéréoisomères}.

Remarque : l'isomère III possède un plan de symétrie interne (entre C2 et C3) ; il est \textbf{achiral} malgré ses deux C* : c'est le composé \emph{méso}, optiquement inactif. Il n'existe donc que \textbf{trois} stéréo-isomères et non quatre ($2^2$), car deux des quatre combinaisons sont identiques.

\textbf{Conclusion} : I et II sont énantiomères (toutes configurations inversées) ; I et III sont diastéréoisomères (une seule configuration inversée) ; III est un composé \emph{méso} achiral.
\end{exoresolu}

\courssub{Fiche méthode 5 : Identifier et nommer les isomères Z/E}

\begin{methode}[Reconnaître et nommer un isomère Z ou E]
\begin{enumerate}
\item Repérer la double liaison \ce{C=C} et vérifier que \textbf{chaque} carbone de la double liaison porte \textbf{deux substituants différents} : sinon, pas d'isomérie Z/E.
\item Sur chaque carbone, déterminer le \textbf{groupe prioritaire} (numéro atomique $Z$ le plus élevé ; ex. \ce{Br} > \ce{Cl} > \ce{C} > \ce{H}).
\item Observer la position relative des deux groupes prioritaires :
\begin{itemize}
\item \textbf{Du même côté} de la double liaison $\Rightarrow$ isomère \cle{Z} (\emph{zusammen}, ensemble) ;
\item \textbf{De part et d'autre} $\Rightarrow$ isomère \cle{E} (\emph{entgegen}, opposé).
\end{itemize}
\item Nommer en faisant précéder le nom de (Z)- ou (E)-.
\end{enumerate}
\textbf{Réflexe} : Z et E sont des \textbf{diastéréoisomères} : leurs propriétés physiques diffèrent (ex. températures d'ébullition).
\end{methode}

\begin{exoresolu}[Isomérie Z/E du but-2-ène]
Le but-2-ène \ce{CH3-CH=CH-CH3} est un gaz obtenu lors du craquage des hydrocarbures (procédé utilisé à la SONARA).
\begin{enumerate}
\item Montrer que le but-2-ène présente une isomérie Z/E.
\item Représenter et nommer les deux stéréo-isomères.
\end{enumerate}
\tcblower
\textbf{1. Existence de l'isomérie Z/E.}

Chaque carbone de la double liaison \ce{C=C} porte deux substituants différents : \ce{-H} et \ce{-CH3}. La rotation autour de la double liaison étant impossible, les deux dispositions spatiales sont distinctes et non interconvertibles : il existe bien une \textbf{isomérie Z/E}.

\textbf{2. Représentation et nommage.}

Sur chaque carbone, le groupe prioritaire est \ce{CH3} (carbone, $Z=6$) devant \ce{H} ($Z=1$).

\textbf{Isomère (Z)} : les deux groupes \ce{CH3} du même côté :
\begin{center}
\chemfig{CH_3-[:30]=[:-30](-[:-90]H)-[:30]CH_3}
\end{center}
\begin{center}
\chemfig{CH_2=C(-CH_3)-CH_3}
\end{center}
Les deux \ce{CH3} prioritaires sont du même côté : c'est le \textbf{(Z)-but-2-ène}.

\textbf{Isomère (E)} : les deux groupes \ce{CH3} de part et d'autre :
\begin{center}
\begin{tabular}{ccc}
\ce{CH3} & & \ce{H} \\
 & \ce{C=C} & \\
\ce{H} & & \ce{CH3}
\end{tabular}
\end{center}
Les deux \ce{CH3} prioritaires sont de part et d'autre : c'est le \textbf{(E)-but-2-ène}.

\textbf{Conclusion} : le but-2-ène existe sous deux diastéréoisomères, (Z)-but-2-ène et (E)-but-2-ène, de propriétés physiques différentes ($\theta_{\text{éb}} = \SI{3,7}{\celsius}$ pour le Z et $\SI{0,9}{\celsius}$ pour le E).
\end{exoresolu}

\courssub{Fiche méthode 6 : Représenter les conformations (Newman, chaise, bateau)}

\begin{methode}[Construire une représentation de Newman]
\begin{enumerate}
\item Choisir la \textbf{liaison \ce{C-C} d'observation} : l'œil regarde le long de cette liaison.
\item Représenter le \textbf{carbone avant} par un \textbf{point} d'où partent 3 liaisons à $120^\circ$.
\item Représenter le \textbf{carbone arrière} par un \textbf{cercle} d'où partent 3 liaisons à $120^\circ$.
\item Placer les substituants :
\begin{itemize}
\item \textbf{Conformation éclipsée} : les liaisons avant et arrière sont \textbf{alignées} (décalage angulaire $0^\circ$) — conformation la \textbf{moins stable} (gêne stérique maximale) ;
\item \textbf{Conformation décalée} : les liaisons arrière sont placées \textbf{entre} celles de l'avant (décalage $60^\circ$) — conformation la \textbf{plus stable}.
\end{itemize}
\item Pour passer de l'une à l'autre : \textbf{rotation de $60^\circ$} du carbone arrière autour de la liaison \ce{C-C}.
\end{enumerate}
\textbf{Pour le cyclohexane} : la forme \cle{chaise} (la plus stable, angles proches de $109^\circ$, toutes les liaisons décalées) s'obtient en relevant deux sommets opposés d'un hexagone, l'un vers le haut, l'autre vers le bas ; la forme \cle{bateau} (moins stable) relève les deux sommets du même côté.
\end{methode}

\begin{exoresolu}[Conformations de l'éthane et du cyclohexane]
\begin{enumerate}
\item Représenter en projection de Newman les conformations éclipsée et décalée de l'éthane \ce{CH3-CH3}, en regardant le long de la liaison \ce{C-C}.
\item Indiquer la conformation la plus stable et justifier.
\item Représenter les conformations chaise et bateau du cyclohexane \ce{C6H12} et indiquer la plus stable.
\end{enumerate}
\tcblower
\textbf{1. Projections de Newman de l'éthane.}

L'œil est placé dans l'axe de la liaison \ce{C-C}. Le carbone avant est figuré par le point central, le carbone arrière par le cercle ; chacun porte trois atomes d'hydrogène.

\textbf{Conformation éclipsée} (hydrogènes alignés) :
\begin{center}
\begin{tikzpicture}[scale=0.9]
\draw[popline, thick] (0,0) circle (0.9);
\foreach \a in {90,210,330}{
  \draw[popblue, thick] (0,0) -- (\a:0.75);
  \node[popdark] at (\a:1.05) {\footnotesize H};
}
\foreach \a in {90,210,330}{
  \draw[poporange, thick, dashed] (\a:0.75) -- (\a:0.9);
}
\node[popmuted] at (0,-1.5) {\footnotesize éclipsée (décalage $0^\circ$)};
\end{tikzpicture}
\hspace{1.5cm}
\begin{tikzpicture}[scale=0.9]
\draw[popline, thick] (0,0) circle (0.9);
\foreach \a in {90,210,330}{
  \draw[popblue, thick] (0,0) -- (\a:0.75);
  \node[popdark] at (\a:1.05) {\footnotesize H};
}
\foreach \a in {30,150,270}{
  \draw[poporange, thick] (\a:0.9) -- (\a:0.55);
  \node[popdark] at (\a:1.2) {\footnotesize H};
}
\node[popmuted] at (0,-1.5) {\footnotesize décalée (décalage $60^\circ$)};
\end{tikzpicture}
\end{center}

\textbf{2. Stabilité.}

La conformation \textbf{décalée} est la plus stable : les liaisons \ce{C-H} (et les doublets liants) des deux carbones sont le plus \textbf{éloignées} possible, ce qui minimise les répulsions électroniques et la gêne stérique. L'éclipsée correspond à un maximum d'énergie (environ $\SI{12}{kJ.mol^{-1}}$ au-dessus de la décalée). La rotation autour de la liaison \ce{C-C} reste quasi libre à température ordinaire.

\textbf{3. Cyclohexane.}

\begin{center}
\begin{tikzpicture}[scale=1.0]
% chaise
\coordinate (A) at (0,0.35);
\coordinate (B) at (0.7,0.0);
\coordinate (C) at (1.4,0.35);
\coordinate (D) at (1.4,1.05);
\coordinate (E) at (0.7,1.4);
\coordinate (F) at (0,1.05);
\draw[popblue, thick] (A)--(B)--(C)--(D)--(E)--(F)--cycle;
\node[popmuted] at (0.7,-0.5) {\footnotesize chaise (stable)};
% bateau
\begin{scope}[xshift=4cm]
\coordinate (A2) at (0,0.35);
\coordinate (B2) at (0.7,0.0);
\coordinate (C2) at (1.4,0.35);
\coordinate (D2) at (1.4,1.05);
\coordinate (E2) at (0.7,1.4);
\coordinate (F2) at (0,1.05);
\draw[poporange, thick] (A2)--(B2)--(C2)--(D2)--(E2)--(F2)--cycle;
\draw[poporange, thick] (0,1.05) -- (0,1.5) -- (1.4,1.5) -- (1.4,1.05);
\node[popmuted] at (0.7,-0.5) {\footnotesize bateau (moins stable)};
\end{scope}
\end{tikzpicture}
\end{center}

Dans la forme \textbf{chaise}, tous les angles de valence sont voisins de $109^\circ$ et toutes les liaisons \ce{C-H} sont décalées : c'est la conformation \textbf{la plus stable}. Dans la forme \textbf{bateau}, des hydrogènes se font face (gêne stérique, interactions « drapeau ») : elle est \textbf{moins stable}.

\textbf{Conclusion} : l'éthane privilégie la conformation décalée ; le cyclohexane adopte préférentiellement la conformation chaise, ce qui explique la grande stabilité des cyclohexanes naturels.
\end{exoresolu}

\begin{attention}[Les 5 erreurs qui coûtent des points au Bac]
\begin{enumerate}
\item \textbf{Confondre isomérie de position et de chaîne.} Avant de conclure, vérifie d'abord la fonction, puis la position du groupe, puis seulement le squelette. Butan-1-ol / butan-2-ol : position ; butan-1-ol / 2-méthylpropan-1-ol : chaîne.
\item \textbf{Déclarer asymétrique un carbone portant deux groupes identiques.} Un C* doit porter \textbf{quatre substituants différents}. Dans \ce{CH3-CH2-CH(OH)-CH3}, le C2 porte deux groupes méthyle ? Non : \ce{-CH3} et \ce{-CH2CH3} sont différents — mais dans \ce{CH3-CH(OH)-CH3}, le C2 porte deux \ce{-CH3} identiques : il n'est \textbf{pas} asymétrique. Vérifie chaque substituant en entier, pas seulement le premier atome.
\item \textbf{Tracer 3 liaisons dans le plan en représentation de Cram.} La règle est stricte : \textbf{2 liaisons dans le plan, 1 devant (plein), 1 derrière (hachuré)}. Trois traits dans le plan rendent la représentation fausse.
\item \textbf{Faire tourner une projection de Fischer de $90^\circ$.} Une rotation de $90^\circ$ transforme la molécule en son \textbf{énantiomère} ; seule une rotation de $180^\circ$ dans le plan conserve la configuration. Ne « retourne » jamais une Fischer pour la comparer à une autre.
\item \textbf{Inverser Z et E ou oublier la condition d'existence.} Z = groupes prioritaires du \textbf{même côté} ; E = de part et d'autre. Et surtout : si un des deux carbones de la double liaison porte \textbf{deux substituants identiques} (ex. \ce{CH2=CH-CH3}), il n'y a \textbf{pas} d'isomérie Z/E. Le dire dans la copie montre ta rigueur.
\end{enumerate}
\end{attention}

\newpage

\coursec{Exercices}

\courssub{Je vérifie mes connaissances}

\begin{exercice}[QCM : les bons réflexes]
Pour chaque question, choisir la (ou les) bonne(s) réponse(s).
\begin{enumerate}
\item Deux isomères de constitution ont :
\begin{enumerate}
\item la même formule brute et la même formule semi-développée ;
\item la même formule brute mais des formules semi-développées différentes ;
\item des formules brutes différentes.
\end{enumerate}
\item Un carbone asymétrique est un carbone :
\begin{enumerate}
\item tétragonal lié à quatre atomes ou groupes d'atomes différents ;
\item porteur d'une double liaison ;
\item lié à au moins un atome d'hydrogène.
\end{enumerate}
\item Deux énantiomères sont :
\begin{enumerate}
\item deux stéréo-isomères images l'un de l'autre dans un miroir et non superposables ;
\item deux stéréo-isomères images l'un de l'autre dans un miroir et superposables ;
\item deux isomères de constitution.
\end{enumerate}
\item Un mélange racémique :
\begin{enumerate}
\item dévie le plan de la lumière polarisée vers la droite ;
\item contient les deux énantiomères en quantités égales et est optiquement inactif ;
\item ne contient qu'un seul énantiomère.
\end{enumerate}
\item Le nombre maximal de stéréo-isomères d'une molécule possédant $n$ carbones asymétriques différents est :
\begin{enumerate}
\item $n$ ; \item $2n$ ; \item $2^n$.
\end{enumerate}
\end{enumerate}
\end{exercice}

\begin{exercice}[Vrai ou faux, en justifiant]
Répondre par \textbf{vrai} ou \textbf{faux}, puis justifier chaque réponse en une ou deux phrases.
\begin{enumerate}
\item Le butan-1-ol et le butan-2-ol sont des isomères de fonction.
\item Toute molécule possédant un carbone asymétrique est chirale.
\item Deux diastéréoisomères sont toujours images l'un de l'autre dans un miroir.
\item Le pent-2-ène présente l'isomérie $Z/E$, mais pas le but-2-ène.
\item Une solution d'un énantiomère pur est dite dextrogyre si elle fait tourner le plan de polarisation de la lumière vers la droite.
\item Les conformations éclipsée et décalée de l'éthane sont deux molécules différentes que l'on peut isoler à température ordinaire.
\end{enumerate}
\end{exercice}

\begin{exercice}[Texte à trous : le polarimètre]
Recopier et compléter le texte suivant avec les mots ou expressions de la liste : \emph{dextrogyre ; pouvoir rotatoire spécifique ; polarimètre ; concentration massique ; lévogyre ; longueur ; Biot ; racémique}.

\medskip
Une substance est dite \textbf{optiquement active} lorsqu'elle fait tourner le plan de polarisation de la lumière polarisée. On mesure l'angle de rotation $\alpha$ à l'aide d'un \ldots\ldots\ldots\ldots . Si le plan tourne vers la droite (sens des aiguilles d'une montre), la substance est dite \ldots\ldots\ldots\ldots ; si le plan tourne vers la gauche, elle est dite \ldots\ldots\ldots\ldots . La loi de \ldots\ldots\ldots\ldots s'écrit $\alpha = [\alpha] \times \ell \times C$, où $\ell$ désigne la \ldots\ldots\ldots\ldots de la cuve, $C$ la \ldots\ldots\ldots\ldots de la solution et $[\alpha]$ le \ldots\ldots\ldots\ldots de la substance. Un mélange \ldots\ldots\ldots\ldots , contenant les deux énantiomères en quantités égales, donne un angle $\alpha$ nul.
\end{exercice}

\begin{exercice}[Définitions express]
Définir en une ou deux phrases chacun des termes suivants, en donnant à chaque fois un exemple simple :
\begin{enumerate}
\item isomères de chaîne ;
\item carbone asymétrique ;
\item énantiomères ;
\item diastéréoisomères ;
\item représentation de Cram ;
\item conformation décalée.
\end{enumerate}
\end{exercice}

\courssub{Je m'entraîne}

\begin{exercice}[Isomères de constitution du \ce{C4H10O}]
On considère les composés organiques de formule brute \ce{C4H10O} ne possédant que des liaisons simples.
\begin{enumerate}
\item Écrire les formules semi-développées de tous les alcools de formule brute \ce{C4H10O} et les nommer.
\item Écrire les formules semi-développées de tous les éthers de formule brute \ce{C4H10O} et les nommer.
\item Pour chaque paire d'isomères, préciser le type d'isomérie mis en jeu (chaîne, position ou fonction).
\item Parmi tous ces isomères, un seul possède un carbone asymétrique. Lequel ? Repérer ce carbone par un astérisque et représenter ses deux énantiomères en perspective de Cram.
\end{enumerate}
\emph{Données :} $M(\ce{C}) = \SI{12}{g.mol^{-1}}$ ; $M(\ce{H}) = \SI{1}{g.mol^{-1}}$ ; $M(\ce{O}) = \SI{16}{g.mol^{-1}}$.
\end{exercice}

\begin{exercice}[Isomérie $Z/E$ : le bon candidat]
On considère le pent-2-ène et le 2-méthylbut-2-ène.
\begin{enumerate}
\item Écrire la formule semi-développée de chacun de ces deux alcènes.
\item Rappeler la condition qu'une double liaison \ce{C=C} doit remplir pour donner lieu à l'isomérie $Z/E$.
\item Montrer qu'un seul des deux alcènes présente l'isomérie $Z/E$. Lequel ?
\item Représenter les deux stéréo-isomères de cet alcène et préciser, pour chacun, s'il s'agit de l'isomère $Z$ ou $E$.
\item Ces deux stéréo-isomères sont-ils des énantiomères ou des diastéréoisomères ? Justifier.
\end{enumerate}
\end{exercice}

\begin{exercice}[Polarimétrie sur le saccharose]
Une solution de saccharose (sucre de canne, très utilisé au Cameroun) de concentration massique $C = \SI{0,20}{g.mL^{-1}}$ est placée dans une cuve de polarimètre de longueur $\ell = \SI{1,0}{dm}$. À la température de \SI{20}{\celsius}, on mesure un angle de rotation $\alpha = +13,3\degres$.
\begin{enumerate}
\item Le saccharose est-il dextrogyre ou lévogyre ? Justifier.
\item Calculer le pouvoir rotatoire spécifique $[\alpha]$ du saccharose à \SI{20}{\celsius}. Préciser son unité.
\item Une solution de saccharose de concentration inconnue, placée dans la même cuve, donne un angle $\alpha' = +8,0\degres$. Calculer sa concentration massique $C'$.
\item Quel angle mesurerait-on avec un mélange racémique d'un sucre dans les mêmes conditions ? Justifier.
\end{enumerate}
\end{exercice}

\begin{exercice}[Newman : le 1,2-dichloroéthane]
On étudie la molécule de 1,2-dichloroéthane \ce{ClCH2-CH2Cl} en regardant selon l'axe de la liaison \ce{C-C}.
\begin{enumerate}
\item Représenter en projection de Newman la conformation éclipsée et la conformation décalée (anti) de cette molécule.
\item Indiquer, en justifiant, laquelle de ces deux conformations est la plus stable.
\item Peut-on séparer ces deux conformations à température ordinaire ? Expliquer pourquoi.
\item Citer la conformation la plus stable du cyclohexane et la représenter.
\end{enumerate}
\end{exercice}

\courssub{Je me prépare au Bac}

\begin{exercice}[L'alanine, une molécule chirale de la vie]
L'acide 2-aminopropanoïque, appelé \textbf{alanine}, de formule semi-développée \chemfig{CH_3-CH(-[2]NH_2)-COOH}, est l'un des acides $\alpha$-aminés constitutifs des protéines. On le trouve notamment dans les protéines du niébé et de l'arachide, très consommés au Cameroun.

\medskip
\emph{Données :} $M(\ce{C}) = \SI{12}{g.mol^{-1}}$ ; $M(\ce{H}) = \SI{1}{g.mol^{-1}}$ ; $M(\ce{O}) = \SI{16}{g.mol^{-1}}$ ; $M(\ce{N}) = \SI{14}{g.mol^{-1}}$.

\begin{enumerate}
\item Calculer la masse molaire moléculaire de l'alanine.
\item Montrer que la molécule d'alanine possède un carbone asymétrique. Le repérer par un astérisque sur la formule semi-développée.
\item Que peut-on en conclure quant à la chiralité de la molécule ?
\item Représenter en perspective de Cram les deux énantiomères de l'alanine, en précisant les conventions utilisées (traits plein, hachuré, plein dans le plan).
\item Représenter ces deux mêmes énantiomères en projection de Fischer, en plaçant le groupe \ce{COOH} en haut.
\item On dissout \SI{4,45}{g} d'alanine dans de l'eau distillée de façon à obtenir \SI{500}{mL} de solution. Calculer la concentration molaire de la solution.
\item La solution précédente contient en réalité les deux énantiomères en quantités strictement égales. Un élève affirme : « cette solution fera tourner le plan de la lumière polarisée ». A-t-il raison ? Justifier en nommant le type de mélange obtenu.
\item Dans les organismes vivants, seul l'un des deux énantiomères de l'alanine (l'énantiomère L) est utilisé pour fabriquer les protéines. Proposer une explication à l'importance de cette sélection naturelle.
\end{enumerate}
\end{exercice}

\begin{exercice}[Une molécule à deux carbones asymétriques]
On considère le 3-bromobutan-2-ol, de formule semi-développée \chemfig{CH_3-CH(-[2]OH)-CH(-[6]Br)-CH_3}. Cette molécule servira de modèle pour étudier les relations entre stéréo-isomères.

\medskip
\emph{Données :} $M(\ce{C}) = \SI{12}{g.mol^{-1}}$ ; $M(\ce{H}) = \SI{1}{g.mol^{-1}}$ ; $M(\ce{O}) = \SI{16}{g.mol^{-1}}$ ; $M(\ce{Br}) = \SI{80}{g.mol^{-1}}$.

\begin{enumerate}
\item Repérer, par un astérisque, les carbones asymétriques de la molécule. Justifier qu'ils sont bien asymétriques et qu'ils sont différents l'un de l'autre.
\item En déduire le nombre maximal de stéréo-isomères de cette molécule. Justifier par la relation appropriée.
\item Représenter en perspective de Cram les quatre stéréo-isomères, en notant $(R\text{-type})$ ou simplement $(+)/(-)$ les configurations des deux carbones pour les distinguer (on notera les formes A, B, C et D).
\item Parmi ces quatre formes, citer un couple d'énantiomères et un couple de diastéréoisomères, en justifiant chaque choix.
\item Classer toutes les relations (énantiomérie ou diastéréoisomérie) entre les formes A, B, C et D deux à deux.
\item On réalise une solution contenant uniquement les formes A et B en quantités égales. Cette solution dévie-t-elle le plan de la lumière polarisée ? Même question pour une solution contenant A et C en quantités égales. Justifier.
\end{enumerate}
\end{exercice}

\begin{exercice}[Le thalidomide : quand la stéréochimie devient un enjeu de santé publique]
\textbf{Document.} Au début des années 1960, le thalidomide était prescrit aux femmes enceintes comme calmant et contre les nausées. Ce médicament, vendu sous forme de mélange racémique, possède un carbone asymétrique. Or, si l'un des énantiomères possède bien l'effet sédatif recherché, l'autre s'est révélé \emph{tératogène} : il provoque de graves malformations du fœtus. Des milliers d'enfants sont nés avec des malformations avant le retrait du médicament. Depuis, l'industrie pharmaceutique cherche à produire des médicaments sous forme d'un seul énantiomère. Ainsi, seul l'énantiomère $(S)$ de l'ibuprofène est actif contre la douleur, et les deux énantiomères de la carvone ont des odeurs différentes : l'un sent la menthe, l'autre le carvi.

\begin{enumerate}
\item Expliquer, à l'aide du document, ce qu'est un mélange racémique et pourquoi il était moins coûteux de vendre le thalidomide sous cette forme.
\item Définir les termes \emph{carbone asymétrique} et \emph{énantiomère}.
\item Expliquer pourquoi deux énantiomères, qui ont les mêmes propriétés physiques (température de fusion, solubilité\ldots) et chimiques, peuvent avoir des effets biologiques très différents dans un organisme vivant, lui-même constitué de molécules chirales.
\item Citer, à partir du document, deux molécules (autres que le thalidomide) dont les énantiomères ont des propriétés biologiques différentes, en précisant ces différences.
\item La séparation des énantiomères d'un mélange racémique (opération appelée \emph{dédoublement}) est délicate car les deux énantiomères ont des propriétés physiques presque identiques. Proposer une raison pour laquelle une simple distillation ne permet pas de les séparer.
\item Rédiger un court paragraphe (cinq à huit lignes) expliquant pourquoi la maîtrise de la stéréochimie est aujourd'hui un enjeu de santé publique, y compris pour les pharmacies et les centres de santé du Cameroun.
\end{enumerate}
\end{exercice}

\newpage

\coursec{Corrigés des exercices}

\begin{corrige}[Exercice 1 — QCM : les bons réflexes]
\begin{enumerate}
\item \textbf{Réponse b.} Deux isomères de constitution ont la \textbf{même formule brute} mais des \textbf{formules semi-développées différentes} : les atomes ne sont pas reliés dans le même ordre. S'ils avaient la même formule semi-développée, ce serait le même composé ; s'ils avaient des formules brutes différentes, ils ne seraient pas isomères.
\item \textbf{Réponse a.} Un carbone asymétrique est un carbone \textbf{tétragonal} (quatre liaisons simples) lié à \textbf{quatre atomes ou groupes d'atomes différents}. Un carbone porteur d'une double liaison est trigonal (plan) et ne peut pas être asymétrique ; la présence d'un hydrogène n'est ni nécessaire ni suffisante.
\item \textbf{Réponse a.} Deux énantiomères sont des stéréo-isomères \textbf{images l'un de l'autre dans un miroir et non superposables} (comme les deux mains). Si les images sont superposables, c'est la même molécule ; les énantiomères ont la même formule semi-développée, ce ne sont donc pas des isomères de constitution.
\item \textbf{Réponse b.} Un mélange racémique contient les deux énantiomères en \textbf{quantités égales} ; les pouvoirs rotatoires, égaux en valeur absolue et de signes opposés, se compensent : le mélange est \textbf{optiquement inactif} ($\alpha = 0$).
\item \textbf{Réponse c.} Le nombre maximal de stéréo-isomères est $2^n$ : chaque carbone asymétrique offre deux configurations possibles, d'où $2 \times 2 \times \cdots \times 2 = 2^n$ combinaisons.
\end{enumerate}
\end{corrige}

\begin{corrige}[Exercice 2 — Vrai ou faux, en justifiant]
\begin{enumerate}
\item \textbf{Faux.} Le butan-1-ol \ce{CH3-CH2-CH2-CH2OH} et le butan-2-ol \ce{CH3-CH2-CHOH-CH3} ont la même fonction alcool ; seule la \emph{position} du groupe hydroxyle change. Ce sont des isomères de \textbf{position}.
\item \textbf{Vrai.} Une molécule possédant \emph{un seul} carbone asymétrique n'a aucun plan de symétrie : elle n'est pas superposable à son image dans un miroir, elle est donc chirale. (Attention : avec \emph{plusieurs} carbones asymétriques, une molécule \emph{méso} possédant un plan de symétrie serait achirale.)
\item \textbf{Faux.} Deux diastéréoisomères sont des stéréo-isomères qui ne sont \textbf{pas} images l'un de l'autre dans un miroir. C'est justement ce qui les distingue des énantiomères.
\item \textbf{Faux.} C'est l'inverse ! Dans le but-2-ène \ce{CH3-CH=CH-CH3}, chaque carbone de la double liaison porte deux groupes différents (\ce{H} et \ce{CH3}) : il présente l'isomérie $Z/E$. Le pent-2-ène \ce{CH3-CH=CH-CH2-CH3} la présente aussi d'ailleurs ; l'affirmation est donc doublement fausse.
\item \textbf{Vrai.} Par définition, une substance est dextrogyre si elle fait tourner le plan de polarisation vers la droite (sens des aiguilles d'une montre, angle $\alpha > 0$) ; lévogyre si elle le fait tourner vers la gauche.
\item \textbf{Faux.} Les conformations éclipsée et décalée de l'éthane se convertissent l'une en l'autre par simple rotation autour de la liaison \ce{C-C}, dont la barrière énergétique est très faible ($\approx \SI{12}{kJ.mol^{-1}}$). À température ordinaire, cette rotation est libre : on ne peut pas isoler les conformères.
\end{enumerate}
\end{corrige}

\begin{corrige}[Exercice 3 — Texte à trous : le polarimètre]
Une substance est dite \textbf{optiquement active} lorsqu'elle fait tourner le plan de polarisation de la lumière polarisée. On mesure l'angle de rotation $\alpha$ à l'aide d'un \textbf{polarimètre}. Si le plan tourne vers la droite (sens des aiguilles d'une montre), la substance est dite \textbf{dextrogyre} ; si le plan tourne vers la gauche, elle est dite \textbf{lévogyre}. La loi de \textbf{Biot} s'écrit $\alpha = [\alpha] \times \ell \times C$, où $\ell$ désigne la \textbf{longueur} de la cuve (en dm), $C$ la \textbf{concentration massique} de la solution (en \si{g.mL^{-1}}) et $[\alpha]$ le \textbf{pouvoir rotatoire spécifique} de la substance. Un mélange \textbf{racémique}, contenant les deux énantiomères en quantités égales, donne un angle $\alpha$ nul.
\end{corrige}

\begin{corrige}[Exercice 4 — Définitions express]
\begin{enumerate}
\item \textbf{Isomères de chaîne :} isomères de constitution qui diffèrent par l'enchaînement des atomes de carbone (chaîne linéaire ou ramifiée). Exemple : le butane \ce{CH3-CH2-CH2-CH3} et le 2-méthylpropane \ce{CH3-CH(CH3)-CH3}, tous deux de formule brute \ce{C4H10}.
\item \textbf{Carbone asymétrique :} atome de carbone tétragonal lié à quatre atomes ou groupes d'atomes tous différents ; on le note C$^*$. Exemple : le carbone 2 du butan-2-ol \chemfig{CH_3-\chembelow{C}{*}(-[2]OH)-CH_2-CH_3} porte \ce{H}, \ce{OH}, \ce{CH3} et \ce{C2H5}.
\item \textbf{Énantiomères :} couple de stéréo-isomères images l'un de l'autre dans un miroir et non superposables. Exemple : les deux formes de l'acide lactique \ce{CH3-CHOH-COOH}.
\item \textbf{Diastéréoisomères :} stéréo-isomères qui ne sont pas images l'un de l'autre dans un miroir. Exemple : les isomères $Z$ et $E$ du but-2-ène.
\item \textbf{Représentation de Cram :} représentation conventionnelle en perspective d'une molécule : les liaisons dans le plan sont des traits pleins fins, celle qui sort du plan vers l'observateur est un trait plein épais (triangulaire), celle qui rentre derrière le plan est un trait hachuré. Exemple : la représentation d'un énantiomère de l'alanine.
\item \textbf{Conformation décalée :} conformation dans laquelle, vue en projection de Newman, les liaisons issues des deux carbones voisins forment entre elles un angle de \ang{60} ; les groupes sont le plus éloignés possible, ce qui minimise les répulsions. Exemple : la conformation décalée de l'éthane, plus stable que l'éclipsée.
\end{enumerate}
\end{corrige}

\begin{corrige}[Exercice 5 — Isomères de constitution du \ce{C4H10O}]
\textbf{1. Les alcools (4 isomères) :}
\begin{center}
\begin{tabularx}{\linewidth}{|l|X|}
\hline
\rowcolor{popblueL} \textbf{Formule semi-développée} & \textbf{Nom} \\
\hline
\ce{CH3-CH2-CH2-CH2OH} & butan-1-ol \\
\hline
\ce{CH3-CH2-CHOH-CH3} & butan-2-ol \\
\hline
\ce{CH3-CH(CH3)-CH2OH} & 2-méthylpropan-1-ol \\
\hline
\ce{CH3-C(CH3)(OH)-CH3} & 2-méthylpropan-2-ol \\
\hline
\end{tabularx}
\end{center}

\textbf{2. Les éthers (3 isomères) :}
\begin{center}
\begin{tabularx}{\linewidth}{|l|X|}
\hline
\rowcolor{popblueL} \textbf{Formule semi-développée} & \textbf{Nom} \\
\hline
\ce{CH3-O-CH2-CH2-CH3} & 1-méthoxypropane \\
\hline
\ce{CH3-O-CH(CH3)-CH3} & 2-méthoxypropane \\
\hline
\ce{CH3-CH2-O-CH2-CH3} & éthoxyéthane \\
\hline
\end{tabularx}
\end{center}

\textbf{3. Types d'isomérie :}
\begin{itemize}
\item butan-1-ol / butan-2-ol : isomérie de \textbf{position} (même chaîne, même fonction, position du \ce{OH} différente) ;
\item butan-1-ol / 2-méthylpropan-1-ol : isomérie de \textbf{chaîne} (chaîne linéaire contre chaîne ramifiée) ;
\item butan-1-ol / éthoxyéthane : isomérie de \textbf{fonction} (alcool contre éther) ;
\item méthoxypropane / 2-méthoxypropane : isomérie de \textbf{chaîne}.
\end{itemize}

\textbf{4. Carbone asymétrique :} seul le \textbf{butan-2-ol} possède un carbone asymétrique : le carbone 2, lié à \ce{H}, \ce{OH}, \ce{CH3} et \ce{CH2CH3} (quatre groupes différents) :
\[ \chemfig{CH_3-\chembelow{C}{*}(-[2]OH)-CH_2-CH_3} \]
Les deux énantiomères en perspective de Cram :
\begin{center}
\begin{tikzpicture}
\node at (0,0) {\chemfig{C(-[2]OH)(-[4]H)(<[6]CH_3)(<:[7]CH_2CH_3)}};
\node[below] at (0,-1.6) {\small Énantiomère 1};
\node at (5.5,0) {\chemfig{C(-[2]OH)(-[4]H)(<:[6]CH_3)(<[7]CH_2CH_3)}};
\node[below] at (5.5,-1.6) {\small Énantiomère 2 (image miroir)};
\draw[dashed, popmuted] (2.75,-1.8) -- (2.75,1.2);
\end{tikzpicture}
\end{center}
Les deux formes sont images l'une de l'autre dans un miroir (les positions des groupes \ce{CH3} et \ce{C2H5} devant/derrière le plan sont interverties) et non superposables.
\end{corrige}

\begin{corrige}[Exercice 6 — Isomérie $Z/E$ : le bon candidat]
\textbf{1. Formules semi-développées :}
\begin{itemize}
\item pent-2-ène : \ce{CH3-CH=CH-CH2-CH3} ;
\item 2-méthylbut-2-ène : \ce{CH3-C(CH3)=CH-CH3}.
\end{itemize}

\textbf{2. Condition d'existence de l'isomérie $Z/E$ :} chacun des deux atomes de carbone de la double liaison \ce{C=C} doit porter \textbf{deux atomes ou groupes d'atomes différents}.

\textbf{3. Application :}
\begin{itemize}
\item Pent-2-ène : le carbone 2 porte \ce{H} et \ce{CH3} (différents) ; le carbone 3 porte \ce{H} et \ce{C2H5} (différents). La condition est remplie : \textbf{isomérie $Z/E$ possible}.
\item 2-méthylbut-2-ène : le carbone 2 porte \textbf{deux groupes \ce{CH3} identiques} : la condition n'est pas remplie, \textbf{pas d'isomérie $Z/E$}.
\end{itemize}
Seul le \textbf{pent-2-ène} présente l'isomérie $Z/E$.

\textbf{4. Représentation des deux stéréo-isomères :}
\begin{center}
\begin{tikzpicture}
\node at (0,0) {\chemfig{CH_3-[:0]C(-[2]H)=[0]C(-[2]H)-[0]CH_2CH_3}};
\node[below] at (0,-1.5) {\small $(Z)$-pent-2-ène : \ce{CH3} et \ce{C2H5} du même côté};
\node at (7.5,0) {\chemfig{CH_3-[:0]C(-[2]H)=[0]C(-[-2]H)-[0]CH_2CH_3}};
\node[below] at (7.5,-1.5) {\small $(E)$-pent-2-ène : \ce{CH3} et \ce{C2H5} de part et d'autre};
\end{tikzpicture}
\end{center}
Dans l'isomère $Z$ (\emph{zusammen}, « ensemble »), les deux groupes prioritaires (ici \ce{CH3} et \ce{C2H5}) sont du \textbf{même côté} de la double liaison ; dans l'isomère $E$ (\emph{entgegen}, « opposé »), ils sont de part et d'autre.

\textbf{5. Nature de la relation :} ces deux stéréo-isomères ne sont \textbf{pas} images l'un de l'autre dans un miroir (aucune rotation ne permet de passer de l'un à l'autre sans rompre la double liaison, et leur image miroir est superposable à chacun d'eux). Ce sont donc des \textbf{diastéréoisomères}.
\end{corrige}

\begin{corrige}[Exercice 7 — Polarimétrie sur le saccharose]
\textbf{1.} L'angle mesuré est \textbf{positif} : $\alpha = +13,3\degres$. Le plan de polarisation tourne donc vers la droite : le saccharose est \textbf{dextrogyre}.

\textbf{2.} D'après la loi de Biot : $\alpha = [\alpha] \times \ell \times C$, d'où :
\[ [\alpha] = \frac{\alpha}{\ell \times C} = \frac{13,3}{1,0 \times 0,20} = \SI{66,5}{\degree.mL.g^{-1}.dm^{-1}} \]
Le pouvoir rotatoire spécifique du saccharose à \SI{20}{\celsius} vaut $[\alpha] = +66,5\degres\cdot\text{mL}\cdot\text{g}^{-1}\cdot\text{dm}^{-1}$ (valeur proche de la valeur tabulée $+66,5$, excellente cohérence).

\textbf{3.} La même cuve et la même substance donnent le même $[\alpha]$ :
\[ C' = \frac{\alpha'}{[\alpha] \times \ell} = \frac{8,0}{66,5 \times 1,0} = \SI{0,12}{g.mL^{-1}} \]
Vérification : $66,5 \times 1,0 \times 0,12 = 7,98 \approx 8,0\degres$. Correct.

\textbf{4.} Avec un mélange racémique, les deux énantiomères sont présents en quantités égales ; leurs pouvoirs rotatoires, opposés, se compensent exactement :
\[ \alpha = [\alpha]\,\ell\,C_{\text{total}} \times \left(\tfrac{1}{2} - \tfrac{1}{2}\right) = 0\degres \]
On mesurerait un \textbf{angle nul} : le mélange est optiquement inactif.
\end{corrige}

\begin{corrige}[Exercice 8 — Newman : le 1,2-dichloroéthane]
\textbf{1. Projections de Newman} (l'observateur regarde selon l'axe \ce{C-C} ; le point central représente le carbone avant, le cercle le carbone arrière) :
\begin{center}
\begin{tikzpicture}[scale=0.9]
% Éclipsée
\draw[popline] (0,0) circle (1);
\foreach \a in {90,210,330}{
  \draw[popdark, thick] (0,0) -- (\a:1);
}
\foreach \a in {90,210,330}{
  \draw[popdark, thick] (\a:1) -- (\a:1.45);
}
\node at (90:1.7) {\ce{Cl}};
\node at (210:1.7) {\ce{H}};
\node at (330:1.7) {\ce{H}};
\node at (0,-2.1) {\small Conformation éclipsée};
% Décalée anti
\begin{scope}[xshift=5.5cm]
\draw[popline] (0,0) circle (1);
\foreach \a in {90,210,330}{
  \draw[popdark, thick] (0,0) -- (\a:1);
}
\foreach \a in {30,150,270}{
  \draw[popdark, thick] (\a:1) -- (\a:1.45);
}
\node at (90:1.25) {\ce{Cl}};
\node at (210:1.25) {\ce{H}};
\node at (330:1.25) {\ce{H}};
\node at (30:1.7) {\ce{H}};
\node at (150:1.7) {\ce{H}};
\node at (270:1.7) {\ce{Cl}};
\node at (0,-2.1) {\small Conformation décalée (anti)};
\end{scope}
\end{tikzpicture}
\end{center}
Dans la conformation éclipsée, les liaisons du carbone arrière sont cachées derrière celles du carbone avant (on les dessine légèrement décalées pour la lisibilité). Dans la conformation décalée \emph{anti}, les deux atomes de chlore sont diamétralement opposés (angle de \ang{180}).

\textbf{2.} La conformation \textbf{décalée (anti)} est la plus stable : les groupes volumineux \ce{Cl} sont le plus éloignés possible, ce qui minimise les répulsions électroniques et stériques. Dans la conformation éclipsée, les doublets des liaisons se font face, ce qui crée une tension d'éclipsage.

\textbf{3.} \textbf{Non}, on ne peut pas les séparer à température ordinaire : la rotation autour de la liaison simple \ce{C-C} est quasi libre (barrière énergétique très faible, de l'ordre de quelques \si{kJ.mol^{-1}}). Les conformations s'interconvertissent en permanence : ce ne sont pas des molécules différentes mais des \textbf{conformères} de la même molécule.

\textbf{4.} La conformation la plus stable du cyclohexane est la conformation \textbf{chaise}, dans laquelle toutes les liaisons sont décalées et les angles proches de \ang{109} :
\begin{center}
\begin{tikzpicture}[scale=0.9, popdark, thick]
\coordinate (C1) at (0,0.7);
\coordinate (C2) at (0.9,0);
\coordinate (C3) at (1.8,0.7);
\coordinate (C4) at (2.7,0);
\coordinate (C5) at (1.8,-0.7);
\coordinate (C6) at (0.9,-1.4);
\draw (C1)--(C2)--(C3)--(C4)--(C5)--(C6)--cycle;
\node[below, popmuted] at (1.35,-1.8) {\small Conformation chaise du cyclohexane};
\end{tikzpicture}
\end{center}
\end{corrige}

\begin{corrige}[Exercice 9 — L'alanine, une molécule chirale de la vie]
\textbf{1. Masse molaire.} Formule brute de l'alanine \ce{CH3-CH(NH2)-COOH} : \ce{C3H7NO2}.
\[ M = 3\,M(\ce{C}) + 7\,M(\ce{H}) + M(\ce{N}) + 2\,M(\ce{O}) = 3\times 12 + 7\times 1 + 14 + 2\times 16 = \SI{89}{g.mol^{-1}} \]

\textbf{2. Carbone asymétrique.} Le carbone 2 est lié à quatre groupes tous différents : \ce{H}, \ce{NH2}, \ce{CH3} et \ce{COOH}. C'est un carbone asymétrique :
\[ \chemfig{CH_3-\chembelow{C}{*}(-[2]NH_2)-COOH} \]

\textbf{3. Chiralité.} La molécule possède un seul carbone asymétrique et aucun plan de symétrie : elle n'est pas superposable à son image dans un miroir. Elle est donc \textbf{chirale} et existe sous deux énantiomères.

\textbf{4. Représentation de Cram.} Conventions : trait plein fin = liaison dans le plan ; trait plein épais = liaison vers l'avant ; trait hachuré = liaison vers l'arrière.
\begin{center}
\begin{tikzpicture}
\node at (0,0) {\chemfig{C(-[2]COOH)(-[6]H)(<[1]NH_2)(<:[7]CH_3)}};
\node[below] at (0,-1.6) {\small Énantiomère 1};
\node at (6,0) {\chemfig{C(-[2]COOH)(-[6]H)(<:[1]NH_2)(<[7]CH_3)}};
\node[below] at (6,-1.6) {\small Énantiomère 2};
\draw[dashed, popmuted] (3,-1.8) -- (3,1.2);
\end{tikzpicture}
\end{center}

\textbf{5. Projection de Fischer.} \ce{COOH} en haut, \ce{CH3} en bas (chaîne verticale vers l'arrière), \ce{NH2} et \ce{H} sur l'horizontale (vers l'avant) :
\begin{center}
\begin{tikzpicture}[scale=0.9, popdark]
% Fischer 1
\draw[thick] (0,-0.8)--(0,0.8); \draw[thick] (-0.8,0)--(0.8,0);
\node[above] at (0,0.8) {\ce{COOH}}; \node[below] at (0,-0.8) {\ce{CH3}};
\node[left] at (-0.8,0) {\ce{NH2}}; \node[right] at (0.8,0) {\ce{H}};
\node[below] at (0,-1.5) {\small Énantiomère 1};
% Fischer 2
\begin{scope}[xshift=5cm]
\draw[thick] (0,-0.8)--(0,0.8); \draw[thick] (-0.8,0)--(0.8,0);
\node[above] at (0,0.8) {\ce{COOH}}; \node[below] at (0,-0.8) {\ce{CH3}};
\node[left] at (-0.8,0) {\ce{H}}; \node[right] at (0.8,0) {\ce{NH2}};
\node[below] at (0,-1.5) {\small Énantiomère 2};
\end{scope}
\draw[dashed, popmuted] (2.5,-1.7) -- (2.5,1.5);
\end{tikzpicture}
\end{center}

\textbf{6. Concentration molaire.}
\[ n = \frac{m}{M} = \frac{4,45}{89} = \SI{0,050}{mol} \qquad ; \qquad c = \frac{n}{V} = \frac{0,050}{0,500} = \SI{0,10}{mol.L^{-1}} \]

\textbf{7.} L'élève a \textbf{tort} : la solution contient les deux énantiomères en quantités égales, c'est un \textbf{mélange racémique}. Les pouvoirs rotatoires des deux énantiomères, égaux en valeur absolue et de signes opposés, se compensent : l'angle de rotation mesuré est \textbf{nul}. La solution est optiquement inactive.

\textbf{8.} Les récepteurs biologiques (enzymes, sites actifs des protéines) sont eux-mêmes \textbf{chiraux} : ils ne reconnaissent qu'une seule configuration spatiale, comme une main droite n'entre que dans un gant droit. L'utilisation d'un seul énantiomère (l'alanine L) garantit que toutes les protéines ont une architecture spatiale cohérente et fonctionnelle ; l'incorporation de l'autre énantiomère fausserait le repliement des protéines.

\medskip
\textbf{Barème indicatif (sur 10) :} Q1 : 1 pt ; Q2 : 1,5 pt ; Q3 : 1 pt ; Q4 : 1,5 pt ; Q5 : 1,5 pt ; Q6 : 1,5 pt ; Q7 : 1 pt ; Q8 : 1 pt.

\textbf{Points de vigilance :} ne pas oublier l'unité de $M$ ; en Fischer, la verticale part \emph{vers l'arrière} et l'horizontale \emph{vers l'avant} ; une permutation de deux groupes donne l'énantiomère, deux permutations redonnent la même molécule ; un mélange racémique n'est \emph{pas} « sans carbone asymétrique », il est optiquement inactif par compensation.
\end{corrige}

\begin{corrige}[Exercice 10 — Une molécule à deux carbones asymétriques]
\textbf{1. Carbones asymétriques.} Dans \ce{CH3-CHOH-CHBr-CH3} :
\begin{itemize}
\item le carbone 2 est lié à \ce{H}, \ce{OH}, \ce{CH3} et \ce{CHBr-CH3} : quatre groupes différents $\Rightarrow$ C$^*$ ;
\item le carbone 3 est lié à \ce{H}, \ce{Br}, \ce{CH3} et \ce{CHOH-CH3} : quatre groupes différents $\Rightarrow$ C$^*$.
\end{itemize}
\[ \chemfig{CH_3-\chembelow{C}{*}(-[2]OH)-\chembelow{C}{*}(-[6]Br)-CH_3} \]
Ils sont \textbf{différents l'un de l'autre} car leurs quatre substituants ne sont pas les mêmes (\ce{OH} d'un côté, \ce{Br} de l'autre) : la molécule ne possède aucun plan de symétrie, il n'existe donc pas de forme \emph{méso}.

\textbf{2. Nombre de stéréo-isomères.} Avec $n = 2$ carbones asymétriques différents : $2^n = 2^2 = \textbf{4 stéréo-isomères}$ au maximum, tous effectivement distincts ici.

\textbf{3. Les quatre formes} (configuration de C2 notée en premier, celle de C3 en second ; $+$ et $-$ désignent les deux configurations opposées d'un même carbone) :
\begin{center}
\begin{tabularx}{\linewidth}{|l|X|X|}
\hline
\rowcolor{popblueL} \textbf{Forme} & \textbf{Configuration C2} & \textbf{Configuration C3} \\
\hline
A & $+$ & $+$ \\
\hline
B & $-$ & $-$ \\
\hline
C & $+$ & $-$ \\
\hline
D & $-$ & $+$ \\
\hline
\end{tabularx}
\end{center}
Représentation de Cram de la forme A (les autres s'obtiennent en inversant un ou deux centres) :
\begin{center}
\begin{tikzpicture}
\node at (0,0) {\chemfig{CH_3-[0]C(-[2]H)(<:[6]OH)-[0]C(-[2]H)(<:[6]Br)-[0]CH_3}};
\node[below] at (0,-1.6) {\small Forme A $(+,+)$};
\end{tikzpicture}
\end{center}

\textbf{4. Couples remarquables :}
\begin{itemize}
\item \textbf{A et B} : les deux centres sont inversés $(+,+)\leftrightarrow(-,-)$ ; A et B sont images dans un miroir et non superposables : ce sont des \textbf{énantiomères}.
\item \textbf{A et C} : seul le centre C3 est inversé ; A et C ne sont pas images miroir : ce sont des \textbf{diastéréoisomères}.
\end{itemize}

\textbf{5. Classement complet des relations deux à deux :}
\begin{center}
\begin{tabularx}{\linewidth}{|l|X|}
\hline
\rowcolor{popblueL} \textbf{Couple} & \textbf{Relation} \\
\hline
A -- B & énantiomères (deux centres inversés) \\
\hline
C -- D & énantiomères (deux centres inversés) \\
\hline
A -- C & diastéréoisomères (un seul centre inversé) \\
\hline
A -- D & diastéréoisomères (un seul centre inversé) \\
\hline
B -- C & diastéréoisomères (un seul centre inversé) \\
\hline
B -- D & diastéréoisomères (un seul centre inversé) \\
\hline
\end{tabularx}
\end{center}

\textbf{6. Activité optique des mélanges :}
\begin{itemize}
\item \textbf{A + B en quantités égales :} A et B sont des énantiomères ; le mélange est \textbf{racémique} : les pouvoirs rotatoires se compensent, la solution \textbf{ne dévie pas} le plan de la lumière polarisée ($\alpha = 0$).
\item \textbf{A + C en quantités égales :} A et C sont des \textbf{diastéréoisomères} ; leurs pouvoirs rotatoires n'ont aucune raison d'être opposés (ce ne sont pas des images miroir) : il n'y a pas compensation, la solution \textbf{dévie} le plan de la lumière polarisée.
\end{itemize}

\medskip
\textbf{Barème indicatif (sur 10) :} Q1 : 2 pts ; Q2 : 1 pt ; Q3 : 2 pts ; Q4 : 2 pts ; Q5 : 1,5 pt ; Q6 : 1,5 pt.

\textbf{Points de vigilance :} la règle des $2^n$ ne donne qu'un \emph{maximum} (une forme \emph{méso} le réduirait, mais elle n'existe pas ici car les deux C$^*$ sont différents) ; inverser \emph{un seul} centre donne un diastéréoisomère, inverser \emph{tous} les centres donne l'énantiomère ; un mélange équimolaire de diastéréoisomères n'est \emph{jamais} racémique.
\end{corrige}

\begin{corrige}[Exercice 11 — Le thalidomide : enjeu de santé publique]
\textbf{1.} Un mélange racémique est un mélange contenant les \textbf{deux énantiomères en quantités égales} ; il est optiquement inactif. Vendre le thalidomide sous cette forme était moins coûteux car la synthèse chimique classique produit naturellement les deux énantiomères en proportions égales : séparer l'énantiomère actif (dédoublement) ou mettre au point une synthèse asymétrique demande des étapes supplémentaires longues et chères.

\textbf{2.} Un \textbf{carbone asymétrique} est un atome de carbone tétragonal lié à quatre atomes ou groupes d'atomes tous différents. Un \textbf{énantiomère} est l'un des deux stéréo-isomères d'un couple de molécules images l'une de l'autre dans un miroir et non superposables.

\textbf{3.} Les énantiomères ont les mêmes propriétés physiques et chimiques \emph{dans un environnement achiral}. Mais l'organisme vivant est un environnement \textbf{chiral} : enzymes, récepteurs et protéines sont eux-mêmes constitués de molécules chirales (acides aminés L, sucres D). Chaque énantiomère s'ajuste différemment à ces récepteurs, comme une main droite et une main gauche dans un gant droit : l'un est reconnu et produit l'effet voulu, l'autre est inactif ou provoque un effet toxique (tératogène pour le thalidomide).

\textbf{4.} D'après le document :
\begin{itemize}
\item \textbf{Ibuprofène :} seul l'énantiomère $(S)$ est actif contre la douleur ; l'autre est (quasi) inactif.
\item \textbf{Carvone :} un énantiomère sent la \textbf{menthe}, l'autre le \textbf{carvi} : même formule, odeurs différentes.
\end{itemize}

\textbf{5.} La distillation sépare des espèces de \textbf{températures d'ébullition différentes}. Or deux énantiomères ont exactement la \textbf{même température d'ébullition} (mêmes propriétés physiques) : ils s'évaporent ensemble et la distillation ne peut pas les séparer. Il faut des méthodes chirales (chromatographie chirale, formation de sels de diastéréoisomères, enzymes).

\textbf{6. Paragraphe de synthèse (exemple de rédaction) :} La tragédie du thalidomide a montré que deux énantiomères, pourtant de même formule brute, peuvent avoir des effets biologiques radicalement opposés : l'un soigne, l'autre malforme. Maîtriser la stéréochimie, c'est donc maîtriser l'efficacité et l'innocuité des médicaments. Aujourd'hui, les industriels développent des synthèses asymétriques et des contrôles qualité polarimétriques pour garantir que le principe actif délivré est le bon stéréo-isomère. Cet enjeu concerne aussi le Cameroun : les pharmacies et centres de santé de Yaoundé, Douala ou des zones rurales doivent pouvoir compter sur des médicaments contrôlés, correctement étiquetés et exempts de l'énantiomère nocif. Former des techniciens et pharmaciens sensibilisés à la stéréochimie, c'est protéger la santé de toute la population.

\medskip
\textbf{Barème indicatif (sur 10) :} Q1 : 1,5 pt ; Q2 : 1,5 pt ; Q3 : 2 pts ; Q4 : 1,5 pt ; Q5 : 1,5 pt ; Q6 : 2 pts.

\textbf{Points de vigilance :} bien distinguer « mêmes propriétés physiques » (vrai dans un environnement achiral) et « mêmes effets biologiques » (faux, car le vivant est chiral) ; le mélange racémique n'est pas dangereux \emph{parce qu'il est inactif optiquement}, mais parce qu'il contient l'énantiomère toxique ; dans la synthèse, mobiliser le vocabulaire exigé : énantiomère, chiralité, mélange racémique, dédoublement.
\end{corrige}

\newpage

\coursec{Fiche bilan}

\begin{popfigure}
\begin{tikzpicture}[scale=0.95, transform shape,
  central/.style={ellipse, fill=popblue, text=white, font=\bfseries\small, align=center, minimum width=3.2cm, minimum height=1.4cm},
  branche/.style={rounded corners, fill=#1, text=popdark, font=\footnotesize, align=center, minimum width=3.4cm, minimum height=1.1cm, inner sep=3pt},
  lien/.style={thick, popmuted}]
\node[central] (C) at (0,0) {Stéréochimie\\isomérie};
\node[branche=popblueL] (A) at (-5.6,2.6) {Isomérie de constitution\\chaîne -- position -- fonction};
\node[branche=poporangeL] (B) at (0,3.4) {Chiralité\\carbone asymétrique \ce{C*}\\énantiomères (Cram)};
\node[branche=popgreenL] (C1) at (5.6,2.6) {Activité optique\\lumière polarisée\\loi de Biot $\alpha=[\alpha]\,\ell\,C$};
\node[branche=poppinkL] (D) at (-5.6,-2.6) {Diastéréoisomères\\Z/E -- deux \ce{C*}\\projection de Fischer};
\node[branche=popgoldL] (E) at (0,-3.4) {Conformations\\Newman : éclipsée / décalée\\cyclohexane : chaise / bateau};
\node[branche=poppurpleL] (F) at (5.6,-2.6) {Enjeux\\thalidomide, ibuprofène\\carvone (arômes)};
\draw[lien] (C) -- (A); \draw[lien] (C) -- (B); \draw[lien] (C) -- (C1);
\draw[lien] (C) -- (D); \draw[lien] (C) -- (E); \draw[lien] (C) -- (F);
\end{tikzpicture}
\legende{Carte mentale de la leçon : les six idées forces de la stéréochimie.}
\end{popfigure}

\begin{bilan}[L'essentiel de la leçon]
\textbf{Définitions clés}
\begin{itemize}
  \item \cle{Isomères} : composés de \textbf{même formule brute} mais de propriétés différentes.
  \item \cle{Isomérie de constitution} : \textbf{chaîne} (squelette différent), \textbf{position} (place du groupe caractéristique ou de la liaison multiple), \textbf{fonction} (nature du groupe caractéristique).
  \item \cle{Carbone asymétrique} \ce{C*} : carbone tétraédrique portant \textbf{quatre groupes différents}.
  \item \cle{Molécule chirale} : non superposable à son image dans un miroir. Deux images non superposables = \cle{énantiomères}.
  \item \cle{Mélange racémique} : mélange équimolaire des deux énantiomères ; optiquement \textbf{inactif}.
  \item \cle{Diastéréoisomères} : stéréo-isomères qui ne sont \textbf{pas} images l'un de l'autre (Z/E ; molécules à plusieurs \ce{C*}).
  \item \cle{Conformations} : stéréo-isomères interconvertibles par simple rotation autour d'une liaison $\sigma$.
\end{itemize}

\textbf{Équations et relations à connaître par cœur}
\begin{itemize}
  \item Loi de Biot : $\alpha = [\alpha]\times \ell \times C$ \quad ($\alpha$ en degrés, $\ell$ en dm, $C$ en \si{g.mL^{-1}} ou \si{kg.L^{-1}} selon l'unité de $[\alpha]$).
  \item Nombre maximal de stéréo-isomères d'une molécule à $n$ carbones asymétriques : $N = 2^{n}$.
  \item Dextrogyre $(+)$ : dévie la lumière polarisée vers la droite ; lévogyre $(-)$ : vers la gauche.
\end{itemize}

\textbf{Représentations et cas types}
\begin{center}
\begin{tabularx}{\linewidth}{|l|X|X|}
\hline
\rowcolor{popblueL} \textbf{Notion} & \textbf{Représentation} & \textbf{Point de vigilance} \\
\hline
Cram & trait plein (plan), coin plein (avant), hachures (arrière) & un seul coin et une seule hachure par \ce{C*} \\
\hline
Fischer & croix : verticale vers l'arrière, horizontale vers l'avant & échanger 2 groupes $\Rightarrow$ énantiomère \\
\hline
Z / E & \ce{C=C} : Z = prioritaires du même côté ; E = opposés & priorité : atome de $Z$ le plus grand \\
\hline
Newman & cercle (C arrière) + point (C avant) & décalée plus stable qu'éclipsée \\
\hline
Cyclohexane & chaise (stable) et bateau (instable) & positions axiales et équatoriales \\
\hline
\end{tabularx}
\end{center}

\textbf{Méthodes express}
\begin{enumerate}
  \item Repérer un \ce{C*} : chercher les carbones $sp^3$ portant 4 substituants différents (jamais \ce{CH2}, \ce{CH3}, ni \ce{C=C}).
  \item Deux stéréo-isomères ? Images en miroir non superposables $\Rightarrow$ énantiomères ; sinon $\Rightarrow$ diastéréoisomères.
  \item Z ou E ? Classer les priorités sur chaque carbone de la double liaison, puis comparer les positions.
  \item Activité optique : un seul \ce{C*} (ou molécule chirale) $\Rightarrow$ active ; racémique ou mésomérique $\Rightarrow$ inactive.
\end{enumerate}
\end{bilan}

\begin{aretenir}[Checklist avant le Bac]
\begin{itemize}
  \item[$\square$] Je sais distinguer isomérie de chaîne, de position et de fonction sur des exemples simples.
  \item[$\square$] Je sais repérer tous les carbones asymétriques d'une molécule et les marquer d'une étoile.
  \item[$\square$] Je sais dessiner une molécule chirale en représentation de Cram (coin plein, hachures).
  \item[$\square$] Je sais dessiner les deux énantiomères d'une molécule en Cram et en projection de Fischer.
  \item[$\square$] Je sais qu'un mélange racémique (50/50) est optiquement inactif.
  \item[$\square$] Je connais la loi de Biot $\alpha=[\alpha]\,\ell\,C$ et le rôle du polarimètre.
  \item[$\square$] Je sais attribuer les descripteurs Z et E à un alcène.
  \item[$\square$] Je sais calculer le nombre de stéréo-isomères $2^{n}$ et citer le cas de l'acide tartrique (forme méso).
  \item[$\square$] Je sais distinguer énantiomères et diastéréoisomères sans hésiter.
  \item[$\square$] Je sais dessiner les projections de Newman éclipsée et décalée de l'éthane.
  \item[$\square$] Je sais esquisser les conformations chaise et bateau du cyclohexane.
  \item[$\square$] Je sais citer un exemple d'enjeu : thalidomide, ibuprofène ou carvone.
\end{itemize}
\end{aretenir}

\findecours

\end{document}
